\begin{filecontents*}[overwrite]{note_refs.bib}
@inproceedings{lambert2022variational,
  author    = {Lambert, Marc and Chewi, Sinho and Bach, Francis and Bonnabel, Silv\`ere and Rigollet, Philippe},
  title     = {Variational inference via {W}asserstein gradient flows},
  booktitle = {Advances in Neural Information Processing Systems},
  volume    = {35},
  year      = {2022}
}
@inproceedings{diao2023forward,
  author    = {Diao, Michael Ziyang and Balasubramanian, Krishna and Chewi, Sinho and Salim, Adil},
  title     = {Forward-backward {G}aussian variational inference via {JKO} in the {B}ures--{W}asserstein space},
  booktitle = {Proceedings of the 40th International Conference on Machine Learning},
  year      = {2023}
}
@article{roeder2017sticking,
  author  = {Roeder, Geoffrey and Wu, Yuhuai and Duvenaud, David K.},
  title   = {Sticking the landing: Simple, lower-variance gradient estimators for variational inference},
  journal = {Advances in Neural Information Processing Systems},
  volume  = {30},
  year    = {2017}
}
@article{chen2023sampling,
  author  = {Chen, Yifan and Huang, Daniel Zhengyu and Huang, Jiaoyang and Reich, Sebastian and Stuart, Andrew M.},
  title   = {Sampling via gradient flows in the space of probability measures},
  journal = {arXiv preprint arXiv:2310.03597},
  year    = {2023}
}
@article{amari1998natural,
  author  = {Amari, Shun-Ichi},
  title   = {Natural gradient works efficiently in learning},
  journal = {Neural Computation},
  volume  = {10},
  number  = {2},
  pages   = {251--276},
  year    = {1998}
}
@article{opper2009variational,
  author  = {Opper, Manfred and Archambeau, C\'edric},
  title   = {The variational {G}aussian approximation revisited},
  journal = {Neural Computation},
  volume  = {21},
  number  = {3},
  pages   = {786--792},
  year    = {2009}
}
@article{kim2024stl,
  author  = {Kim, Kyurae and Ma, Yian and Gardner, Jacob R.},
  title   = {On the convergence of black-box variational inference and the sticking-the-landing estimator},
  journal = {arXiv preprint arXiv:2307.14642},
  year    = {2024}
}
@article{laurent2000adaptive,
  author  = {Laurent, B{\'e}atrice and Massart, Pascal},
  title   = {Adaptive estimation of a quadratic functional by model selection},
  journal = {Annals of Statistics},
  volume  = {28},
  number  = {5},
  pages   = {1302--1338},
  year    = {2000}
}
@article{domke2023provable,
  author  = {Domke, Justin and Garrigos, Guillaume and Gower, Robert},
  title   = {Provable convergence guarantees for black-box variational inference},
  journal = {Advances in Neural Information Processing Systems},
  volume  = {36},
  year    = {2023}
}
\end{filecontents*}

\documentclass[11pt]{article}

% ============================================================
% Packages
% ============================================================
\usepackage[margin=1in]{geometry}
\usepackage{amsmath,amssymb,amsthm,mathtools}
\usepackage{bm}
\usepackage{xcolor}
\usepackage{algorithm}
\usepackage{algpseudocode}
\usepackage[numbers,sort&compress]{natbib}
\usepackage{hyperref}
\usepackage{aliascnt}
\usepackage[nameinlink,capitalize]{cleveref}

\hypersetup{
    colorlinks=true,
    linkcolor=blue!60!black,
    citecolor=blue!60!black,
    urlcolor=blue!60!black,
    pdftitle={Covariance Bootstrap, Quadratic Rescue, and Gaussian-Core Sharp Regions for Gaussian Natural Gradient Variational Inference},
    pdfauthor={}
}

% ============================================================
% Theorem environments
% ============================================================
\theoremstyle{plain}
\newtheorem{theorem}{Theorem}[section]

\newaliascnt{proposition}{theorem}
\newtheorem{proposition}[proposition]{Proposition}
\aliascntresetthe{proposition}

\newaliascnt{lemma}{theorem}
\newtheorem{lemma}[lemma]{Lemma}
\aliascntresetthe{lemma}

\newaliascnt{corollary}{theorem}
\newtheorem{corollary}[corollary]{Corollary}
\aliascntresetthe{corollary}

\theoremstyle{definition}
\newaliascnt{assumption}{theorem}
\newtheorem{assumption}[assumption]{Assumption}
\aliascntresetthe{assumption}

\newaliascnt{definition}{theorem}
\newtheorem{definition}[definition]{Definition}
\aliascntresetthe{definition}

\theoremstyle{remark}
\newaliascnt{remark}{theorem}
\newtheorem{remark}[remark]{Remark}
\aliascntresetthe{remark}

\crefname{theorem}{Theorem}{Theorems}
\Crefname{theorem}{Theorem}{Theorems}
\crefname{proposition}{Proposition}{Propositions}
\Crefname{proposition}{Proposition}{Propositions}
\crefname{lemma}{Lemma}{Lemmas}
\Crefname{lemma}{Lemma}{Lemmas}
\crefname{corollary}{Corollary}{Corollaries}
\Crefname{corollary}{Corollary}{Corollaries}
\crefname{assumption}{Assumption}{Assumptions}
\Crefname{assumption}{Assumption}{Assumptions}
\crefname{definition}{Definition}{Definitions}
\Crefname{definition}{Definition}{Definitions}
\crefname{remark}{Remark}{Remarks}
\Crefname{remark}{Remark}{Remarks}
\crefname{algorithm}{Algorithm}{Algorithms}
\Crefname{algorithm}{Algorithm}{Algorithms}
\crefname{section}{Section}{Sections}
\Crefname{section}{Section}{Sections}
\crefname{appendix}{Appendix}{Appendices}
\Crefname{appendix}{Appendix}{Appendices}

% ============================================================
% Macros
% ============================================================
\newcommand{\R}{\mathbb{R}}
\newcommand{\N}{\mathcal{N}}
\newcommand{\E}{\mathbb{E}}
\newcommand{\Prob}{\mathbb{P}}
\newcommand{\Var}{\operatorname{Var}}
\newcommand{\KL}{\mathrm{KL}}
\newcommand{\Tr}{\operatorname{Tr}}
\newcommand{\SPD}{\mathbb{S}_{++}}
\newcommand{\Sym}{\mathbb{S}}
\newcommand{\Id}{I}
\newcommand{\dd}{\mathrm{d}}
\newcommand{\grad}{\operatorname{grad}}
\newcommand{\op}{\mathrm{op}}
\newcommand{\Frob}{\mathrm{F}}
\newcommand{\post}{\mathrm{post}}
\newcommand{\FR}{\mathrm{FR}}
\newcommand{\BW}{\mathrm{BW}}
\newcommand{\Riem}{\mathrm{R}}
\newcommand{\calE}{\mathcal{E}}
\newcommand{\PG}{\mathcal{P}_{\mathrm G}}
\newcommand{\Aspace}{\mathcal{A}}
\newcommand{\Fil}{\mathcal{F}}
\newcommand{\STL}{\mathrm{STL}}
\newcommand{\norm}[1]{\left\lVert #1\right\rVert}
\newcommand{\inner}[2]{\left\langle #1,#2\right\rangle}
\newcommand{\logp}{\log_+}
\newcommand{\glob}{\mathrm{glob}}
\newcommand{\cov}{\mathrm{cov}}
\newcommand{\loc}{\mathrm{loc}}
\newcommand{\Gradnorm}{\mathsf{G}}
\newcommand{\LH}{L_H}
\newcommand{\barLH}{\overline{L}_H}
\newcommand{\LB}{L_\Lambda}
\newcommand{\Retr}{\operatorname{Retr}}
\newcommand{\Kresc}{\mathsf{K}}
\newcommand{\Fresc}{\mathsf{F}}
\newcommand{\diag}{\operatorname{diag}}
\newcommand{\Ureg}{\mathcal{U}}
\newcommand{\Hreg}{\mathcal{H}}
\newcommand{\FGauss}{F_{\mathrm G}}
\newcommand{\gstar}{\gamma_\star}
\newcommand{\gund}{\underline{\gamma}}
\newcommand{\Lstar}{L_\star}
\newcommand{\sharpthr}{\Delta^{\sharp}}
\newcommand{\Kng}{K_{\mathrm{ng}}}
\newcommand{\Kdisc}{K^{\mathrm{disc}}}
\newcommand{\LHs}{L_{H,\star}}

% ============================================================
% Title
% ============================================================
\title{Covariance Bootstrap, Quadratic Rescue, and Gaussian-Core Sharp
Regions for Gaussian Natural Gradient Variational Inference:\\
Deterministic and Price/Hessian Stochastic Theory}
\author{}
\date{}

\begin{document}
\maketitle

\begin{abstract}
We analyze Gaussian variational inference along the Fisher--Rao natural gradient flow and its two
discretizations---a Riemannian matrix-exponential retraction scheme and a KL/Bregman resolvent
scheme---throughout in the intrinsic Fisher--Rao metric on Gaussian space, which carries the factor
$\tfrac12$ on the covariance block. Deterministically, a \emph{covariance bootstrap} replaces the
frozen covariance floor of the standard global proofs by an explicit time-varying lower envelope,
improving the initialization dependence from $1/(\beta\lambda_{0,\min})$ to
$\logp(1/(\beta\lambda_{0,\min}))$---order-optimal already for one-dimensional Gaussian targets---and
organizing the complexity into covariance burn-in, global localization, and local convergence; the
burn-in is removable by a short Bures--Wasserstein transport warm-start or a single quadratic-rescue
Hessian query. The local stage rests on a \emph{Gaussian-core decomposition} $F=\FGauss+H$ that
dissipates the exact Fisher--Rao field of the Gaussian target and Taylor-bounds only the non-Gaussian
remainder $H$, yielding certified energy-sublevel regions that are sharp for Gaussian targets in every
dimension and degrade continuously through a non-Gaussian modulus vanishing at Gaussianity.
Stochastically, under a Price/Hessian oracle returning $(\nabla V(X),\nabla^2V(X))$, the quadratic
rescue initializes $C_0$ on $[\beta^{-1}\Id,\alpha^{-1}\Id]$, and the sticking-the-landing iterations
keep the covariance pathwise on $[(2\beta)^{-1}\Id,\alpha^{-1}\Id]$ for the Riemannian scheme and on
$[\beta^{-1}\Id,\alpha^{-1}\Id]$ for the KL scheme, with no stopping time, clipping, or containment
hypothesis. At one sample per iteration, the constant-stepsize global bounds contract geometrically
toward explicit residual levels on iteration scales $\widetilde O(\kappa^3)$ under strong
log-concavity alone and $\widetilde O(\kappa^2)$ under an additional Hessian-Lipschitz assumption;
their intrinsic additive terms vanish for Gaussian targets, and a common decreasing-stepsize schedule
gives floor-free $O(1/N)$ convergence. Inside the certified region, a
stochastic Gaussian-core local theory gives a mean-square contraction on the non-exit event with a
target-adaptive additive floor proportional to the squared whitened Hessian-Lipschitz modulus
$\LHs^2$, paired with a finite-horizon exit-probability bound and a global-to-local entry corollary.
The local spectral scale obeys $\Gamma=O(1+\sqrt{N_\theta}\log\kappa)$, so at the certified stepsize
with a sufficiently large batch the local iteration complexity is
$O\bigl((1+\sqrt{N_\theta}\log\kappa)^2\log\tfrac1\epsilon\bigr)$, becoming $\kappa$-independent,
$O(\log\tfrac1\epsilon)$, in the small-$\Gamma$ (in particular near-Gaussian) regime.
\end{abstract}

\tableofcontents

\section{Introduction}
\label{sec:intro}

Let $\rho_{\post}(\theta)\propto\exp(-V(\theta))$ be a target on $\R^{N_\theta}$ and let
\[
    \PG:=\bigl\{\rho_a=\N(m,C):a=(m,C)\in\Aspace:=\R^{N_\theta}\times\SPD(N_\theta)\bigr\}
\]
be the Gaussian variational family. Gaussian variational inference minimizes
\[
    \calE(a):=\KL(\rho_a\,\|\,\rho_{\post}),
    \qquad
    a_\star=(m_\star,C_\star)\in\operatorname*{argmin}_{a\in\Aspace}\calE(a),
\]
and the natural gradient flow of $\calE$ on $\PG$ is the Fisher--Rao gradient flow. This note
develops the deterministic and stochastic theory of that flow and its two discretizations---the
Riemannian retraction scheme and the KL/Bregman scheme---and contributes the following.

\paragraph{A covariance bootstrap for the deterministic theory.}
Global proofs for the Fisher--Rao flow customarily freeze the lower covariance bound at
$\lambda_{\min}=\min\{\lambda_{0,\min},1/\beta\}$ throughout the global phase. Tracking the growth
of the covariance through an explicit time-varying covariance lower envelope---an exactly solvable
comparison envelope for the KL recursion and a conservative comparison envelope for the Riemannian
update---improves the
dependence on the initial covariance from $1/(\beta\lambda_{0,\min})$ to $\logp(1/(\beta\lambda_{0,\min}))$ and
exhibits the iteration complexity as
\[
    \boxed{\;
    \text{covariance burn-in}
    \;+\;
    \text{global localization}
    \;+\;
    \text{local convergence}.
    \;}
\]

\paragraph{Two deterministic burn-in removers.}
The covariance burn-in can be removed by either of two optional deterministic initializers
(\Cref{sec:transport-bootstrap}). The first is a short Bures--Wasserstein transport warm-start: in the
underdispersed regime, Fisher--Rao covariance growth is multiplicative while Bures--Wasserstein growth
is additive and scale-independent, so using transport only as a warm-start and then switching
permanently to Fisher--Rao removes the burn-in with an expectation-level oracle and no pointwise
Hessian. The second, under a pointwise Hessian oracle, is the \emph{quadratic rescue} of the
stochastic pipeline used here in its deterministic role: a single Hessian query lands the covariance
directly on the curvature band $[\tfrac1\beta,\tfrac1\alpha]$, so the burn-in count is zero from the
first step (\Cref{thm:disc-rescue-to-FR}), and is exact for Gaussian targets. The two initializers
trade oracle strength against landing scale (\Cref{rem:two-initializers}); neither affects the local
factor.

\paragraph{A Gaussian-sharp local region via the Gaussian core.}
The local convergence theory (\Cref{sec:local}) is built on the Gaussian-core decomposition
$F=\FGauss+H$ with $\FGauss(m,C)=(-Cm,\,C-C^2)$ the exact Fisher--Rao field of the standard
Gaussian target. The Gaussian core is dissipated \emph{exactly}, and only the non-Gaussian
remainder $H$ is Taylor-bounded through a second-order score calculus whose modulus $\Kng(\delta)$
vanishes for Gaussian targets. The consequence is a certified energy-sublevel region that equals
the exact supremal Gaussian threshold $\sharpthr_{\mathrm G}(\rho)$ in every dimension, rather than
merely respecting it, and degrades continuously for general strongly log-concave targets.

\paragraph{A quadratic-rescued Price/Hessian stochastic theory.}
The stochastic pipeline is
\[
    \boxed{\;
    \text{one Price/Hessian quadratic rescue}
    \;\longrightarrow\;
    \text{minibatch Price/Hessian STL Fisher--Rao iterations.}
    \;}
\]
The rescue is a single \emph{deterministic} Price/Hessian query at the initial mean: with
$G_{\mathrm{in}}=\nabla V(m_{\mathrm{in}})$ and $H_{\mathrm{in}}=\nabla^2V(m_{\mathrm{in}})$, it
sets $m_0=m_{\mathrm{in}}-H_{\mathrm{in}}^{-1}G_{\mathrm{in}}$ and $C_0=H_{\mathrm{in}}^{-1}$. By
operator antitonicity this initializes the covariance on the curvature band
$\tfrac1\beta\Id\preceq C_0\preceq\tfrac1\alpha\Id$ pathwise, and for a Gaussian target it recovers
the exact optimizer in one query (\Cref{thm:gauss-recovery}). The subsequent iterations use STL
for the mean/score and sampled Hessians for the covariance; because the sampled Hessian obeys the
almost-sure spectral bound $\alpha\Id\preceq\nabla^2V(X)\preceq\beta\Id$, the Riemannian band
$\tfrac1{2\beta}\Id\preceq C_n\preceq\tfrac1\alpha\Id$ is forward invariant pathwise when
$\Delta t\leq1/\kappa$, while the KL scheme preserves the sharper initialization band for every
$\Delta t>0$. The global theorems therefore hold in plain expectation with no stopping time, clipping,
or containment hypothesis.

\paragraph{Two stochastic variance regimes.}
We treat the intrinsic Hessian fluctuation $\Psi$ in two ways. Under
\Cref{assump:logconcave-smooth} alone, a self-bounding estimate
(\Cref{lem:Psi-self-bound}) yields a certified residual level $O(\kappa N_\theta)$ and a
single-sample iteration scale $\widetilde O(\kappa^3)$ for contraction to a prescribed accuracy above
that level. Under an additional Hessian-Lipschitz assumption
(\Cref{assump:global-Hess-Lip}), the fluctuation is controlled directly by the nonquadraticity of
the target (\Cref{lem:Psi-Hg}), giving a certified residual level
$O(\kappa N_\theta^2\barLH^2/B)$ and a single-sample iteration scale
$\widetilde O(\kappa^2)$ above that level; the improvement is exactly
$\widetilde O(\kappa^3)\to\widetilde O(\kappa^2)$ (\Cref{sec:regime-comparison}). The intrinsic additive
floors vanish for Gaussian targets, although the coarse self-bounding corollaries do not retain this
exact cancellation after substituting their universal bound on $\Psi$; a common decreasing-stepsize
schedule gives floor-free convergence for non-Gaussian targets (\Cref{thm:decreasing}).

\paragraph{Stochastic local convergence.}
The global stochastic theorems contract the expected gap at the global $\kappa$-dependent rate toward
their displayed residual levels. Once an iterate has
entered the certified local region, we prove a stochastic analogue of the deterministic local theory
(\Cref{sec:stoch-local}): a mean-square contraction on the non-exit event at the same order as the
deterministic local contraction (per-step factor $1-\mu_\bullet\Delta t$ with
$\mu_{\Riem}=1/(2(4+\Gamma))$ and $\mu_{\KL}=1/(2(4+2\Delta t+\Gamma))$; the deterministic squared-norm factor is $(1-\mu_\bullet\Delta t)^2$), with a
target-adaptive additive floor $O\!\bigl(\Delta t\,(4+\Gamma)\,N_\theta^2\LHs^2/B\bigr)$
proportional to the \emph{squared} whitened Hessian-Lipschitz modulus $\LHs^2$ ($\LHs\leq\barLH$), a finite-horizon
high-probability containment bound, and a Markov-based global-to-local entry corollary
(\Cref{cor:stoch-three-stage}). The bound is a killed-process quantity, not a stopped-process bound;
because the STL mean-score noise is not of bounded support (even for a Gaussian target, away from
$C=\Id$ it equals $(C^{-1/2}-C^{1/2})Z$), pathwise invariance of a bounded local sublevel cannot be
guaranteed at fixed stepsize (without projection, clipping, or a safeguard), so non-exit contraction
plus exit-probability control is the honest resolution rather than infinite-horizon containment. For Gaussian targets the \emph{additive} floor vanishes
(\Cref{cor:stoch-local-gauss}); the noise is pathwise zero (hence deterministic contraction, no
exit-probability control) whenever the covariance is initialized matched, $C_0=C_\star$, even if
$m_0\neq m_\star$---quadratic rescue is the special case $a_0=a_\star$---while generic $C_0\neq C_\star$
retains unbounded mean noise and is covered only by the non-exit and exit-probability bounds. The
local spectral scale satisfies $\Gamma=O(1+\sqrt{N_\theta}\log\kappa)$ (a matrix quadratic-form
estimate; \Cref{lem:linear-diagonal-bounds}), so under the universal local window and the large-batch
branch $B\gtrsim V_1$ (with $V_1=24+6N_\theta\LHs^2$) the iteration complexity is
$O((1+\sqrt{N_\theta}\log\kappa)^2\log\tfrac1\epsilon)$
(large-$\kappa$ shorthand $O(N_\theta\log^2\kappa\,\log\tfrac1\epsilon)$; for general $B$, the extra
factor $\max\{1,V_1/B\}$); under uniform local
conditioning ($\Gamma=O(1)$, $B\gtrsim V_1$) it becomes $O(\log\tfrac1\epsilon)$, independent of
$\kappa$. No intermediate $O(\log\kappa)$ theorem is claimed (\Cref{rem:local-condnum}). The
energy-descent stepsize bound $1/(\beta_\star\bar\Lambda_\delta)$ is \emph{not} needed for this local
phase---invariance is supplied by the exit time, not deterministic descent
(\Cref{lem:local-map-contraction})---which is what keeps the local rate free of a spurious
polynomial-$\kappa$ factor.

\paragraph{Scope.}
The stochastic results presuppose a Price/Hessian oracle; the gradient-only covariance estimator,
whose samplewise spectrum is unbounded, is outside the scope of the pathwise arguments used here
(\Cref{rem:price}). The explicit constants in the stochastic local theory ($V_0,V_1$, the aggregate
coefficients $64,68$, and the exponential-remainder constant $e/2$) are conservative and not
optimized; only their orders are used (\Cref{rem:stoch-local-scope}).

\section{Setup and notation}
\label{sec:setup}

We write the energy gap
\[
    \Delta(a):=\calE(a)-\calE(a_\star),
    \qquad
    \Delta_0:=\Delta(a_0),
    \qquad
    \Delta_n:=\Delta(a_n),
\]
and denote by $\Delta t>0$ the Fisher--Rao stepsize (written $h$ in one-step lemmas). Throughout,
$\Delta t$ carries the differential $t$, so there is no clash with the energy gap $\Delta(\cdot)$;
$N_\theta$ is the parameter dimension.

\begin{assumption}[Strong log-concavity and smoothness]
\label{assump:logconcave-smooth}
There exist $0<\alpha\leq\beta<\infty$ with
\[
    \alpha\Id\preceq\nabla^2V(\theta)\preceq\beta\Id,
    \qquad
    \forall\theta\in\R^{N_\theta}.
\]
We write $\kappa:=\beta/\alpha\geq1$.
\end{assumption}

\Cref{assump:logconcave-smooth} is the standing regularity assumption throughout. The local phase
of \Cref{sec:local} additionally uses whitened bounds $\alpha_\star,\beta_\star$ defined there, and
the sharpened stochastic theory of \Cref{sec:stoch-Hlip} additionally uses a Hessian-Lipschitz
assumption (\Cref{assump:global-Hess-Lip}).

For $a=(m,C)\in\Aspace$ we write
\begin{equation}
\label{eq:gAR}
    g(a):=-\E_{\rho_a}[\nabla V(X)],
    \qquad
    A(a):=\E_{\rho_a}[\nabla^2V(X)],
    \qquad
    R(a):=C^{-1}-A(a),
\end{equation}
and call $R(a)$ the \emph{precision residual}. \Cref{assump:logconcave-smooth} gives
$\alpha\Id\preceq A(a)\preceq\beta\Id$. At the minimizer,
\begin{equation}
\label{eq:first-order}
    g(a_\star)=0,
    \qquad
    A(a_\star)=C_\star^{-1},
    \qquad\text{i.e.}\qquad
    R(a_\star)=0 .
\end{equation}
We use the energy splitting $\calE=\calE_1+\calE_2$,
\begin{equation}
\label{eq:splitting}
    \calE_1(a)=\E_{\rho_a}[\log\rho_a]=-\tfrac12\log\det C+\text{const.},
    \qquad
    \calE_2(a)=-\E_{\rho_a}[\log\rho_{\post}]=\E_{\rho_a}[V]+\log Z .
\end{equation}

\subsection{Initial covariance and the band constant}

We assume $\lambda_{0,\min}\Id\preceq C_0\preceq\lambda_{0,\max}\Id$ and set
\begin{equation}
\label{eq:Lambda}
    \lambda_0:=\lambda_{0,\min},
    \qquad
    \Lambda:=\max\Bigl\{\lambda_{0,\max},\frac1\alpha\Bigr\},
    \qquad
    \LB:=\beta\Lambda,
    \qquad
    \logp x:=\max\{\log x,0\}.
\end{equation}
Since $\Lambda\geq1/\alpha$ we always have $\LB\geq\kappa\geq1$.

\subsection{The Gaussian Fisher--Rao metric}
\label{sec:metric}

\begin{definition}[Gaussian Fisher--Rao metric]
\label{def:FR-metric}
For $a=(m,C)\in\Aspace$ and tangent vectors
$(u,X),(v,Y)\in\R^{N_\theta}\times\Sym(N_\theta)$, the intrinsic (information-geometric)
Fisher--Rao metric on $\PG$ is
\begin{equation}
\label{eq:FR-metric}
    g_a\bigl((u,X),(v,Y)\bigr)
    :=
    u^\top C^{-1}v+\frac12\Tr\bigl(C^{-1}XC^{-1}Y\bigr),
\end{equation}
with induced tangent norm
\begin{equation}
\label{eq:tangent-norm}
    \norm{(u,X)}_{a,\FR}^2
    :=
    u^\top C^{-1}u+\frac12\norm{C^{-1/2}XC^{-1/2}}_{\Frob}^2 .
\end{equation}
We abbreviate $\inner{\cdot}{\cdot}_a:=g_a(\cdot,\cdot)$ and $\norm{\cdot}_a:=\norm{\cdot}_{a,\FR}$.
\end{definition}

The factor $\tfrac12$ on the covariance block is the one induced by the Fisher information of the
Gaussian family.

\begin{lemma}[Gaussian covariance derivative; Price's theorem, directional form]
\label{lem:price-derivative}
Let $C\in\SPD(N_\theta)$, $Y\in\Sym(N_\theta)$, and let $V\in C^2$ satisfy
$\norm{\nabla^2V}_\op\leq\beta$. Write $a_\epsilon:=(m,C+\epsilon Y)$. Then
\begin{equation}
\label{eq:price-directional}
    \frac{\dd}{\dd\epsilon}\bigg|_{\epsilon=0}\E_{\rho_{a_\epsilon}}[V]
    =\frac12\Tr\bigl(Y\,A(a)\bigr),
    \qquad
    \frac{\dd}{\dd\epsilon}\bigg|_{\epsilon=0}\calE_1(a_\epsilon)
    =-\frac12\Tr\bigl(YC^{-1}\bigr),
\end{equation}
and consequently
\begin{equation}
\label{eq:E-directional}
    \frac{\dd}{\dd\epsilon}\bigg|_{\epsilon=0}\calE(a_\epsilon)
    =\frac12\Tr\Bigl(Y\bigl(A(a)-C^{-1}\bigr)\Bigr)
    =-\frac12\Tr\bigl(Y R(a)\bigr).
\end{equation}
\end{lemma}

\begin{proof}
Let $p_\epsilon$ be the density of $\N(m,C+\epsilon Y)$; for $|\epsilon|$ small,
$C+\epsilon Y\succ0$ and $p_\epsilon$ is Schwartz. Its Fourier transform is
$\widehat p_\epsilon(\xi)=\exp\bigl(i\,m^\top\xi-\tfrac12\xi^\top(C+\epsilon Y)\xi\bigr)$, so
\[
    \partial_\epsilon\widehat p_\epsilon(\xi)\big|_{\epsilon=0}
    =-\tfrac12\bigl(\xi^\top Y\xi\bigr)\widehat p_0(\xi).
\]
The Fourier transform of $\partial_i\partial_jp_0$ is $-\xi_i\xi_j\widehat p_0$, hence the Fourier
transform of
\[
    \tfrac12\Tr\bigl(Y\nabla_x^2p_0\bigr)=\tfrac12\sum_{i,j}Y_{ij}\partial_i\partial_jp_0
\]
is exactly $-\tfrac12(\xi^\top Y\xi)\widehat p_0(\xi)$. By Fourier inversion,
\begin{equation}
\label{eq:heat-identity}
    \partial_\epsilon p_\epsilon\big|_{\epsilon=0}=\tfrac12\Tr\bigl(Y\,\nabla_x^2p_0\bigr).
\end{equation}
Since $\norm{\nabla^2V}_\op\leq\beta$, $V$ grows at most quadratically, while $p_\epsilon$ and its
first two $x$-derivatives decay faster than any polynomial, uniformly for $\epsilon$ near $0$;
this justifies differentiating under the integral sign and integrating by parts twice with
vanishing boundary terms. Therefore
\[
    \frac{\dd}{\dd\epsilon}\bigg|_{0}\int V\,p_\epsilon\,\dd x
    =\int V\cdot\tfrac12\Tr\bigl(Y\nabla_x^2p_0\bigr)\dd x
    =\tfrac12\int\Tr\bigl(Y\,\nabla^2V\bigr)p_0\,\dd x
    =\tfrac12\Tr\bigl(Y\,A(a)\bigr).
\]
For the entropy term, Jacobi's formula gives
$\tfrac{\dd}{\dd\epsilon}\log\det(C+\epsilon Y)\big|_0=\Tr(C^{-1}Y)$, and
$\calE_1(a_\epsilon)=-\tfrac12\log\det(C+\epsilon Y)+\text{const.}$ Adding the two contributions
and using \eqref{eq:splitting} gives \eqref{eq:E-directional}.
\end{proof}

\begin{lemma}[Fisher--Rao gradient, flow, and gradient norm]
\label{lem:FR-gradient}
Under \Cref{def:FR-metric},
\begin{equation}
\label{eq:FR-grad}
    \grad\calE(a)=\bigl(-Cg(a),\,-CR(a)C\bigr),
\end{equation}
so that the Fisher--Rao (natural) gradient flow $\dot a_t=-\grad\calE(a_t)$ is
\begin{equation}
\label{eq:NG-flow}
    \frac{\dd m_t}{\dd t}=C_tg(a_t),
    \qquad
    \frac{\dd C_t}{\dd t}=C_t-C_tA(a_t)C_t .
\end{equation}
Moreover, defining the Fisher--Rao gradient norm
\begin{equation}
\label{eq:G-norm}
    \Gradnorm(a)^2:=g(a)^\top Cg(a)+\frac12\norm{C^{1/2}R(a)C^{1/2}}_{\Frob}^2 ,
\end{equation}
one has $\norm{\grad\calE(a)}_a^2=\Gradnorm(a)^2$ and, along \eqref{eq:NG-flow},
$\tfrac{\dd}{\dd t}\Delta(a_t)=-\Gradnorm(a_t)^2$.
\end{lemma}

\begin{proof}
Fix a tangent vector $(v,Y)$. Since $\calE_1$ does not depend on $m$ and
$\tfrac{\dd}{\dd\epsilon}\big|_0\E_{\rho_{(m+\epsilon v,C)}}[V]=\E_{\rho_a}[\nabla V]^\top v=-g(a)^\top v$,
while the covariance direction is given by \Cref{lem:price-derivative}, the differential of
$\calE$ at $a$ is
\begin{equation}
\label{eq:dE}
    \dd\calE(a)[(v,Y)]=-g(a)^\top v-\tfrac12\Tr\bigl(R(a)Y\bigr).
\end{equation}
The gradient $(u^\star,U^\star):=\grad\calE(a)$ is defined by
$g_a\bigl((u^\star,U^\star),(v,Y)\bigr)=\dd\calE(a)[(v,Y)]$ for all $(v,Y)$. Matching the mean
block, $u^{\star\top}C^{-1}v=-g(a)^\top v$ for all $v$, so $u^\star=-Cg(a)$. Matching the
covariance block,
$\tfrac12\Tr(C^{-1}U^\star C^{-1}Y)=-\tfrac12\Tr(R(a)Y)$ for all $Y\in\Sym(N_\theta)$; both
$C^{-1}U^\star C^{-1}$ and $R(a)$ are symmetric and $(S,Y)\mapsto\Tr(SY)$ is nondegenerate on
$\Sym(N_\theta)$, so $C^{-1}U^\star C^{-1}=-R(a)$, i.e.\ $U^\star=-CR(a)C$. This is
\eqref{eq:FR-grad}. Then $\dot a=-\grad\calE(a)$ gives $\dot m=Cg(a)$ and
$\dot C=CR(a)C=C(C^{-1}-A(a))C=C-CA(a)C$, which is \eqref{eq:NG-flow}. Finally, with
$u^\star=-Cg$, $U^\star=-CRC$,
\[
    \norm{\grad\calE(a)}_a^2
    =(Cg)^\top C^{-1}(Cg)+\frac12\norm{C^{-1/2}(CRC)C^{-1/2}}_{\Frob}^2
    =g^\top Cg+\frac12\norm{C^{1/2}RC^{1/2}}_{\Frob}^2
    =\Gradnorm(a)^2 .
\]
The dissipation identity follows from the chain rule,
$\tfrac{\dd}{\dd t}\Delta(a_t)=\dd\calE(a_t)[-\grad\calE(a_t)]=-\norm{\grad\calE(a_t)}_{a_t}^2$.
\end{proof}

\subsection{Deterministic schemes}

The deterministic \emph{Riemannian retraction scheme} (covariance-exponential scheme) is
\begin{equation}
\label{eq:Riem-det}
    m_{n+1}=m_n+\Delta t\,C_ng(a_n),
    \qquad
    C_{n+1}=C_n^{1/2}\exp\bigl(\Delta t(\Id-C_n^{1/2}A(a_n)C_n^{1/2})\bigr)C_n^{1/2},
\end{equation}
and the deterministic \emph{KL/Bregman scheme} is
\begin{equation}
\label{eq:KL-det}
    m_{n+1}=m_n+\Delta t\,C_ng(a_n),
    \qquad
    C_{n+1}=(1+\Delta t)\bigl(C_n^{-1}+\Delta t\,A(a_n)\bigr)^{-1},
\end{equation}
equivalently $C_{n+1}^{-1}=\tfrac1{1+\Delta t}\bigl(C_n^{-1}+\Delta t\,A(a_n)\bigr)$.

\begin{remark}[Terminology: retraction, not geodesic]
\label{rem:not-geodesic}
The covariance block $C\mapsto C^{1/2}\exp(\cdot)C^{1/2}$ is the exponential map of the
affine-invariant metric on $\SPD(N_\theta)$, but \eqref{eq:Riem-det} pairs it with a Euclidean
step in the mean, and we do not claim that the resulting map is the exponential map of the full
Gaussian Fisher--Rao metric \eqref{eq:FR-metric}. All statements are proved for the retraction
\eqref{eq:retraction} that the algorithms use; we therefore say \emph{Riemannian retraction
scheme} rather than ``geodesic scheme''.
\end{remark}

\section{Two basic inequalities}
\label{sec:basic}

\begin{lemma}[Bures--Wasserstein PL inequality on Gaussians]
\label{lem:W2-PL}
Under \Cref{assump:logconcave-smooth}, for every $a\in\Aspace$,
\[
    \norm{\grad^{W_2}\calE(a)}_{W_2}^2
    =\norm{g(a)}^2+\Tr\bigl(R(a)CR(a)\bigr)
    \geq2\alpha\,\Delta(a),
\]
where $\grad^{W_2}\calE(a)=(-g(a),-R(a))$ is the Gaussian-restricted Bures--Wasserstein gradient.
\end{lemma}

This is the Polyak--{\L}ojasiewicz inequality for the restricted gradient, following from the
$\alpha$-geodesic convexity of $\calE$ along Bures--Wasserstein geodesics between Gaussians; see
\citet{lambert2022variational,diao2023forward}. Here $(-g(a),-R(a))$ denotes the mean-velocity/dual
(potential) representative of the Bures--Wasserstein gradient, in which the squared metric norm takes
the displayed form; the covariance-coordinate tangent of the corresponding gradient flow is
$-(R(a)C+CR(a))$, cf.\ \eqref{eq:W-flow}, not $-R(a)$ itself.

\begin{lemma}[Fisher--Rao/Wasserstein comparison]
\label{lem:FR-W2}
Let $a=(m,C)$ with $C\succeq\ell\Id$, $\ell>0$. Then
\begin{equation}
\label{eq:FR-W2}
    \boxed{\;
    \Gradnorm(a)^2
    \;\geq\;
    \frac{\ell}{2}\norm{\grad^{W_2}\calE(a)}_{W_2}^2
    \;\geq\;
    \alpha\ell\,\Delta(a).
    \;}
\end{equation}
Moreover $\Gradnorm(a)^2\leq\norm{C}_\op\norm{\grad^{W_2}\calE(a)}_{W_2}^2$.
\end{lemma}

\begin{proof}
Using $\norm{C^{1/2}R(a)C^{1/2}}_{\Frob}^2=\Tr(CR(a)CR(a))$, \eqref{eq:G-norm} reads
$\Gradnorm(a)^2=g^\top Cg+\tfrac12\Tr(CRCR)$. If $C\succeq\ell\Id$ then
$g^\top Cg\geq\ell\norm{g}^2$. Also $R CR\succeq0$, and for $N\succeq0$,
$\Tr(CN)=\Tr(N^{1/2}CN^{1/2})\geq\ell\Tr(N)$; applying this with $N=RCR$,
$\tfrac12\Tr(CRCR)\geq\tfrac\ell2\Tr(RCR)$. Hence
\[
    \Gradnorm(a)^2\geq\ell\norm{g}^2+\frac\ell2\Tr(RCR)
    \geq\frac\ell2\Bigl(\norm{g}^2+\Tr(RCR)\Bigr)
    =\frac\ell2\norm{\grad^{W_2}\calE(a)}_{W_2}^2 ,
\]
and \Cref{lem:W2-PL} gives $\Gradnorm(a)^2\geq\tfrac\ell2\cdot2\alpha\Delta(a)=\alpha\ell\Delta(a)$.
For the upper bound, with $L_C=\norm{C}_\op$, $g^\top Cg\leq L_C\norm{g}^2$ and
$\tfrac12\Tr(CRCR)\leq\tfrac{L_C}2\Tr(RCR)\leq L_C\Tr(RCR)$.
\end{proof}

\section{Continuous-time covariance bootstrap}
\label{sec:cont}

\begin{lemma}[Precision Duhamel formula]
\label{lem:precision-duhamel}
Let $P_t:=C_t^{-1}$. Along \eqref{eq:NG-flow},
$\dot P_t=A(a_t)-P_t$, hence $P_t=e^{-t}P_0+\int_0^te^{-(t-s)}A(a_s)\,\dd s$.
\end{lemma}

\begin{proof}
$\dot P_t=-P_t\dot C_tP_t=-P_t(C_t-C_tA(a_t)C_t)P_t=-P_t+A(a_t)$. Multiplying
$\dot P_t+P_t=A(a_t)$ by $e^t$ and integrating gives the formula.
\end{proof}

\begin{lemma}[Continuous covariance bootstrap]
\label{lem:cont-cov}
Under \Cref{assump:logconcave-smooth}, along \eqref{eq:NG-flow},
\begin{equation}
\label{eq:ell-t}
    C_t\succeq\ell(t)\Id,
    \qquad
    \ell(t)=\frac{\lambda_0e^t}{1+\beta\lambda_0(e^t-1)},
    \qquad\text{and}\qquad
    C_t\preceq\Lambda\Id .
\end{equation}
In particular $C_t\succeq\tfrac1{2\beta}\Id$ for all $t\geq\logp\frac1{\beta\lambda_0}$.
\end{lemma}

\begin{proof}
By \Cref{lem:precision-duhamel}, $A(a_s)\preceq\beta\Id$ and $P_0\preceq\lambda_0^{-1}\Id$,
\[
    P_t\preceq\bigl(e^{-t}\lambda_0^{-1}+\beta(1-e^{-t})\bigr)\Id ,
\]
and inverting gives $C_t\succeq\ell(t)\Id$ with
$\ell(t)=\lambda_0e^t/\bigl(1+\beta\lambda_0(e^t-1)\bigr)$. For the upper bound, if $\mu_i(t)$ is
an eigenvalue of $C_t$ with unit eigenvector $u_i(t)$, then $\dot\mu_i=\mu_i(1-\mu_ih_i)$ with
$h_i=u_i^\top A(a_t)u_i\in[\alpha,\beta]$, so $\dot\mu_i\leq0$ whenever $\mu_i\geq1/\alpha$; no
eigenvalue crosses $1/\alpha$ from below, whence $C_t\preceq\Lambda\Id$. (At repeated eigenvalues,
here and in all later band arguments, this is read via the Dini derivative of
$t\mapsto\lambda_{\max}(C_t)$ along a unit vector in the extremal eigenspace, which satisfies the
same scalar differential inequality; eigenvector differentiability is not needed.) Finally, if
$e^{-t}\lambda_0^{-1}\leq\beta$, then $e^{-t}\lambda_0^{-1}+\beta(1-e^{-t})\leq2\beta$.
\end{proof}

\begin{theorem}[Continuous-time global convergence; deterministic]
\label{thm:cont-global}
Under \Cref{assump:logconcave-smooth}, the flow \eqref{eq:NG-flow} satisfies
\[
    \boxed{\;
    \Delta(a_t)\leq\Delta_0\bigl(1+\beta\lambda_0(e^t-1)\bigr)^{-1/\kappa}.
    \;}
\]
Consequently, for $\delta\in(0,\Delta_0)$, the hitting time
$T_{\mathrm c}^{\glob}(\delta):=\inf\{t\geq0:\Delta(a_t)\leq\delta\}$ obeys
\[
    T_{\mathrm c}^{\glob}(\delta)
    \leq\log\!\left(1+\frac{(\Delta_0/\delta)^{\kappa}-1}{\beta\lambda_0}\right)
    =O\!\left(\logp\frac1{\beta\lambda_0}+\kappa\logp\frac{\Delta_0}{\delta}\right).
\]
\end{theorem}

\begin{proof}
By \Cref{lem:FR-gradient}, $\tfrac{\dd}{\dd t}\Delta(a_t)=-\Gradnorm(a_t)^2$. By
\Cref{lem:cont-cov}, $C_t\succeq\ell(t)\Id$, so \Cref{lem:FR-W2} gives
$\tfrac{\dd}{\dd t}\Delta(a_t)\leq-\alpha\ell(t)\Delta(a_t)$, and Gr\"onwall yields
$\Delta(a_t)\leq\Delta_0\exp\bigl(-\alpha\int_0^t\ell(s)\dd s\bigr)$. Substituting $v=e^s$,
\[
    \int_0^t\ell(s)\,\dd s
    =\int_0^t\frac{\lambda_0e^s}{1+\beta\lambda_0(e^s-1)}\,\dd s
    =\frac1\beta\Bigl[\log\bigl(1+\beta\lambda_0(e^s-1)\bigr)\Bigr]_0^t
    =\frac1\beta\log\bigl(1+\beta\lambda_0(e^t-1)\bigr),
\]
so $\Delta(a_t)\leq\Delta_0(1+\beta\lambda_0(e^t-1))^{-\alpha/\beta}$, which is the claim since
$\alpha/\beta=1/\kappa$. Solving
$\Delta_0(1+\beta\lambda_0(e^T-1))^{-1/\kappa}=\delta$ gives
$e^T=1+\bigl((\Delta_0/\delta)^\kappa-1\bigr)/(\beta\lambda_0)$. The big-$O$ form follows from
$\log(1+xy)\leq\logp x+\logp y+\log2$ for $x,y\geq0$ and
$\log\bigl((\Delta_0/\delta)^\kappa\bigr)=\kappa\log(\Delta_0/\delta)$.
\end{proof}

\begin{proposition}[Order-optimality of the covariance burn-in]
\label{prop:burnin-optimal}
Let $N_\theta=1$, $V(\theta)=\tfrac12\theta^2$ (so $\alpha=\beta=1$), $m_0=0$, and
$C_0=\lambda_0\in(0,\tfrac12)$. Then the flow \eqref{eq:NG-flow} has the exact solution
$c_t=\lambda_0e^t/\bigl(1+\lambda_0(e^t-1)\bigr)$---the envelope \eqref{eq:ell-t} with equality---so
$c_t\geq\tfrac12$ if and only if $t\geq\log\tfrac{1-\lambda_0}{\lambda_0}$, and for every $\delta>0$
with lower band root $\ell_\delta<1$ (\Cref{cor:band}, $\varphi(\ell_\delta)=\delta$),
\[
    \Delta(a_t)\leq\delta
    \quad\Longrightarrow\quad
    t\;\geq\;\log\frac1{\lambda_0}-\log\frac{1-\ell_\delta}{\ell_\delta}-\log2 .
\]
Hence the $\logp(1/(\beta\lambda_0))$ burn-in stage of \Cref{thm:cont-global} is order-optimal for
pure Fisher--Rao dynamics: it cannot be removed or improved beyond constants without changing the
dynamics (which is exactly what the initializers of \Cref{sec:transport-bootstrap} do).
\end{proposition}

\begin{proof}
Here $g\equiv-m$, $A\equiv1$, so $m_t\equiv0$ and $\dot c_t=c_t-c_t^2$, the logistic equation, with
the stated exact solution. Monotonicity gives $c_t<1$ for all $t$, and by
\Cref{lem:entropy-coerc} (an equality for Gaussian targets) $\Delta(a_t)=\varphi(c_t)$, so
$\Delta(a_t)\leq\delta$ forces $c_t\geq\ell_\delta$. Solving
$\lambda_0e^t\geq\ell_\delta\bigl(1+\lambda_0(e^t-1)\bigr)$ gives
$e^t\geq\ell_\delta(1-\lambda_0)/\bigl(\lambda_0(1-\ell_\delta)\bigr)$, and $1-\lambda_0\geq\tfrac12$
yields the display.
\end{proof}

\section{The Riemannian one-step descent lemma}
\label{sec:descent}

Throughout this section $a=(m,C)$ is fixed and
\begin{equation}
\label{eq:S-def}
    S:=\Id-C^{1/2}A(a)C^{1/2}=C^{1/2}R(a)C^{1/2},
\end{equation}
so that, by \eqref{eq:G-norm}, $\Gradnorm(a)^2=g(a)^\top Cg(a)+\tfrac12\norm{S}_{\Frob}^2$. We write
the retraction
\begin{equation}
\label{eq:retraction}
    \Retr_a(u,U):=\Bigl(m+u,\;C^{1/2}\exp\bigl(C^{-1/2}UC^{-1/2}\bigr)C^{1/2}\Bigr),
\end{equation}
which reproduces \eqref{eq:Riem-det} with stepsize $h$ on taking
$(u,U)=-h\,\grad\calE(a)$, that is $u=hCg$ and $U=hCRC$,
since $C^{-1/2}(hCRC)C^{-1/2}=hS$.

\subsection{Two scalar inequalities}

\begin{lemma}[Scalar inequality, deterministic step]
\label{lem:scalar-det}
Let $K\geq1$ and $0<h\leq1/K$. Then for every $s\in[1-K,1]$,
\begin{equation}
\label{eq:scalar-det}
    \bigl(e^{hs/2}-1\bigr)(1-s)+\frac K2\bigl(e^{hs/2}-1\bigr)^2-\frac h2s
    \;\leq\;
    -\frac h2\Bigl(1-\frac{Kh}2\Bigr)s^2 .
\end{equation}
\end{lemma}

\begin{proof}
Put $u:=hs/2$ and
\[
    \Xi(s):=\bigl(e^u-1\bigr)(1-s)+\frac K2\bigl(e^u-1\bigr)^2-\frac h2s+\frac h2\Bigl(1-\frac{Kh}2\Bigr)s^2 ,
\]
so \eqref{eq:scalar-det} is $\Xi(s)\leq0$.

\emph{Step 0 (range of $u$).} Since $s\in[1-K,1]$ and $0<h\leq1/K$ with $K\geq1$,
$u\leq\tfrac h2\leq\tfrac1{2K}\leq\tfrac12$ and
$u\geq\tfrac{h(1-K)}2\geq\tfrac{1-K}{2K}=\tfrac1{2K}-\tfrac12\geq-\tfrac12$, so
$u\in[-\tfrac12,\tfrac12]$, with the same sign as $s$.

\emph{Step 1 (exact reduction).} Note $\tfrac h2s=u$, $\tfrac h2s^2=us$, and
$\tfrac h2\cdot\tfrac{Kh}2s^2=Ku^2$. With $D(u):=e^u-1-u\geq0$ and $w:=e^u-1$ we have
$u-w=-D(u)$, hence
\[
    \Xi(s)=D(u)+s(u-w)+\frac K2w^2-Ku^2
    =(1-s)D(u)+\frac K2\Bigl[(e^u-1)^2-2u^2\Bigr].
\]
Setting $P(u):=2u^2-(e^u-1)^2$, the claim is equivalent to
\begin{equation}
\label{eq:scalar-reduced}
    (1-s)\,D(u)\leq\frac K2\,P(u),\qquad u=\tfrac{hs}2 .
\end{equation}

\emph{Step 2 ($P\geq0$ on $[-\tfrac12,\tfrac12]$).} For $u\leq0$, $|e^u-1|\leq|u|$, so
$(e^u-1)^2\leq u^2\leq2u^2$. For $u\in[0,\tfrac12]$, convexity of $u\mapsto e^u-1$ and the chord
bound give $e^u-1\leq2(e^{1/2}-1)u\leq\sqrt2\,u$, whence $(e^u-1)^2\leq2u^2$.

\emph{Step 3 (case $s\leq0$).} Then $u\leq0$, and $s\geq1-K$ gives $K\geq1-s>0$. Since $P(u)\geq0$,
$\tfrac K2P(u)\geq\tfrac{1-s}2P(u)$, so it suffices that $2D(u)\leq P(u)$ on $[-\tfrac12,0]$.
Let $Q(u):=P(u)-2D(u)=2u^2-(e^u-1)^2-2e^u+2+2u$. Then $Q(0)=0$,
$Q'(u)=4u-2e^{2u}+2$ with $Q'(0)=0$, and $Q''(u)=4-4e^{2u}\geq0$ for $u\leq0$. So $Q'$ is
nondecreasing on $[-\tfrac12,0]$ with $Q'(0)=0$, hence $Q'\leq0$ there, hence $Q$ is
nonincreasing and $Q(u)\geq Q(0)=0$.

\emph{Step 4 (case $s>0$).} Then $u>0$, and $h\leq1/K$ gives $s=2u/h\geq2uK$; combined with $s\leq1$
this forces $2uK\leq1$, i.e.\ $1-2uK\geq0$ (the branch $1-2uK<0$ cannot occur). Multiplying
$1-s\leq1-2uK$ by $D(u)\geq0$ gives $(1-s)D(u)\leq(1-2uK)D(u)$, so to establish
\eqref{eq:scalar-reduced} it suffices to prove
$(1-2uK)D(u)\leq\tfrac K2P(u)$. Rearranged, this is $f(K)\geq0$ with
$f(K):=K\bigl(\tfrac{P(u)}2+2uD(u)\bigr)-D(u)$. The coefficient of $K$ is nonnegative by Step 2
and $u,D(u)\geq0$, so $f$ is nondecreasing; since $K\geq1$ it suffices that $f(1)\geq0$, i.e.
\[
    G(u):=2u^2-(e^u-1)^2+(4u-2)\bigl(e^u-1-u\bigr)\geq0,\qquad u\in[0,\tfrac12].
\]
Write $G(u)=-2g(u)$ with $g(u):=\tfrac12e^{2u}-2ue^{u}+u^2+u-\tfrac12$. Now $g(0)=0$,
$g'(u)=e^{2u}-2e^u-2ue^u+2u+1$ with $g'(0)=0$, $g''(u)=2e^{2u}-4e^u-2ue^u+2$ with $g''(0)=0$, and
$g'''(u)=4e^{2u}-6e^u-2ue^u=2e^u\bigl(2e^u-3-u\bigr)$.
On $[0,\tfrac12]$ the map $u\mapsto2e^u-3-u$ has derivative $2e^u-1>0$, hence is increasing, with
maximal value $2e^{1/2}-3.5=-0.2026\ldots<0$. Thus $g'''<0$, so $g''$ is strictly decreasing with
$g''(0)=0$, hence $g''\leq0$; so $g'\leq0$; so $g\leq0$ on $[0,\tfrac12]$, giving $G\geq0$.
\end{proof}

\begin{lemma}[Scalar inequality, arbitrary retraction direction]
\label{lem:scalar-stoch}
For every $w\in[-1,1]$,
\begin{equation}
\label{eq:scalar-stoch}
    \Bigl(e^{w/2}-1-\frac w2\Bigr)+\frac12\bigl(e^{w/2}-1\bigr)^2\;\leq\;\frac{w^2}2 .
\end{equation}
\end{lemma}

\begin{proof}
Put $v:=w/2\in[-\tfrac12,\tfrac12]$. Since
$(e^v-1)+\tfrac12(e^v-1)^2=\tfrac12(e^{2v}-1)$, the left side equals $\tfrac12(e^{2v}-1)-v$, and
the claim is $\tfrac12(e^{2v}-1)-v\leq2v^2$, i.e.\ $e^x\leq1+x+x^2$ with $x:=2v\in[-1,1]$. This
holds because
$e^x-1-x=\sum_{k\geq2}\tfrac{x^k}{k!}\leq x^2\sum_{k\geq2}\tfrac{|x|^{k-2}}{k!}\leq(e-2)x^2<x^2$
for $|x|\leq1$.
\end{proof}

\subsection{Deterministic one-step descent}

\begin{lemma}[Deterministic Riemannian one-step descent]
\label{lem:det-descent}
Assume $0\preceq\nabla^2V\preceq\beta\Id$ and $C\preceq\Lambda\Id$; set $K:=\beta\Lambda$ and
assume $K\geq1$. Let $0<h\leq1/K$ and
$m^+:=m+hCg(a)$, $C^+:=C^{1/2}\exp(hS)C^{1/2}$, $a^+:=(m^+,C^+)$. Then
\begin{equation}
\label{eq:det-descent}
    \boxed{\;
    \calE(a^+)\;\leq\;\calE(a)-h\Bigl(1-\frac{Kh}2\Bigr)\Gradnorm(a)^2 .
    \;}
\end{equation}
\end{lemma}

\begin{proof}
\emph{Step 1 (coupling).} Let $\xi\sim\N(0,\Id)$ and set $X:=m+C^{1/2}\xi\sim\rho_a$ and
$Y:=m+hCg(a)+C^{1/2}e^{hS/2}\xi$. Since $e^{hS/2}$ is symmetric, $Y$ is Gaussian with mean $m^+$
and covariance $C^{1/2}e^{hS/2}e^{hS/2}C^{1/2}=C^+$, i.e.\ $Y\sim\rho_{a^+}$.

\emph{Step 2 (exact entropy change).} By \eqref{eq:splitting},
$\calE_1(a^+)-\calE_1(a)=-\tfrac12\log\det e^{hS}=-\frac h2\Tr S$.

\emph{Step 3 ($\beta$-smoothness).} $\calE_2(a^+)-\calE_2(a)=\E[V(Y)-V(X)]$, and
$V(Y)\leq V(X)+\nabla V(X)^\top(Y-X)+\tfrac\beta2\norm{Y-X}^2$ with
$Y-X=hCg(a)+C^{1/2}E\xi$, $E:=e^{hS/2}-\Id$.
By $\E[\nabla V(X)]=-g(a)$, $\E[\nabla V(X)^\top hCg(a)]=-h\,g^\top Cg$. By Price's identity
$\E[\xi\,\nabla V(X)^\top]=C^{1/2}A(a)$ (\Cref{rem:price}),
\[
    \E\bigl[\nabla V(X)^\top C^{1/2}E\xi\bigr]
    =\Tr\bigl(E\,C^{1/2}A(a)C^{1/2}\bigr)
    =\Tr\bigl(E(\Id-S)\bigr).
\]
Since $\E[\xi]=0$ the cross term vanishes and
$\tfrac\beta2\E\norm{Y-X}^2=\tfrac\beta2h^2g^\top C^2g+\tfrac\beta2\Tr(ECE)$. With
$C\preceq\Lambda\Id$, $g^\top C^2g\leq\Lambda g^\top Cg$ and $\Tr(ECE)\leq\Lambda\norm{E}_{\Frob}^2$,
so $\tfrac\beta2\E\norm{Y-X}^2\leq\tfrac K2h^2g^\top Cg+\tfrac K2\norm{E}_{\Frob}^2$.

\emph{Step 4 (assembly).} Adding Steps 2 and 3,
\[
    \calE(a^+)-\calE(a)
    \leq
    \underbrace{-h\,g^\top Cg+\frac K2h^2g^\top Cg}_{M}
    +\underbrace{-\frac h2\Tr S+\Tr\bigl(E(\Id-S)\bigr)+\frac K2\norm{E}_{\Frob}^2}_{\Phi},
\]
and $M=-h\bigl(1-\tfrac{Kh}2\bigr)g^\top Cg$.

\emph{Step 5 (covariance part).} Let $S=\sum_is_iv_iv_i^\top$. Since $0\preceq A(a)\preceq\beta\Id$
and $C\preceq\Lambda\Id$, $0\preceq C^{1/2}A(a)C^{1/2}\preceq\beta C\preceq K\Id$, so
$s_i\in[1-K,1]$. As $E$ and $\Id-S$ share the eigenbasis $\{v_i\}$,
$\Phi=\sum_i\phi(s_i)$ with
$\phi(s):=-\tfrac h2s+\bigl(e^{hs/2}-1\bigr)(1-s)+\tfrac K2\bigl(e^{hs/2}-1\bigr)^2$.
By \Cref{lem:scalar-det}, $\phi(s_i)\leq-\tfrac h2\bigl(1-\tfrac{Kh}2\bigr)s_i^2$, so
$\Phi\leq-h\bigl(1-\tfrac{Kh}2\bigr)\cdot\tfrac12\norm{S}_{\Frob}^2$.

\emph{Step 6.} Adding $M$ and the bound on $\Phi$ gives \eqref{eq:det-descent}.
\end{proof}

\subsection{Retraction smoothness for arbitrary directions}

\begin{lemma}[Retraction smoothness, constant $2\beta\Lambda$]
\label{lem:retr-smooth}
Assume $0\preceq\nabla^2V\preceq\beta\Id$ and $C\preceq\Lambda\Id$; set $K:=\beta\Lambda$ and
$L:=2K$. Let $\zeta=(u,U)$ be any tangent vector at $a$ with
\begin{equation}
\label{eq:scope}
    \norm{W}_\op\leq1,\qquad W:=C^{-1/2}UC^{-1/2}.
\end{equation}
Then
\begin{equation}
\label{eq:retr-smooth}
    \boxed{\;
    \calE\bigl(\Retr_a(\zeta)\bigr)
    \;\leq\;
    \calE(a)+\inner{\grad\calE(a)}{\zeta}_a+\frac L2\norm{\zeta}_a^2 .
    \;}
\end{equation}
\end{lemma}

\begin{proof}
\emph{Step 1.} Let $\xi\sim\N(0,\Id)$, $X=m+C^{1/2}\xi\sim\rho_a$ and $Y:=m+u+C^{1/2}e^{W/2}\xi$,
so $Y\sim\rho_{a^+}$ with $a^+=\Retr_a(\zeta)$.

\emph{Step 2.} As before $\calE_1(a^+)-\calE_1(a)=-\tfrac12\Tr W$.

\emph{Step 3.} With $E:=e^{W/2}-\Id$ and $Y-X=u+C^{1/2}E\xi$, $\beta$-smoothness,
$\E[\nabla V(X)]=-g(a)$, Price's identity and $\E[\xi]=0$ give
\[
    \calE_2(a^+)-\calE_2(a)
    \leq-g(a)^\top u+\Tr\bigl(E\,C^{1/2}A(a)C^{1/2}\bigr)
    +\frac\beta2\norm{u}^2+\frac\beta2\Tr(ECE),
\]
and $\tfrac\beta2\Tr(ECE)\leq\tfrac K2\norm{E}_{\Frob}^2$ since $C\preceq\Lambda\Id$.

\emph{Step 4.} By \Cref{lem:FR-gradient} and \eqref{eq:FR-metric},
$\inner{\grad\calE(a)}{\zeta}_a=-g^\top u-\tfrac12\Tr(SW)$ and
$\norm{\zeta}_a^2=u^\top C^{-1}u+\tfrac12\norm{W}_{\Frob}^2$, using
$\Tr(RU)=\Tr(C^{1/2}RC^{1/2}W)=\Tr(SW)$.

\emph{Step 5 (mean block).} $C\preceq\Lambda\Id$ gives $C^{-1}\succeq\Lambda^{-1}\Id$, so
$\tfrac\beta2\norm{u}^2\leq\tfrac K2u^\top C^{-1}u\leq\tfrac L2u^\top C^{-1}u$.

\emph{Step 6 (covariance block).} The remaining terms are
$\Theta:=-\tfrac12\Tr W+\Tr\bigl(E(\Id-S)\bigr)+\tfrac K2\norm{E}_{\Frob}^2+\tfrac12\Tr(SW)$, and
we must show $\Theta\leq\tfrac L4\norm{W}_{\Frob}^2$. Using
$\Tr(E)-\Tr(ES)+\tfrac12\Tr(SW)-\tfrac12\Tr(W)=\Tr[(E-\tfrac W2)(\Id-S)]$, we obtain the exact
identity
$\Theta=\Tr\bigl[(e^{W/2}-\Id-\tfrac W2)(\Id-S)\bigr]+\tfrac K2\norm{E}_{\Frob}^2$.
Put $N:=e^{W/2}-\Id-\tfrac W2\succeq0$ and $M:=\Id-S=C^{1/2}A(a)C^{1/2}$, so $0\preceq M\preceq K\Id$.
Then $\Tr(NM)=\Tr(N^{1/2}MN^{1/2})\leq K\Tr(N)$, since $N^{1/2}MN^{1/2}\preceq KN$; no commutation
between $W$ and $S$ is used, which is the source of the extra factor $2$ in $L$. Diagonalizing
$W=\sum_iw_iq_iq_i^\top$ with $|w_i|\leq1$ by \eqref{eq:scope} and applying \Cref{lem:scalar-stoch},
\[
    \Theta\leq K\sum_i\left[\Bigl(e^{w_i/2}-1-\frac{w_i}2\Bigr)+\frac12\bigl(e^{w_i/2}-1\bigr)^2\right]
    \leq K\sum_i\frac{w_i^2}2=\frac K2\norm{W}_{\Frob}^2=\frac L4\norm{W}_{\Frob}^2 .
\]
Combining Steps 5 and 6 gives \eqref{eq:retr-smooth}.
\end{proof}

\begin{remark}[Sharper retraction constant]
\label{rem:retr-sharp}
The left side of \eqref{eq:scalar-stoch} simplifies exactly to $\tfrac12(e^w-1-w)$, and
$e^w-1-w\leq(e-2)w^2$ on $[-1,1]$, so \Cref{lem:scalar-stoch} holds with the right side
$\tfrac{e-2}2w^2$ and \Cref{lem:retr-smooth} with $L=2(e-2)\beta\Lambda\approx1.437\,\beta\Lambda$
(the mean block still satisfies Step~5 since $2(e-2)>1$); a radius-dependent refinement is
$L(r)=2\beta\Lambda\,(e^r-1-r)/r^2$ under $\norm W_\op\leq r$. The stochastic one-step estimates below
retain the conservative $L=2\beta\Lambda$; the improvement rescales stepsizes and floors by the
constant $e-2\approx0.718$ without affecting any order.
\end{remark}

\section{Riemannian retraction scheme: deterministic global theory}
\label{sec:Riem-global}

\begin{lemma}[Riemannian covariance bootstrap]
\label{lem:Riem-cov}
Assume \Cref{assump:logconcave-smooth}, $\lambda_{0,\min}\Id\preceq C_0\preceq\lambda_{0,\max}\Id$
and $0<\Delta t\leq1/\LB$. Then $C_n\preceq\Lambda\Id$ for all $n\geq0$. Define
\begin{equation}
\label{eq:Riem-envelope}
    L_0:=\min\Bigl\{\lambda_0,\frac1{2\beta}\Bigr\},
    \qquad
    L_{n+1}:=\frac{e^{\Delta t}L_n}{1+(e-1)\Delta t\,\beta L_n}.
\end{equation}
Then $C_n\succeq L_n\Id$ for all $n\geq0$, and with $y_n:=(\beta L_n)^{-1}$,
\begin{equation}
\label{eq:Riem-y}
    y_n=y_\infty+(y_0-y_\infty)e^{-n\Delta t},
    \qquad
    y_\infty:=\frac{(e-1)\Delta t}{e^{\Delta t}-1}\in[1,e-1)\subset[1,2).
\end{equation}
Consequently $C_n\succeq\tfrac1{2\beta}\Id$ for all
\begin{equation}
\label{eq:Ncov-minus}
    n\geq N_{\cov,-}^{\Riem}
    :=\left\lceil\frac1{\Delta t}\logp\frac{y_0-y_\infty}{2-y_\infty}\right\rceil
    =O\!\left(\frac1{\Delta t}\logp\frac1{\beta\lambda_0}\right),
\end{equation}
and, with $z_0:=\tfrac1{\alpha\lambda_{0,\max}}$ and $z_\infty:=\tfrac{\Delta t}{e^{\Delta t}-1}\in[\tfrac1{e-1},1)$,
$C_n\preceq\tfrac2\alpha\Id$ for all $n\geq N_{\cov,+}^{\Riem}=O(1/\Delta t)$, where
$N_{\cov,+}^{\Riem}=0$ if $z_0\geq\tfrac12$ and
$N_{\cov,+}^{\Riem}=\lceil\Delta t^{-1}\log\tfrac{z_\infty-z_0}{z_\infty-1/2}\rceil$ otherwise.
\end{lemma}

\begin{proof}
Write $B_n:=C_n^{1/2}A(a_n)C_n^{1/2}$, so $C_{n+1}=C_n^{1/2}e^{\Delta t(\Id-B_n)}C_n^{1/2}$ and
$C_{n+1}^{-1}=C_n^{-1/2}e^{\Delta t(B_n-\Id)}C_n^{-1/2}$.

\emph{Upper bound.} From $e^X\succeq\Id+X$,
$C_{n+1}^{-1}\succeq(1-\Delta t)C_n^{-1}+\Delta t\,A(a_n)$. Assuming $C_n\preceq\Lambda\Id$, then
$C_n^{-1}\succeq\Lambda^{-1}\Id$ and $A(a_n)\succeq\alpha\Id\succeq\Lambda^{-1}\Id$; as $\Delta t\leq1$,
$C_{n+1}^{-1}\succeq\Lambda^{-1}\Id$, i.e.\ $C_{n+1}\preceq\Lambda\Id$.

\emph{Lower envelope.} Since $A(a_n)\preceq\beta\Id$ and $C_n\preceq\Lambda\Id$,
$0\preceq B_n\preceq\LB\Id$, so $0\preceq\Delta t\,B_n\preceq\Id$ (as $\Delta t\leq1/\LB$). The
scalar inequality $e^x\leq1+(e-1)x$ on $[0,1]$ gives
$e^{\Delta tB_n}\preceq\Id+(e-1)\Delta t\,B_n$, hence
$C_{n+1}^{-1}\preceq e^{-\Delta t}\bigl[C_n^{-1}+(e-1)\Delta t\,A(a_n)\bigr]$. If $C_n\succeq L_n\Id$
then $C_{n+1}\succeq L_{n+1}\Id$ with \eqref{eq:Riem-envelope}. In $y_n=(\beta L_n)^{-1}$ this is
$y_{n+1}=e^{-\Delta t}(y_n+(e-1)\Delta t)$, with solution \eqref{eq:Riem-y}. The condition
$y_n\leq2$ gives \eqref{eq:Ncov-minus}; note $L_0\leq1/(2\beta)$ forces $y_0\geq2>y_\infty$.

\emph{Upper envelope.} From $e^X\succeq\Id+X$ again,
$C_{n+1}^{-1}\succeq e^{-\Delta t}(C_n^{-1}+\Delta t\,\alpha\Id)$, so
$z_n:=\lambda_{\min}(C_n^{-1})/\alpha$ dominates $z_\infty+(z_0-z_\infty)e^{-n\Delta t}$; the
condition $C_n\preceq\tfrac2\alpha\Id$ is $z_n\geq\tfrac12$, giving $N_{\cov,+}^{\Riem}$.
\end{proof}

\begin{theorem}[Deterministic Riemannian global contraction]
\label{thm:Riem-global}
Assume \Cref{assump:logconcave-smooth} and $0<\Delta t\leq1/\LB$. Along \eqref{eq:Riem-det}, with
$L_n$ from \eqref{eq:Riem-envelope}:
\begin{enumerate}
\item For every $n\geq0$,
$\Delta_{n+1}\leq\bigl(1-\alpha L_n\Delta t(1-\tfrac{\LB\Delta t}2)\bigr)\Delta_n$.
\item For $n\geq N_{\cov,-}^{\Riem}$ one has $L_n\geq\tfrac1{2\beta}$ and
$\Delta_{n+1}\leq\bigl(1-\tfrac{\Delta t}{4\kappa}\bigr)\Delta_n$.
\item For $\delta\in(0,\Delta_0)$,
$N_{\Riem}^{\glob}(\delta)\leq N_{\cov,-}^{\Riem}+\lceil\tfrac{4\kappa}{\Delta t}\log\tfrac{\Delta_0}\delta\rceil
=O\bigl(\tfrac1{\Delta t}\logp\tfrac1{\beta\lambda_0}+\tfrac\kappa{\Delta t}\logp\tfrac{\Delta_0}\delta\bigr)$.
\end{enumerate}
\end{theorem}

\begin{proof}
\emph{(1).} By \Cref{lem:Riem-cov}, $C_n\preceq\Lambda\Id$, so \Cref{lem:det-descent} applies with
$K=\LB\geq1$ and $h=\Delta t\leq1/K$:
$\Delta_{n+1}\leq\Delta_n-\Delta t(1-\tfrac{\LB\Delta t}2)\Gradnorm(a_n)^2$. Also $C_n\succeq L_n\Id$,
so \Cref{lem:FR-W2} gives $\Gradnorm(a_n)^2\geq\alpha L_n\Delta_n$, and $1-\tfrac{\LB\Delta t}2\geq\tfrac12$.
\emph{(2).} $L_n\geq\tfrac1{2\beta}$ gives $\alpha L_n\Delta t(1-\tfrac{\LB\Delta t}2)\geq\tfrac{\alpha\Delta t}{2\beta}\cdot\tfrac12=\tfrac{\Delta t}{4\kappa}$.
\emph{(3).} $\Delta_n$ is nonincreasing, and $(1-x)^N\leq e^{-Nx}$ finishes with \eqref{eq:Ncov-minus}.
\end{proof}

\section{KL/Bregman scheme: deterministic global theory}
\label{sec:KL-global}

\begin{lemma}[KL covariance bootstrap]
\label{lem:KL-cov}
Assume \Cref{assump:logconcave-smooth} and $\lambda_{0,\min}\Id\preceq C_0\preceq\lambda_{0,\max}\Id$.
Along \eqref{eq:KL-det}, $C_n\preceq\Lambda\Id$ for all $n\geq0$. With
$L_0:=\min\{\lambda_0,\tfrac1{2\beta}\}$ and $L_{n+1}=\tfrac{(1+\Delta t)L_n}{1+\Delta t\,\beta L_n}$,
one has $C_n\succeq L_n\Id$ with
$L_n=\bigl(\beta[1+(\tfrac1{\beta L_0}-1)(1+\Delta t)^{-n}]\bigr)^{-1}$. Hence $C_n\succeq\tfrac1{2\beta}\Id$
for all
$n\geq N_{\cov,-}^{\KL}:=\lceil\tfrac{\logp(\tfrac1{\beta L_0}-1)}{\log(1+\Delta t)}\rceil
=O(\tfrac1{\Delta t}\logp\tfrac1{\beta\lambda_0})$,
and $C_n\preceq\tfrac2\alpha\Id$ for all $n\geq N_{\cov,+}^{\KL}:=\lceil\tfrac{\log2}{\log(1+\Delta t)}\rceil=O(1/\Delta t)$.
\end{lemma}

\begin{proof}
The update is $C_{n+1}^{-1}=\tfrac1{1+\Delta t}(C_n^{-1}+\Delta t\,A(a_n))$ with
$\alpha\Id\preceq A(a_n)\preceq\beta\Id$. \emph{Upper:} $C_n\preceq\Lambda\Id$ gives
$C_{n+1}^{-1}\succeq\tfrac1{1+\Delta t}(\Lambda^{-1}+\Delta t\Lambda^{-1})\Id=\Lambda^{-1}\Id$.
\emph{Lower envelope:} $C_n\succeq L_n\Id$ gives $C_{n+1}\succeq L_{n+1}\Id$; in $y_n=(\beta L_n)^{-1}$,
$y_{n+1}=1+\tfrac{y_n-1}{1+\Delta t}$, so $y_n=1+(y_0-1)(1+\Delta t)^{-n}$, and $y_n\leq2$ gives
$N_{\cov,-}^{\KL}$. \emph{Upper envelope:} $z_n:=\lambda_{\min}(C_n^{-1})/\alpha$ dominates
$1+(z_0-1)(1+\Delta t)^{-n}$; $z_n\geq\tfrac12$ once $(1+\Delta t)^{-n}\leq\tfrac12$, giving $N_{\cov,+}^{\KL}$.
\end{proof}

\begin{proposition}[KL/Bregman one-step estimate]
\label{prop:KL-onestep}
Let $0<\Delta t\leq1/\LB$ and $L_n\Id\preceq C_n\preceq\Lambda\Id$. Then, in the forward divergence,
\[
    \calE(a_{n+1})\leq\calE(a_n)-\frac1{\Delta t}\KL(\rho_{a_{n+1}}\|\rho_{a_n}),
    \qquad
    \KL(\rho_{a_{n+1}}\|\rho_{a_n})
    \geq\frac{L_n\Delta t^2}{4(1+\Delta t\,\LB)^2}\norm{\grad^{W_2}\calE(a_n)}_{W_2}^2 .
\]
\end{proposition}

The KL map \eqref{eq:KL-det} is a \emph{forward--backward} (proximal-gradient) step in the Gaussian
KL/Bregman geometry, not a full-gradient mirror step: it linearizes the potential energy $\calE_2$
explicitly and treats the entropy $\calE_1$ implicitly, which is what produces the implicit resolvent
covariance update $C_{n+1}=(1+\Delta t)(C_n^{-1}+\Delta t\,A(a_n))^{-1}$. \Cref{app:KL-onestep} proves
the proposition in full and self-contained form: the variational characterization and explicit maps
(\Cref{lem:KL-argmin}), a scalar inequality for the covariance factor (\Cref{lem:scalar-KL}), and the
one-step estimate itself (\Cref{lem:KL-onestep-app})---both the descent margin, proved directly by a
Gaussian coupling of $\rho_{a_n}$ and $\rho_{a_{n+1}}$ together with \Cref{lem:scalar-KL}, and the
lower-bound factor $\tfrac{L_n\Delta t^2}{4(1+\Delta t\LB)^2}$. Unlike the $W_2^2$-proximal
Bures--Wasserstein forward--backward scheme of \citet{diao2023forward,lambert2022variational} (whose
transport warm-start is used separately in \Cref{sec:transport-bootstrap}), the descent here is in the
\emph{forward} KL divergence $\KL(\rho_{a_{n+1}}\|\rho_{a_n})$ and the constants are specific to the
$\tfrac12$-Fisher--Rao normalization.

\begin{theorem}[Deterministic KL/Bregman global contraction]
\label{thm:KL-global}
Assume \Cref{assump:logconcave-smooth} and $0<\Delta t\leq1/\LB$. Along \eqref{eq:KL-det}:
\begin{enumerate}
\item For every $n\geq0$,
$\Delta_{n+1}\leq\bigl(1-\tfrac{\alpha L_n\Delta t}{2(1+\Delta t\,\LB)^2}\bigr)\Delta_n$.
\item For $n\geq N_{\cov,-}^{\KL}$, $\Delta_{n+1}\leq\bigl(1-\tfrac{\Delta t}{16\kappa}\bigr)\Delta_n$.
\item $N_{\KL}^{\glob}(\delta)\leq N_{\cov,-}^{\KL}+\lceil\tfrac{16\kappa}{\Delta t}\log\tfrac{\Delta_0}\delta\rceil
=O\bigl(\tfrac1{\Delta t}\logp\tfrac1{\beta\lambda_0}+\tfrac\kappa{\Delta t}\logp\tfrac{\Delta_0}\delta\bigr)$.
\end{enumerate}
\end{theorem}

\begin{proof}
\emph{(1).} By \Cref{prop:KL-onestep} and \Cref{lem:W2-PL},
$\calE(a_{n+1})-\calE(a_n)\leq-\tfrac{L_n\Delta t}{4(1+\Delta t\LB)^2}\cdot2\alpha\Delta_n$.
\emph{(2).} $L_n\geq\tfrac1{2\beta}$ and $\Delta t\LB\leq1$ give $(1+\Delta t\LB)^2\leq4$, so the
per-step factor is at most $1-\tfrac{\alpha\Delta t}{2\cdot4\cdot2\beta}=1-\tfrac{\Delta t}{16\kappa}$.
\emph{(3).} As in \Cref{thm:Riem-global}(3).
\end{proof}

\section{Local convergence: the Gaussian core and a Gaussian-sharp region}
\label{sec:local}

The organizing principle of the local theory is a \emph{Gaussian-core decomposition}. The natural
gradient field of the standard Gaussian target is nonlinear in the covariance but can be dissipated
exactly; treating it as a Taylor remainder discards a large, genuinely dissipative contribution. We
write
\[
    F=\FGauss+H,
    \qquad
    \FGauss(m,C):=(-Cm,\;C-C^2),
\]
retain $\FGauss$ exactly, and Taylor-bound only the non-Gaussian perturbation $H$. Because
$H\equiv0$ whenever the target is Gaussian, the certified energy-sublevel region is sharp for
Gaussian targets.

\subsection{Whitened coordinates and the linearized gap}
\label{subsec:whitened}

Define $\widehat m:=C_\star^{-1/2}(m-m_\star)$, $\widehat C:=C_\star^{-1/2}CC_\star^{-1/2}$, and
$\widehat V(z):=V(m_\star+C_\star^{1/2}z)$. Let $0<\alpha_\star\leq\beta_\star<\infty$ satisfy
$\alpha_\star\Id\preceq\nabla^2\widehat V(z)\preceq\beta_\star\Id$, with
$\kappa_\star:=\beta_\star/\alpha_\star$. By \eqref{eq:first-order},
$\E_{\N(0,\Id)}[\nabla\widehat V]=0$ and $\E_{\N(0,\Id)}[\nabla^2\widehat V]=\Id$, so
$\alpha_\star\leq1\leq\beta_\star$, and one may take $\alpha_\star\geq1/\kappa$,
$\beta_\star\leq\kappa$, $\kappa_\star\leq\kappa^2$. Throughout this section we work in whitened
coordinates: $a=(m,C)$ denotes the whitened parameter, $a_\star=(0,\Id)$. Statements transfer back
through the fixed equilibrium norms
\begin{equation}
\label{eq:local-norms}
    \norm{a-a_\star}_\star^2:=\norm{\widehat m}^2+\tfrac12\norm{\widehat C-\Id}_{\Frob}^2,
    \qquad
    \norm{a-a_\star}_{\star,\Delta t}^2:=\norm{\widehat m}^2+\tfrac{1+\Delta t}2\norm{\widehat C-\Id}_{\Frob}^2 .
\end{equation}
We write $\xi:=a-a_\star$, $D(a):=\norm{\xi}_\star^2$, and $F(a):=(Cg(a),\,C-CA(a)C)$, so
\eqref{eq:NG-flow} reads $\dot a=F(a)$, $F(a_\star)=0$. Set
\begin{equation}
\label{eq:Gamma}
    t_0:=\max\{1,\log(4\beta_\star)\},
    \qquad
    \Gamma:=\min\Bigl\{\beta_\star-1,\;\sqrt{N_\theta}\bigl(4\sqrt{2t_0}+4t_0+4\bigr)\Bigr\},
    \qquad
    \gund:=\frac1{4+\Gamma} ,
\end{equation}
so that $\Gamma=O\bigl(1+\sqrt{N_\theta}\log\kappa_\star\bigr)=O\bigl(1+\sqrt{N_\theta}\log\kappa\bigr)$.

\begin{definition}[Linearized score operators]
\label{def:linear-score-ops}
Let $Z\sim\N(0,\Id)$ and $B(Z):=\nabla^2\widehat V(Z)$, so $\E B(Z)=\Id$. For $u\in\R^{N_\theta}$,
$X\in\Sym(N_\theta)$ define
$\mathcal T[X]_k:=\E[\Tr(XB(Z))Z_k]$, $\mathcal T^*[u]_{ij}:=\E[(u^\top Z)B_{ij}(Z)]$, and
$\mathcal H[X]_{ij}:=\E[B_{ij}(Z)(Z^\top XZ-\Tr X)]$. Gaussian score identities give
$u^\top\mathcal T[X]=\Tr(\mathcal T^*[u]X)$ and $\Tr(\mathcal H[X]Y)=\Tr(X\mathcal H[Y])$.
\end{definition}

\begin{lemma}[Diagonal-mode bounds]
\label{lem:linear-diagonal-bounds}
Fix an orthogonal basis in which $X=\diag(\lambda)$, and set $\widetilde T_{ki}:=\E[B_{ii}(Z)Z_k]$,
$G_{ij}:=\E[B_{ii}(Z)Z_j^2]$, $D:=G-\mathbf1\mathbf1^\top$. Then $D$ is symmetric,
\begin{equation}
\label{eq:T-Schur}
    \widetilde T^\top\widetilde T\preceq\Id+D,
    \qquad\text{and}\qquad
    \lambda_{\max}(D)\leq\Gamma .
\end{equation}
\end{lemma}

\begin{proof}
First note that the symmetry of $D$ and the entrywise identity used below follow directly from
Gaussian integration by parts. Indeed, since $\E B=\Id$, for $i\neq j$ two integrations by parts give
\[
    D_{ij}=\E\bigl[B_{ii}(Z)(Z_j^2-1)\bigr]
    =\E\bigl[\widehat V(Z)(Z_i^2-1)(Z_j^2-1)\bigr]=D_{ji},
\]
and integrating one $i$- and one $j$-derivative gives
$\E[B_{ij}(Z)Z_iZ_j]=D_{ij}$. For $i=j$,
$\E[B_{ii}(Z)Z_i^2]=1+D_{ii}$ is immediate from $\E B_{ii}=1$. These integrations are justified by
the bounded Hessian and Gaussian tails (equivalently, by mollification). Thus $D$ is symmetric and
$\E[B_{ij}(Z)Z_iZ_j]=\delta_{ij}+D_{ij}=\delta_{ij}+G_{ij}-1$ for all $i,j$.

For $\eta(Z):=u+\diag(\lambda)Z$, expanding with this identity gives
\begin{equation}
\label{eq:eta-B-identity}
    \E[\eta(Z)^\top B(Z)\eta(Z)]=\|u\|^2+2u^\top\widetilde T\lambda+\|\lambda\|^2+\lambda^\top D\lambda,
\end{equation}
Because $B(Z)\succeq0$, the block quadratic form is nonnegative; its Schur complement with respect to
the identity mean block gives
\eqref{eq:T-Schur}(left). Taking $u=0$ and $B(Z)\preceq\beta_\star\Id$ gives
$\|\lambda\|^2+\lambda^\top D\lambda\leq\beta_\star\|\lambda\|^2$, so
$\lambda_{\max}(D)\leq\beta_\star-1$. For the dimension-dependent bound, normalize $\|\lambda\|=1$,
set $X:=\diag(\lambda)$ in the fixed basis, and
$Q_X(Z):=Z^\top XZ-\Tr X=\sum_j\lambda_j(Z_j^2-1)$. Since
$D_{ij}=\E[B_{ii}(Z)(Z_j^2-1)]$ (using $\E B=\Id$), the quadratic form has the exact matrix
representation
\[
    \lambda^\top D\lambda=\E\bigl[\Tr\bigl(XB(Z)\bigr)\,Q_X(Z)\bigr].
\]
Because $B(Z)\succeq0$, splitting $X=X_+-X_-$ gives $|\Tr(XB)|\leq\Tr(|X|B)$ pointwise, with
$\E\Tr(|X|B)=\Tr|X|=\|\lambda\|_1\leq\sqrt{N_\theta}$ and, almost surely,
$\Tr(|X|B)\leq\beta_\star\Tr|X|\leq\beta_\star\sqrt{N_\theta}$. The Laurent--Massart bound
\citep{laurent2000adaptive} for $\sum_ja_j(Z_j^2-1)$ with $a\succeq0$, applied to the positive and
negative parts of $\lambda$ (so $\|\lambda_+\|+\|\lambda_-\|\leq\sqrt2$, $\|\lambda\|_\infty\leq1$),
gives the two-sided tail
$\Prob\bigl(|Q_X|\geq2\sqrt{2t}+2t\bigr)\leq4e^{-t}$ for all $t\geq0$. With $t_0$ from
\eqref{eq:Gamma} and $T:=2\sqrt{2t_0}+2t_0$, split the expectation. On $\{|Q_X|\leq T\}$,
$\E[\Tr(|X|B)\,|Q_X|]\leq T\,\E\Tr(|X|B)\leq T\sqrt{N_\theta}$. On the complement, the pointwise
bound and the tail integral
\[
    \E\bigl[|Q_X|;|Q_X|>T\bigr]
    \leq T\cdot4e^{-t_0}+\int_{t_0}^\infty4e^{-t}\Bigl(\frac{\sqrt2}{\sqrt t}+2\Bigr)\dd t
    \leq4e^{-t_0}\bigl(T+\sqrt2+2\bigr)\leq4e^{-t_0}(T+4)
\]
(using $t_0\geq1$) give, since $4\beta_\star e^{-t_0}\leq1$, a contribution at most
$\beta_\star\sqrt{N_\theta}\cdot4e^{-t_0}(T+4)\leq\sqrt{N_\theta}(T+4)$. Hence
$|\lambda^\top D\lambda|\leq\sqrt{N_\theta}(2T+4)=\sqrt{N_\theta}(4\sqrt{2t_0}+4t_0+4)$ uniformly over
$\|\lambda\|=1$ (and over the choice of diagonalizing basis, since the representation is basis-free
in $X$). Combining gives $\lambda_{\max}(D)\leq\Gamma$.
\end{proof}

\begin{proposition}[Linearized gap]
\label{prop:linearized}
Let $J_\star:=DF(a_\star)$ and $\Lstar:=-J_\star$. Then
\[
    \Lstar(u,X)=\Bigl(u+\tfrac12\mathcal T[X],\;X+\mathcal T^*[u]+\tfrac12\mathcal H[X]\Bigr),
\]
$\Lstar$ is self-adjoint in $\inner{\cdot}{\cdot}_\star$, and with
$\gstar:=\lambda_{\min}(\Lstar)$, $\Lambda_\star:=\lambda_{\max}(\Lstar)$,
\begin{equation}
\label{eq:linear-gap-bounds}
    \boxed{\;\max\Bigl\{\alpha_\star,\frac1{4+\Gamma}\Bigr\}\leq\gstar\leq\Lambda_\star\leq\min\Bigl\{\beta_\star,\frac{4+\Gamma}2\Bigr\} .\;}
\end{equation}
For a Gaussian target, $\mathcal T=0$, $\mathcal H=0$, and $\gstar=\Lambda_\star=1$.
\end{proposition}

\begin{proof}
Bonnet--Price differentiation at $a_\star=(0,\Id)$ gives
$\E_{\N(u,\Id+X)}[\nabla\widehat V]=u+\tfrac12\mathcal T[X]+O(\|(u,X)\|_\star^2)$ and
$\E_{\N(u,\Id+X)}[\nabla^2\widehat V]=\Id+\mathcal T^*[u]+\tfrac12\mathcal H[X]+O(\|(u,X)\|_\star^2)$;
substituting into $F(m,C)=(-C\E[\nabla\widehat V],\,C-C\E[\nabla^2\widehat V]C)$ gives the stated
$\Lstar$. Self-adjointness follows from the adjoint identities of \Cref{def:linear-score-ops}. Its
quadratic form, after diagonalizing $X=\diag(\lambda)$, is
\[
    Q(u,\diag\lambda)=\|u\|^2+u^\top\widetilde T\lambda+\tfrac12\|\lambda\|^2+\tfrac14\lambda^\top D\lambda.
\]
Let $\nu:=1/(4+\Gamma)$. Completing the square in $u$, coercivity $Q\geq\nu\|\cdot\|_\star^2$ reduces,
via \eqref{eq:T-Schur}, to $\tfrac{1-\nu}2-\tfrac{1+\nu\Gamma}{4(1-\nu)}\geq0$, i.e.\
$1-\nu(4+\Gamma)+2\nu^2\geq0$, which holds at $\nu=1/(4+\Gamma)$ (the left side equals $2\nu^2>0$).
For the upper bound, \Cref{lem:linear-diagonal-bounds} gives
$Q(u,\diag\lambda)\leq\|u\|^2+\sqrt{1+\Gamma}\|u\|\|\lambda\|+(\tfrac12+\tfrac\Gamma4)\|\lambda\|^2\leq\tfrac{4+\Gamma}2(\|u\|^2+\tfrac12\|\lambda\|^2)$.
Finally, the direct spectral sandwich: taking $u\to2u$ in the exact identity
\eqref{eq:eta-B-identity} and adding $\tfrac14\|\lambda\|^2$,
\[
    Q(u,\diag\lambda)
    =\frac14\,\E\bigl[(2u+XZ)^\top B(Z)\,(2u+XZ)\bigr]+\frac14\|\lambda\|^2,
    \qquad X=\diag(\lambda),
\]
and sandwiching $\alpha_\star\Id\preceq B(Z)\preceq\beta_\star\Id$ with
$\E\|2u+XZ\|^2=4\|u\|^2+\|\lambda\|^2$ gives
$\alpha_\star\|u\|^2+\tfrac{1+\alpha_\star}4\|\lambda\|^2\leq Q\leq\beta_\star\|u\|^2+\tfrac{1+\beta_\star}4\|\lambda\|^2$;
since $\alpha_\star\leq1\leq\beta_\star$, $\tfrac{1+\alpha_\star}4\geq\tfrac{\alpha_\star}2$ and
$\tfrac{1+\beta_\star}4\leq\tfrac{\beta_\star}2$, whence
$\alpha_\star\|\xi\|_\star^2\leq Q(\xi)\leq\beta_\star\|\xi\|_\star^2$. The spectral theorem gives
\eqref{eq:linear-gap-bounds}. The two sides of the maximum and minimum are complementary: the
$\Gamma$-branch is universal, while the $(\alpha_\star,\beta_\star)$-branch is sharper for
well-conditioned targets even when the universal $\Gamma$ estimate is pessimistic.
\end{proof}

\begin{proposition}[Linearized contractions of the deterministic maps]
\label{prop:linearized-maps}
Let $T_{\Riem,h}$, $T_{\KL,h}$ be the maps \eqref{eq:Riem-det}, \eqref{eq:KL-det} in whitened
coordinates. If $0<h\leq2/(4+\Gamma)$, then
\[
    \|DT_{\Riem,h}(a_\star)\xi\|_\star\leq\Bigl(1-\tfrac h{4+\Gamma}\Bigr)\|\xi\|_\star,
    \qquad
    \|DT_{\KL,h}(a_\star)\xi\|_{\star,h}\leq\Bigl(1-\tfrac h{4+2h+\Gamma}\Bigr)\|\xi\|_{\star,h}.
\]
\end{proposition}

\begin{proof}
For the Riemannian map, $DT_{\Riem,h}(a_\star)=\Id-h\Lstar$, whose spectrum lies in
$[0,1-h/(4+\Gamma)]$ since $h\Lambda_\star\leq1$ by \eqref{eq:linear-gap-bounds}. For the KL map,
$M_{\KL,h}:=DT_{\KL,h}(a_\star)$ is self-adjoint in $\inner{\cdot}{\cdot}_{\star,h}$ and satisfies
$\|\xi\|_{\star,h}^2-\inner{M_{\KL,h}\xi}{\xi}_{\star,h}=hQ(\xi)$ with $Q$ as in
\Cref{prop:linearized}. The weighted coercivity $Q(\xi)\geq\tfrac1{4+2h+\Gamma}\|\xi\|_{\star,h}^2$
follows by the same completing-the-square argument, reducing to
$1-\nu(4+2h+\Gamma)+2(1+h)\nu^2\geq0$ at $\nu=1/(4+2h+\Gamma)$. Since $Q(\xi)\leq\tfrac{4+\Gamma}2\|\xi\|_{\star,h}^2$
and $h\leq2/(4+\Gamma)$, the spectrum of $M_{\KL,h}$ is nonnegative and bounded above by
$1-h/(4+2h+\Gamma)$.
\end{proof}

\subsection{Energy geometry: entropy-enhanced coercivity and bands}

\begin{lemma}[Entropy-enhanced coercivity]
\label{lem:entropy-coerc}
Write $B:=C^{1/2}$ with eigenvalues $s_1,\dots,s_{N_\theta}$ and
$\phi_{\alpha_\star}(s):=s-1-\log s+\tfrac{\alpha_\star}2(s-1)^2$. Then
\begin{equation}
\label{eq:entropy-coerc}
    \Delta(m,C)\geq\frac{\alpha_\star}2\norm{m}^2+\sum_{i=1}^{N_\theta}\phi_{\alpha_\star}(s_i).
\end{equation}
For a Gaussian target ($\alpha_\star=1$) this is an equality and
$\Delta_{\mathrm G}(m,C)=\tfrac12\norm{m}^2+\sum_i\varphi(c_i)$ with $\varphi(c):=\tfrac12(c-1-\log c)$
and $c_i$ the eigenvalues of $C$.
\end{lemma}

\begin{proof}
With $X=m+BZ$, $\Delta=\E[\widehat V(m+BZ)-\widehat V(Z)]-\log\det B$. By $\alpha_\star$-strong
convexity, $\widehat V(m+BZ)-\widehat V(Z)\geq\nabla\widehat V(Z)^\top(m+(B-\Id)Z)+\tfrac{\alpha_\star}2\norm{m+(B-\Id)Z}^2$.
Taking expectations: $\E[\nabla\widehat V(Z)]^\top m=0$; $\E[\nabla\widehat V(Z)^\top(B-\Id)Z]=\Tr(B-\Id)$
by Gaussian integration by parts; and $\E\norm{m+(B-\Id)Z}^2=\norm m^2+\norm{B-\Id}_{\Frob}^2$.
Hence $\Delta\geq\tfrac{\alpha_\star}2\norm m^2+\Tr(B-\Id)-\log\det B+\tfrac{\alpha_\star}2\norm{B-\Id}_{\Frob}^2$,
which diagonalizes to \eqref{eq:entropy-coerc}. For a Gaussian target the convexity step is an
equality, and $\phi_1(s)=\tfrac12(s^2-1-\log s^2)=\varphi(c)$ with $c=s^2$.
\end{proof}

\begin{corollary}[Covariance band and norm conversion]
\label{cor:band}
$\phi_{\alpha_\star}$ is strictly decreasing on $(0,1)$, increasing on $(1,\infty)$, with minimum
$0$ at $1$. For $\delta>0$ let $0<s_-(\delta)<1<s_+(\delta)$ solve $\phi_{\alpha_\star}(s)=\delta$,
and set $\ell_\delta:=s_-(\delta)^2$, $U_\delta:=s_+(\delta)^2$. Then $\Delta(a)\leq\delta$ implies
$\ell_\delta\Id\preceq C\preceq U_\delta\Id$. With
\[
    \vartheta(\delta):=\sup_{s\in[s_-,s_+],\,s\neq1}\frac{(s^2-1)^2}{2\phi_{\alpha_\star}(s)}
    \leq\frac{(1+s_+)^2}{\alpha_\star+s_+^{-2}},
    \qquad
    \mathfrak b(\delta)^2:=\max\Bigl\{\tfrac2{\alpha_\star},\vartheta(\delta)\Bigr\},
\]
one has $\Delta(a)\leq\delta\Rightarrow\norm{a-a_\star}_\star\leq\mathfrak b(\delta)\sqrt\delta$, and
$\norm{a-a_\star}_{\star,\Delta t}\leq\mathfrak b_{\Delta t}(\delta)\sqrt\delta$ with
$\mathfrak b_{\Delta t}(\delta)^2:=\max\{2/\alpha_\star,(1+\Delta t)\vartheta(\delta)\}$. For a
Gaussian target, $\ell_\delta,U_\delta$ are the roots of $\varphi(c)=\delta$.
\end{corollary}

\begin{proof}
Each $\phi_{\alpha_\star}(s_i)\leq\delta$ by \eqref{eq:entropy-coerc}, so $s_i\in[s_-,s_+]$; squaring
gives the band. Writing $E_m:=\tfrac{\alpha_\star}2\norm m^2$, $E_C:=\sum_i\phi_{\alpha_\star}(s_i)$,
$E_m+E_C\leq\delta$, one has $\norm m^2=\tfrac2{\alpha_\star}E_m$ and
$\tfrac12\norm{C-\Id}_{\Frob}^2\leq\vartheta(\delta)E_C$, so
$\norm{a-a_\star}_\star^2\leq\mathfrak b(\delta)^2\delta$; the bound on $\vartheta$ uses
$\phi_{\alpha_\star}(s)\geq\tfrac{(s-1)^2}2(s_+^{-2}+\alpha_\star)$.
\end{proof}

\begin{definition}[Regions, hull, moduli]
\label{def:region}
For $\delta>0$, $\Ureg_\delta:=\{a:\Delta(a)\leq\delta\}$ and
$\Hreg_\delta:=\{a_\star+s(a-a_\star):a\in\Ureg_\delta,s\in[0,1]\}$. Set
$\bar\Lambda_\delta:=\max\{U_\delta,1/\alpha_\star\}$ and the non-Gaussian moduli
\[
    \mathcal B_\delta:=\sup_{a\in\Hreg_\delta}\norm{A(a)-\Id}_{\op},
    \quad
    \Sigma_\delta:=\sup_{a\in\Hreg_\delta}\bigl(\E_{\rho_a}\norm{\nabla^2\widehat V(X)-A(a)}_{\Frob}^2\bigr)^{1/2},
    \quad
    \bar\Sigma_\delta:=\min\{\Sigma_\delta,\beta_\star-\alpha_\star\}.
\]
All three vanish for Gaussian targets and are nondecreasing in $\delta$.
\end{definition}

\begin{lemma}[Geometry and invariance]
\label{lem:region-geom}
Let $\delta>0$, $a=(m,C)\in\Ureg_\delta$. Then (i) $\ell_\delta\Id\preceq C\preceq U_\delta\Id$ and
$\norm{a-a_\star}_\star\leq\mathfrak b(\delta)\sqrt\delta$; (ii) every $a'\in\Hreg_\delta$ satisfies
the same bounds; (iii) $\Ureg_\delta$ is compact; (iv) $\Ureg_\delta$ is forward invariant for the
flow, and for any scheme step with $\calE(a_{n+1})\leq\calE(a_n)$.
\end{lemma}

\begin{proof}
(i) is \Cref{cor:band}. (ii): $C'=(1-s)\Id+sC$ has eigenvalues in $[\ell_\delta,U_\delta]$ (convex
combinations of $1$ and $\lambda_i(C)$), and $\norm{a'-a_\star}_\star=s\norm\xi_\star$. (iii): closed
and bounded by (i) and $\norm m\leq\sqrt{2\delta/\alpha_\star}$. (iv): $\tfrac{\dd}{\dd t}\Delta=-\Gradnorm^2\leq0$;
a descent step keeps $\Delta(a_{n+1})\leq\delta$.
\end{proof}

\subsection{The exact Gaussian core}

\begin{lemma}[Gaussian-core dissipation]
\label{lem:gauss-core}
Let $\FGauss(m,C):=(-Cm,\,C-C^2)$, the natural gradient field of the standard Gaussian target
$\widehat V_{\mathrm G}(z)=\tfrac12\norm z^2$. Then:
\begin{enumerate}
\item $D\FGauss(a_\star)=-\Id$ on $\R^{N_\theta}\times\Sym(N_\theta)$.
\item For every $a=(m,C)$, with $\xi=a-a_\star$,
$\inner{\xi}{\FGauss(a)}_\star=-m^\top Cm-\tfrac12\Tr\bigl(C(C-\Id)^2\bigr)=-\tfrac12q_{\mathrm G}(a)D(a)$,
where $q_{\mathrm G}(a):=\dfrac{2m^\top Cm+\Tr(C(C-\Id)^2)}{\norm m^2+\tfrac12\norm{C-\Id}_{\Frob}^2}$.
\item $q_{\mathrm G}(a)\geq2\lambda_{\min}(C)$, with equality when $m=0$ and $C-\Id$ is supported on
a minimal eigendirection of $C$; hence on $\Ureg_\delta$,
$\inner{\xi}{\FGauss(a)}_\star\leq-\ell_\delta D(a)$.
\end{enumerate}
\end{lemma}

\begin{proof}
(1) At $a_\star=(0,\Id)$ with $\zeta=(u,U)$: $D(-Cm)[\zeta]=-(U\cdot0+\Id\cdot u)=-u$ and
$D(C-C^2)[\zeta]=U-(U\Id+\Id U)=-U$. (2) Write $X=C-\Id$; then $C-C^2=-C(C-\Id)=-(X+X^2)$, so
$\inner{\xi}{\FGauss}_\star=-m^\top Cm+\tfrac12\Tr(X\cdot(-(X+X^2)))=-m^\top Cm-\tfrac12\Tr(CX^2)$,
using $\Id+X=C$; and $\Tr(CX^2)=\Tr(C(C-\Id)^2)$. (3) $2m^\top Cm\geq2\lambda_{\min}(C)\norm m^2$ and,
since $C$ and $X$ commute, $\Tr(CX^2)\geq\lambda_{\min}(C)\norm X_{\Frob}^2$; the numerator is thus
$\geq2\lambda_{\min}(C)D(a)$, with equality as stated. On $\Ureg_\delta$,
$\lambda_{\min}(C)\geq\ell_\delta$.
\end{proof}

\subsection{The non-Gaussian perturbation}

Write $H:=F-\FGauss$. Since $g_{\mathrm G}(a)=-m$ and $A_{\mathrm G}(a)=\Id$,
\begin{equation}
\label{eq:H-def}
    H(a)=\bigl(C(g(a)+m),\;-C(A(a)-\Id)C\bigr),\qquad H(a_\star)=0 .
\end{equation}
Both components vanish identically when $\widehat V=\widehat V_{\mathrm G}$.

\begin{lemma}[Score identities and the sharp second-score constant]
\label{lem:second-score}
Fix $a=(m,C)$, $\zeta=(u,U)$, $v:=C^{-1/2}u$, $S:=C^{-1/2}UC^{-1/2}$,
$t:=\norm\zeta_a=(\norm v^2+\tfrac12\norm S_{\Frob}^2)^{1/2}$, $X=m+C^{1/2}Z$. With
$q_\zeta:=v^\top Z+\tfrac12(Z^\top SZ-\Tr S)$ and
$\chi_\zeta:=q_\zeta^2+(\tfrac12\Tr S^2-\norm v^2-2v^\top SZ-Z^\top S^2Z)$, for $\phi$ of polynomial
growth and $\Phi(a):=\E_{\rho_a}[\phi(X)]$,
\[
    D\Phi(a)[\zeta]=\E[\phi\,q_\zeta],
    \qquad
    D^2\Phi(a)[\zeta,\zeta]=\E[\phi\,\chi_\zeta].
\]
Moreover $q_\zeta$ has Hermite chaoses $1,2$ and $\chi_\zeta$ has chaoses $2,3,4$; in particular
$\E[q_\zeta]=\E[\chi_\zeta]=0$, $\E[Z\chi_\zeta]=0$, $\E[q_\zeta^2]=t^2$, and
$\norm{\chi_\zeta}_{L^2}\leq\sqrt6\,t^2$ (sharp, attained by rank-one covariance directions).
\end{lemma}

\begin{proof}
Deferred to \Cref{app:Kng}; the differentiation identities follow from differentiating
$\log\rho_{a_s}$ under the integral, and the sharp constant from the Hermite chaos decomposition of
$\chi_\zeta$ (with $c=v^\top S^2v\leq ab$ and $d_4=\Tr S^4\leq b^2$).
\end{proof}

\begin{lemma}[Non-Gaussian second-differential modulus]
\label{lem:Kng}
Define
\[
    \Kng(\delta):=\sup_{a\in\Hreg_\delta}\;\sup_{\zeta\neq0}\;\frac{\norm{D^2H(a)[\zeta,\zeta]}_\star}{\norm\zeta_\star^2}.
\]
Then
\[
    \Kng(\delta)\leq\bigl(K_{\mathrm{ng},m}(\delta)^2+\tfrac12K_{\mathrm{ng},C}(\delta)^2\bigr)^{1/2},
\]
with
\[
    K_{\mathrm{ng},m}(\delta):=\frac{U_\delta^{3/2}}{\ell_\delta^2}\bigl(\sqrt3\bar\Sigma_\delta+\Sigma_\delta+\sqrt{\Sigma_\delta^2+2\mathcal B_\delta^2}\bigr),
    \quad
    K_{\mathrm{ng},C}(\delta):=\frac{U_\delta^2}{\ell_\delta^2}\bigl((4\sqrt2+\sqrt6)\Sigma_\delta+4\mathcal B_\delta\bigr).
\]
In particular $\Kng(\delta)=0$ for Gaussian targets, and, with the spectral-band factors
$U_\delta,\ell_\delta$ of \Cref{cor:band} held fixed,
$\Kng(\delta)=O\!\bigl(\tfrac{U_\delta^2}{\ell_\delta^2}(\mathcal B_\delta+\Sigma_\delta)\bigr)$; the
band factors must be retained in the $O(\cdot)$ since they degrade as $\delta$ grows.
\end{lemma}

\begin{proof}
Deferred to \Cref{app:Kng}. In outline: converting $t\leq\ell_\delta^{-1}\norm\zeta_\star$, the mean
block $D^2H_m=2U\,D(g+m)[\zeta]+C\,D^2g[\zeta,\zeta]$ is controlled by the Ornstein--Uhlenbeck inverse
of $q_\zeta$ (first term) and of $\chi_\zeta$ (second term, giving the operator/Frobenius minimum
$\bar\Sigma_\delta$), and the covariance block $H_C=-CWC$ by the product rule together with the score
bounds $\norm{DW[\zeta]}_{\Frob}\leq\Sigma_\delta t$ and $\norm{D^2W[\zeta,\zeta]}_{\Frob}\leq\sqrt6\Sigma_\delta t^2$
of \Cref{lem:second-score}.
\end{proof}

\begin{corollary}[Non-Gaussian Taylor remainder]
\label{cor:taylor-ng}
For $a=a_\star+\xi\in\Ureg_\delta$, write $F(a)=J_\star\xi+N_{\mathrm G}(\xi)+R_{\mathrm{ng}}(\xi)$
with $N_{\mathrm G}(\xi):=\FGauss(a)-J_{\mathrm G}\xi$ ($J_{\mathrm G}=-\Id$) and
$R_{\mathrm{ng}}(\xi):=H(a)-DH(a_\star)\xi$. Then
$\norm{R_{\mathrm{ng}}(\xi)}_\star\leq\tfrac12\Kng(\delta)\norm\xi_\star^2$.
\end{corollary}

\begin{proof}
$R_{\mathrm{ng}}(\xi)=\int_0^1(1-s)D^2H(a_\star+s\xi)[\xi,\xi]\dd s$ (as $H(a_\star)=0$); each
segment point lies in $\Hreg_\delta$, so \Cref{lem:Kng} applies pointwise.
\end{proof}

\subsection{Main continuous-time local theorem and sharpness}

\begin{theorem}[Gaussian-core local convergence of the flow]
\label{thm:core-flow}
Let $0<\gamma_0\leq\gstar$ and $\delta>0$, and set
\begin{equation}
\label{eq:rho-core}
    \rho_{\mathrm{core}}(\delta):=2\gamma_0-2+2\ell_\delta-\Kng(\delta)\,\mathfrak b(\delta)\sqrt\delta .
\end{equation}
If $\rho_{\mathrm{core}}(\delta)>0$ and $\Delta(a_0)\leq\delta$, then the flow \eqref{eq:NG-flow}
remains in $\Ureg_\delta$ and $D(a_t)\leq e^{-\rho_{\mathrm{core}}(\delta)t}D(a_0)$ for all $t\geq0$.
\end{theorem}

\begin{proof}
$\Ureg_\delta$ is forward invariant and compact (\Cref{lem:region-geom}). Since $\norm\cdot_\star$ is
a fixed norm, $\tfrac12\tfrac{\dd}{\dd t}D=\inner{\xi_t}{\FGauss(a_t)}_\star+\inner{\xi_t}{H(a_t)}_\star$.
By \Cref{lem:gauss-core}(3), $\inner{\xi_t}{\FGauss(a_t)}_\star\leq-\ell_\delta D(a_t)$. Next
$H(a_t)=B_\star\xi_t+R_{\mathrm{ng}}(\xi_t)$ with $B_\star=DH(a_\star)=J_\star+\Id$ self-adjoint, so
$\inner{\xi_t}{B_\star\xi_t}_\star=\inner{\xi_t}{J_\star\xi_t}_\star+\norm{\xi_t}_\star^2\leq(1-\gamma_0)D(a_t)$,
using $\inner{\xi}{J_\star\xi}_\star\leq-\gstar\norm\xi_\star^2\leq-\gamma_0\norm\xi_\star^2$. Finally,
by \Cref{cor:taylor-ng} and $\norm{\xi_t}_\star\leq\mathfrak b(\delta)\sqrt\delta$,
$\inner{\xi_t}{R_{\mathrm{ng}}(\xi_t)}_\star\leq\tfrac12\Kng(\delta)\mathfrak b(\delta)\sqrt\delta\,D(a_t)$.
Adding, $\tfrac{\dd}{\dd t}D\leq-\rho_{\mathrm{core}}(\delta)D$; Gr\"onwall gives the claim.
\end{proof}

\begin{corollary}[Universal certified level]
\label{cor:core-universal}
With $\gamma_0=\gund$, $\rho_{\mathrm{core}}(\delta)\to2\gund$ as $\delta\downarrow0$, so for every
$\rho<2\gund$ there is $\delta_{\mathrm c}(\rho)>0$ with $\rho_{\mathrm{core}}(\delta)\geq\rho$ on
$(0,\delta_{\mathrm c}(\rho)]$.
\end{corollary}

\begin{corollary}[Exact sharpness for Gaussian targets]
\label{cor:gauss-sharp}
Let the target be Gaussian. Then $\gstar=1$, $\Kng\equiv0$, and \Cref{thm:core-flow} holds with the
exact rate $\rho_{\mathrm{core}}(\delta)=2\ell_\delta$. Consequently, for every $\rho\in(0,2)$, the
largest energy sublevel on which rate $\rho$ is certified is $\{\Delta\leq\sharpthr_{\mathrm G}(\rho)\}$
with
\begin{equation}
\label{eq:gauss-sharp-thresh}
    \sharpthr_{\mathrm G}(\rho)=\varphi\Bigl(\tfrac\rho2\Bigr)=\tfrac12\Bigl(\tfrac\rho2-1-\log\tfrac\rho2\Bigr),
\end{equation}
independent of dimension, and equals the exact supremal threshold $\sharpthr(\rho)$
(\Cref{prop:sharp}; the defining infimum is approached, not attained, along limiting
minimal-eigendirection covariance modes).
\end{corollary}

\begin{proof}
For a Gaussian target $F=\FGauss$, $H\equiv0$, so $\gstar=1$ and $\Kng\equiv0$, whence
$\rho_{\mathrm{core}}=2\ell_\delta$. Rate $\rho$ is certified iff $\ell_\delta\geq\rho/2$; since
$\ell_\delta$ is the lower root of $\varphi(c)=\delta$ and $\varphi$ decreases on $(0,1)$, this is
$\delta\leq\varphi(\rho/2)$.
\end{proof}

\begin{definition}[Instantaneous rate and maximal threshold]
\label{def:sharp}
For $a\neq a_\star$, set
\[
    q_{\mathrm c}(a):=\frac{-2\inner{a-a_\star}{F(a)}_\star}{\norm{a-a_\star}_\star^2},
    \qquad
    \sharpthr(\rho):=\inf\{\Delta(a):a\neq a_\star,\;q_{\mathrm c}(a)<\rho\}.
\]
\end{definition}

\begin{proposition}[The benchmark is exact and dominates the certified level]
\label{prop:sharp}
Let $\rho>0$. (1) If $\Delta(a_0)<\sharpthr(\rho)$ then $D(a_t)\leq e^{-\rho t}D(a_0)$ for all
$t\geq0$. (2) For every $\delta>\sharpthr(\rho)$ there is $a_0$ with $\Delta(a_0)\leq\delta$ and
instantaneous rate $<\rho$; so $\sharpthr(\rho)$ is the largest energy threshold guaranteeing rate
$\rho$. (3) For a Gaussian target, $\sharpthr(\rho)=\sharpthr_{\mathrm G}(\rho)=\varphi(\rho/2)$ for
$\rho\in(0,2)$, in every dimension, and the certified level of \Cref{cor:gauss-sharp} attains it.
\end{proposition}

\begin{proof}
(1) By energy descent $\Delta(a_t)\leq\Delta(a_0)<\sharpthr(\rho)$, so $q_{\mathrm c}(a_t)\geq\rho$;
thus $\tfrac{\dd}{\dd t}D=-q_{\mathrm c}(a_t)D\leq-\rho D$. (2) is the definition of the infimum.
(3) For a Gaussian target $q_{\mathrm c}=q_{\mathrm G}\geq2\lambda_{\min}(C)$ with equality when
$m=0$ and $C-\Id$ lies in the minimal eigendirection. If $q_{\mathrm c}(a)<\rho$ then some
eigenvalue $c_j<\rho/2$, so $\Delta_{\mathrm G}(a)\geq\varphi(c_j)>\varphi(\rho/2)$; conversely $m=0$,
$C=\diag(c,1,\dots,1)$ with $c\uparrow\rho/2$ has $q_{\mathrm G}=2c<\rho$ and
$\Delta_{\mathrm G}\to\varphi(\rho/2)$.
\end{proof}

\begin{remark}[Rate versus region]
\label{rem:rate-region}
For Gaussian targets both region and rate are exact. The universal bound $\gund=1/(4+\Gamma)$ is
used only for general targets as a lower bound on $\gstar$; when the exact gap is available it
should be used in \eqref{eq:rho-core}. As $\delta\downarrow0$, $\rho_{\mathrm{core}}(\delta)\to2\gstar$,
the exact asymptotic squared-norm rate.
\end{remark}

\subsection{Deterministic discrete schemes: exact Gaussian intervals}
\label{subsec:disc-local}

\begin{proposition}[Exact Gaussian KL map and sharp interval]
\label{prop:kl-gauss}
Let the target be Gaussian, so \eqref{eq:KL-det} becomes
$m_{n+1}=(\Id-\Delta t\,C_n)m_n$, $C_{n+1}=(1+\Delta t)C_n(\Id+\Delta t\,C_n)^{-1}$. In the
eigenbasis of $C_n$ the scalar factors are
$r_m(c)=(1-\Delta t\,c)^2$, $r_C(c)=(1+\Delta t\,c)^{-2}$, $c^+=\tfrac{(1+\Delta t)c}{1+\Delta t\,c}$.
Fix $0<\Delta t<2$ and $q\in(0,1)$ with $q\geq q_0(\Delta t):=\max\{(1+\Delta t)^{-2},(1-\Delta t)^2\}$,
and set $I_{\Delta t}^{\mathrm{KL}}(q):=[\tfrac{q^{-1/2}-1}{\Delta t},\tfrac{1+\sqrt q}{\Delta t}]\ni1$.
Then $r_m(c),r_C(c)\leq q$ on $I_{\Delta t}^{\mathrm{KL}}(q)$; $c\mapsto c^+$ maps
$I_{\Delta t}^{\mathrm{KL}}(q)$ into itself; and if $\operatorname{spec}(C_n)\subseteq I_{\Delta t}^{\mathrm{KL}}(q)$
then $D_{\Delta t}(a_{n+k})\leq q^kD_{\Delta t}(a_n)$ for all $k\geq0$. For $\Delta t\leq\sqrt2$,
$q_0(\Delta t)=(1+\Delta t)^{-2}$.
\end{proposition}

\begin{proof}
The maps share the eigenbasis of $C_n$; $c^+-1=\tfrac{c-1}{1+\Delta t\,c}$ gives $r_C(c)=(1+\Delta t\,c)^{-2}$,
and the mean scales by $1-\Delta t\,c$. Then $r_m(c)\leq q\Leftrightarrow c\in[\tfrac{1-\sqrt q}{\Delta t},\tfrac{1+\sqrt q}{\Delta t}]$
and $r_C(c)\leq q\Leftrightarrow c\geq\tfrac{q^{-1/2}-1}{\Delta t}$; since $q^{-1/2}-1\geq1-\sqrt q$, the
covariance lower bound dominates, giving $I_{\Delta t}^{\mathrm{KL}}(q)$, and $1$ lies in it iff
$q\geq q_0(\Delta t)$. The map $c^+$ lies strictly between $c$ and $1$, so an interval containing $1$
is preserved; each scalar factor is $\leq q$, so $D_{\Delta t}$ contracts by $q$.
\end{proof}

\begin{proposition}[General-target discrete local convergence]
\label{thm:disc-local}
Fix a scheme $\bullet\in\{\Riem,\KL\}$ with certified linearized gap $\gamma_\bullet$
($\gamma_{\Riem}=\gund$, $\gamma_{\KL}=\tfrac1{4+2\Delta t+\Gamma}$) and norm $\norm\cdot_{\star,\bullet}$
($\norm\cdot_\star$ resp.\ $\norm\cdot_{\star,\Delta t}$), with associated norm-conversion constant
$\mathfrak b_{\star,\bullet}(\delta)$ defined by $\Delta(a)\leq\delta\Rightarrow\norm{a-a_\star}_{\star,\bullet}\leq\mathfrak b_{\star,\bullet}(\delta)\sqrt\delta$:
explicitly $\mathfrak b_{\star,\Riem}(\delta):=\mathfrak b(\delta)$ and
$\mathfrak b_{\star,\KL}(\delta):=\mathfrak b_{\Delta t}(\delta)$, both of \Cref{cor:band}, with the clean
comparison $\mathfrak b_{\Delta t}(\delta)\leq\sqrt{1+\Delta t}\,\mathfrak b(\delta)$ (immediate from
$\mathfrak b_{\Delta t}(\delta)^2=\max\{2/\alpha_\star,(1+\Delta t)\vartheta(\delta)\}\leq(1+\Delta t)\mathfrak b(\delta)^2$).
For the scheme map $T$ put
$\Kdisc(\delta):=\Delta t^{-1}\sup_{a\in\Hreg_\delta}\norm{D^2T(a)}_{\star,\bullet}$, finite by
\Cref{rem:Kdisc-finite}, where $\norm{D^2T(a)}_{\star,\bullet}:=\sup_{\norm\xi_{\star,\bullet}=1}\norm{D^2T(a)[\xi,\xi]}_{\star,\bullet}$
is the operator norm of the bilinear map $D^2T(a)$ in the scheme norm. Suppose
\[
    0<\Delta t\leq\min\Bigl\{\tfrac2{4+\Gamma},\tfrac1{\beta_\star\bar\Lambda_\delta}\Bigr\},
    \qquad
    \Kdisc(\delta)\,\mathfrak b_{\star,\bullet}(\delta)\sqrt\delta\leq\gamma_\bullet .
\]
If $\Delta(a_0)\leq\delta$, then the iterates remain in $\Ureg_\delta$ and
$\norm{a_n-a_\star}_{\star,\bullet}\leq(1-\tfrac12\Delta t\,\gamma_\bullet)^n\norm{a_0-a_\star}_{\star,\bullet}$.
Since $\Kdisc(\delta)\to\Kdisc(0)<\infty$ and $\mathfrak b_{\star,\bullet}(\delta)\sqrt\delta\to0$, the
radius condition holds for all small $\delta$.
\end{proposition}

\begin{proof}
Invariance: \Cref{lem:det-descent} resp.\ \Cref{prop:KL-onestep} give $\calE(a_{n+1})\leq\calE(a_n)$
under $\Delta t\leq1/(\beta_\star\bar\Lambda_\delta)$, so $a_{n+1}\in\Ureg_\delta$ by
\Cref{lem:region-geom}(iv). Write $\xi_n:=a_n-a_\star$ and use the integral (vector-valued) Taylor
form, valid without a single intermediate point,
\[
    a_{n+1}-a_\star=T(a_\star+\xi_n)-T(a_\star)
    =M_\bullet\xi_n+\int_0^1(1-s)\,D^2T(a_\star+s\xi_n)[\xi_n,\xi_n]\dd s,
    \qquad M_\bullet=DT(a_\star);
\]
the linear part contracts by $1-\Delta t\gamma_\bullet$ by \Cref{prop:linearized-maps}, and, since
$a_\star+s\xi_n\in\Hreg_\delta$ for $s\in[0,1]$ by the definition of the star-shaped hull $\Hreg_\delta$
(the segment from $a_\star$ to any $a\in\Hreg_\delta$ lies in $\Hreg_\delta$), the remainder is
bounded by $\int_0^1(1-s)\dd s\cdot\sup_{a\in\Hreg_\delta}\norm{D^2T(a)}_{\star,\bullet}\norm{\xi_n}_{\star,\bullet}^2
=\tfrac12\Delta t\Kdisc(\delta)\norm{\xi_n}_{\star,\bullet}^2\leq\tfrac12\Delta t\gamma_\bullet\norm{\xi_n}_{\star,\bullet}$.
Adding gives per-step factor $1-\tfrac12\Delta t\gamma_\bullet$.
\end{proof}

\begin{lemma}[One-step local map contraction before exit]
\label{lem:local-map-contraction}
Fix $\bullet\in\{\Riem,\KL\}$ and a level $\delta>0$ satisfying the radius condition
$\Kdisc(\delta)\,\mathfrak b_{\star,\bullet}(\delta)\sqrt\delta\leq\gamma_\bullet$ of \Cref{thm:disc-local}.
If $0<\Delta t\leq\tfrac2{4+\Gamma}$, then for \emph{every} $a\in\Ureg_\delta$ the deterministic scheme
map $T_\bullet$ satisfies
\[
    \norm{T_\bullet(a)-a_\star}_{\star,\bullet}\leq\Bigl(1-\tfrac12\Delta t\,\gamma_\bullet\Bigr)\norm{a-a_\star}_{\star,\bullet}.
\]
No energy-descent hypothesis is used; in particular the bound requires \emph{not}
$\Delta t\leq1/(\beta_\star\bar\Lambda_\delta)$, and no forward-invariance is asserted---this is purely
the one-step contraction of $T_\bullet$ at the single point $a$.
\end{lemma}

\begin{proof}
This is the contraction step of the proof of \Cref{thm:disc-local}, isolated. With $\xi:=a-a_\star$ and
$a\in\Ureg_\delta\subseteq\Hreg_\delta$, the integral Taylor identity
$T_\bullet(a)-a_\star=M_\bullet\xi+\int_0^1(1-s)D^2T_\bullet(a_\star+s\xi)[\xi,\xi]\dd s$ holds
($M_\bullet=DT_\bullet(a_\star)$). The linear part obeys
$\norm{M_\bullet\xi}_{\star,\bullet}\leq(1-\Delta t\gamma_\bullet)\norm\xi_{\star,\bullet}$ by
\Cref{prop:linearized-maps}, whose only stepsize requirement is $\Delta t\leq\tfrac2{4+\Gamma}$. Since
$a_\star+s\xi\in\Hreg_\delta$ for $s\in[0,1]$ (star-shaped hull), the remainder is bounded by
$\tfrac12\Delta t\,\Kdisc(\delta)\norm\xi_{\star,\bullet}^2\leq\tfrac12\Delta t\,\gamma_\bullet\norm\xi_{\star,\bullet}$,
using the radius condition and $\norm\xi_{\star,\bullet}\leq\mathfrak b_{\star,\bullet}(\delta)\sqrt\delta$.
Adding gives the factor $1-\tfrac12\Delta t\,\gamma_\bullet$. The invariance sentence of
\Cref{thm:disc-local}---the only place $\Delta t\leq1/(\beta_\star\bar\Lambda_\delta)$ enters---is not
invoked.
\end{proof}

\begin{remark}[Finiteness and continuity of $\Kdisc$]
\label{rem:Kdisc-finite}
The bound $\Kdisc(\delta)<\infty$ requires more than compactness of $\Hreg_\delta$: it needs
$D^2T$ \emph{continuous} on a neighborhood of the hull. Under \Cref{assump:logconcave-smooth} the
Gaussian-expectation maps $g(a)=-\E_{\rho_a}[\nabla V]$ and $A(a)=\E_{\rho_a}[\nabla^2V]$ are obtained
by differentiating $\E_{\rho_a}[\,\cdot\,]$ under the integral, which is licensed for observables of at
most polynomial growth (here $\nabla V$ grows linearly); each parameter derivative brings down a
Gaussian score polynomial, so $g$ and $A$ are $C^\infty$ in $(m,C)$ on every spectrally bounded compact
subset of $\Aspace$. Both scheme maps $T$ (the Riemannian matrix-exponential retraction and the KL
resolvent $C\mapsto(1+\Delta t)(C^{-1}+\Delta t A)^{-1}$) are smooth functions of $g,A,C$ with $C$
bounded away from $0$ and $\infty$ on $\Hreg_\delta$, hence $C^2$ there with $D^2T$ continuous. The
supremum of a continuous function over the compact hull is finite, giving $\Kdisc(\delta)<\infty$;
continuity in $\delta$ and shrinking of the hull as $\delta\downarrow0$ give
$\Kdisc(\delta)\to\Kdisc(0)<\infty$.
\end{remark}

\subsection{Online entry criteria}
\label{subsec:entry}

The energy test $\Delta(a_n)\leq\delta$ is not directly computable, since $\calE(a_\star)$ is
unknown. The following sufficient tests use only the current iterate and known target quantities,
in original coordinates. They are \emph{exact-expectation} tests: they presuppose the population
quantities $g(a)$ and $A(a)$; a Monte Carlo implementation would replace them by estimators and
require a concentration-adjusted (confidence) version of the gate. Moreover the threshold
$\delta_{\mathrm{loc}}$ they compare against is only \emph{implicitly} available in general (it must
first be certified admissible for \Cref{thm:disc-local}, whose constant $\Kdisc(\delta)$ is bounded
but not given in closed form), so the test is executable only after such a $\delta_{\mathrm{loc}}$
has been fixed. The conservative bounds of \Cref{rem:threshold-computable} give a precomputable
\emph{continuous-time} threshold; for the discrete schemes the admissible level remains implicit
until an explicit upper bound on $\Kdisc(\delta)$ is supplied.

\begin{proposition}[Exact-expectation residual gate, original coordinates]
\label{prop:entry-residual}
At $a=(m,C)$ set $g(a)=-\E_{\rho_a}[\nabla V(X)]$, $A(a)=\E_{\rho_a}[\nabla^2V(X)]$, and
$\mathcal R_{\mathrm{BW}}(a)^2:=\norm{g(a)}^2+\Tr[(C^{-1}-A(a))C(C^{-1}-A(a))]$. Then
$\mathcal R_{\mathrm{BW}}(a)^2\geq2\alpha\Delta(a)$, so for any a priori certified local level
$\delta_{\mathrm{loc}}$,
\[
    \mathcal R_{\mathrm{BW}}(a_n)^2\leq2\alpha\,\delta_{\mathrm{loc}}
    \;\Longrightarrow\;
    \Delta(a_n)\leq\delta_{\mathrm{loc}} .
\]
The left-hand side requires neither $C_\star$, $m_\star$, nor $\calE(a_\star)$. When a covariance
floor $C_n\succeq\ell_n\Id$ is available, the Fisher--Rao gradient norm gives an alternative
intrinsic test, $\Gradnorm(a_n)^2\leq\alpha\ell_n\delta_{\mathrm{loc}}\Rightarrow\Delta(a_n)\leq\delta_{\mathrm{loc}}$,
which by \Cref{lem:FR-W2} is more restrictive than the Bures--Wasserstein gate. For the continuous
flow and either deterministic discretization, once either test triggers the corresponding
energy-descent property keeps $\Delta(a_k)\leq\delta_{\mathrm{loc}}$ at all subsequent times or
iterations. For the stochastic schemes, neither pathwise nor expected containment follows from this
gate alone: the one-step recursion $\E[\Delta_{n+1}\mid\Fil_n]\leq q\Delta_n+\eta$ with $\eta>0$ gives
only $\E[\Delta_{n+1}\mid\Fil_n]\leq q\delta_{\mathrm{loc}}+\eta$, which is $\leq\delta_{\mathrm{loc}}$
only if $\delta_{\mathrm{loc}}\geq\eta/(1-q)$, i.e.\ if the threshold dominates the corresponding
stochastic floor. An expected-containment statement therefore requires either that condition or a
separate stochastic local-stability argument.
\end{proposition}

\begin{proof}
$\mathcal R_{\mathrm{BW}}(a)^2=\norm{\grad^{W_2}\calE(a)}_{W_2}^2\geq2\alpha\Delta(a)$ by
\Cref{lem:W2-PL}; the Fisher--Rao variant is \Cref{lem:FR-W2}; forward invariance is
\Cref{lem:region-geom}(iv).
\end{proof}

\begin{remark}[Precomputable conservative threshold]
\label{rem:threshold-computable}
The level $\delta_{\mathrm{loc}}$ is target-adaptive through
$\alpha_\star,\beta_\star,\Sigma_\delta,\mathcal B_\delta,\Kng$. A precomputable conservative
choice follows from the a priori bounds $\alpha_\star\geq\kappa^{-1}$, $\beta_\star\leq\kappa$,
$\Sigma_\delta\leq\sqrt{N_\theta}(\beta_\star-\alpha_\star)$,
$\mathcal B_\delta\leq\max\{\beta_\star-1,1-\alpha_\star\}$, which reduce the \emph{continuous-time}
admissibility condition $\rho_{\mathrm{core}}(\delta)>0$ of \Cref{thm:core-flow} to a scalar
inequality in $\alpha,\beta,N_\theta$. This does \emph{not} make the \emph{discrete} admissibility
condition $\Kdisc(\delta)\,\mathfrak b_{\star,\bullet}(\delta)\sqrt\delta\leq\gamma_\bullet$ explicit,
because no closed-form upper bound on $\Kdisc(\delta)=\Delta t^{-1}\sup_{a\in\Hreg_\delta}\norm{D^2T(a)}_{\star,\bullet}$
is supplied here; an explicit bound on $\Kdisc(\delta)$ in terms of the target parameters is still
required for a fully precomputable discrete switching threshold.
\end{remark}

\section{Removing the covariance burn-in: transport and quadratic rescue}
\label{sec:transport-bootstrap}

The Fisher--Rao schemes converge globally from any positive-definite initialization, incurring only
the logarithmic covariance burn-in $\logp\tfrac1{\beta\lambda_0}$ of
\Cref{thm:cont-global,thm:Riem-global,thm:KL-global}. When removing that stage is desirable, two
optional \emph{deterministic} initializers are available, and we develop both here. The first is a
short Bures--Wasserstein transport warm-start (\Cref{subsec:transport-warmstart}): it is
expectation-level, uses only $g(a)$ and $A(a)$, and requires no pointwise Hessian. The second, under
a pointwise Hessian oracle, is the \emph{quadratic rescue} (\Cref{subsec:rescue-det}); this is
precisely the single-query initializer that later drives the stochastic pipeline of
\Cref{sec:stoch-schemes}, and we give it here in its deterministic role. Neither initializer is part
of the mandatory algorithm and neither affects the local discretization factor; each merely deletes
the $\logp\tfrac1{\beta\lambda_0}$ stage. They differ in oracle model and in geometry, as recorded in
\Cref{rem:two-initializers}.

\subsection{Bures--Wasserstein transport warm-start}
\label{subsec:transport-warmstart}

For $\dot C=C-CAC$, an eigenvalue $\mu_i$ obeys $\dot\mu_i=\mu_i(1-\mu_ih_i)$, $h_i\in[\alpha,\beta]$;
in the underdispersed regime $\mu_i\ll1/\beta$ this gives $\dot\mu_i\approx\mu_i$ (multiplicative),
so reaching $1/\beta$ from $\lambda_0$ costs time $\asymp\log(1/(\beta\lambda_0))$. The Gaussian
Bures--Wasserstein flow
\begin{equation}
\label{eq:W-flow}
    \dot m_t=g(a_t),\qquad \dot C_t=2\Id-C_tA(a_t)-A(a_t)C_t
\end{equation}
gives $\dot\mu_i=2(1-\mu_ih_i)\approx2$ (additive, scale-independent). Transport is thus efficient
for the covariance warm-up $0\to O(1/\beta)$, Fisher--Rao for shape calibration $1/\beta\to1/\alpha$.
The flow \eqref{eq:W-flow} is not affine invariant, so the scale $1/\beta$ is stated in the given
coordinates.

\begin{lemma}[Wasserstein covariance bootstrap]
\label{lem:W-cov-bootstrap}
Under \Cref{assump:logconcave-smooth}, along \eqref{eq:W-flow},
$C_t\succeq\ell_W(t)\Id$ with $\ell_W(t):=\tfrac1\beta+(\lambda_0-\tfrac1\beta)e^{-2\beta t}$, and
$C_t\preceq\Lambda\Id$. Consequently, for $c\in(0,1)$,
$T_W^{\mathrm c}(c):=\tfrac1{2\beta}\log\tfrac{1-\beta\lambda_0}{1-c}$ if $\lambda_0<c/\beta$ (and $0$
otherwise) satisfies $C_{T_W^{\mathrm c}(c)}\succeq\tfrac c\beta\Id$; the warm-up time is $O(1/\beta)$,
with no logarithmic dependence on $\lambda_0$.
\end{lemma}

\begin{proof}
For an eigenvalue $\mu_i$, $\dot\mu_i=2(1-\mu_ih_i)\geq2(1-\beta\mu_i)$; scalar comparison with
$\dot\ell_W=2(1-\beta\ell_W)$, $\ell_W(0)=\lambda_0$, gives $\mu_i\geq\ell_W$. If $\mu_i\geq1/\alpha$
then $\dot\mu_i\leq2(1-\alpha\mu_i)\leq0$, so $C_t\preceq\Lambda\Id$. Finally $\ell_W(t)\geq c/\beta$
iff $e^{-2\beta t}\leq\tfrac{1-c}{1-\beta\lambda_0}$.
\end{proof}

\begin{lemma}[One-step Wasserstein covariance floor]
\label{lem:one-step-W}
Given $a=(m,C)$, $\eta>0$, put $M:=\Id-\eta A(a)$, $m^+:=m+\eta g(a)$, $\widetilde C:=MCM$, and
\begin{equation}
\label{eq:W-FB-cov}
    C^+:=\tfrac12\bigl(\widetilde C+2\eta\Id+[\widetilde C(\widetilde C+4\eta\Id)]^{1/2}\bigr).
\end{equation}
If $C\preceq\Lambda\Id$ and $0<\eta\leq1/\beta$, then $C^+\succeq\eta\Id$, $C^+\preceq\Lambda\Id$, and
the Bures--Wasserstein forward--backward descent inequality gives $\calE(m^+,C^+)\leq\calE(m,C)$.
\end{lemma}

\begin{proof}
$\psi_\eta(x):=\tfrac12(x+2\eta+\sqrt{x(x+4\eta)})$ is increasing with $\psi_\eta(0)=\eta$, and
\eqref{eq:W-FB-cov} applies it to the eigenvalues of $\widetilde C\succeq0$, so $C^+\succeq\eta\Id$.
Since $0\preceq M\preceq(1-\eta\alpha)\Id$, $\widetilde C\preceq(1-\eta\alpha)^2\Lambda\Id\preceq(1-\eta/\Lambda)^2\Lambda\Id$;
$\psi_\eta((1-\eta/L)^2L)=L$ at $L=\Lambda$ with monotonicity gives $C^+\preceq\Lambda\Id$. The step
is the Bures--Wasserstein forward--backward step; under $\eta\leq1/\beta$ the standard descent
inequality \citep{diao2023forward} gives energy decrease.
\end{proof}

\begin{theorem}[Discrete transport-warm-started natural gradient; deterministic]
\label{thm:disc-W-to-FR}
Assume \Cref{assump:logconcave-smooth}. Fix $c\in(0,1]$, $\eta:=c/\beta$. If $\lambda_0<c/\beta$
perform one step \eqref{eq:W-FB-cov}; otherwise skip it. Then run \eqref{eq:Riem-det} with
$0<\Delta t\leq\min\{1/\LB,\,2/(4+\Gamma),\,1/(\beta_\star\bar\Lambda_{\delta_{\Riem}})\}$. For small
$\epsilon>0$,
\[
    \boxed{
    N_{W\to\Riem}(\epsilon;c)
    =\mathbf 1_{\{\lambda_0<c/\beta\}}
    +O\!\left(\frac\kappa{c\,\Delta t}\logp\frac{\Delta_0}{\delta_{\Riem}}+\frac{4+\Gamma}{\Delta t}\log\frac1\epsilon\right),
    }
\]
with $\delta_{\Riem}$ a level admissible for \Cref{thm:disc-local}. The same holds for
\eqref{eq:KL-det} with $\delta_{\Riem}\to\delta_{\KL}$ and $4+\Gamma\to4+2\Delta t+\Gamma$.
\end{theorem}

\begin{proof}
If taken, \Cref{lem:one-step-W} gives $C_{\mathrm b}\succeq\tfrac c\beta\Id$, $C_{\mathrm b}\preceq\Lambda\Id$,
$\Delta(a_{\mathrm b})\leq\Delta_0$. The covariance floors of \Cref{lem:Riem-cov,lem:KL-cov} started
from $c/(2\beta)$ are nondecreasing, so \Cref{thm:Riem-global}(1) with frozen floor $c/(2\beta)$
gives $\Delta_{n+1}\leq(1-\tfrac{c\Delta t}{4\kappa})\Delta_n$ (KL: $c\Delta t/(16\kappa)$), whence
localization to $\delta_{\Riem}$ costs $O(\tfrac\kappa{c\Delta t}\logp\tfrac{\Delta_0}{\delta_{\Riem}})$;
then \Cref{thm:disc-local}.
\end{proof}

\begin{remark}
\label{rem:transport-caveat}
The transport warm-start removes the small-initial-covariance burn-in. It does not improve the
local discretization factor $(4+\Gamma)/\Delta t$; the local stiffness is separate from the global
covariance warm-up.
\end{remark}

\subsection{Deterministic quadratic rescue}
\label{subsec:rescue-det}

The second burn-in remover is available whenever a pointwise Hessian can be queried at a single
point. It is exactly the initializer used by the stochastic pipeline of \Cref{sec:stoch-schemes};
we state and analyse it here in its deterministic role.

\begin{definition}[Quadratic rescue]
\label{def:rescue}
Given $m_{\mathrm{in}}\in\R^{N_\theta}$, query the oracle at $m_{\mathrm{in}}$:
$G_{\mathrm{in}}:=\nabla V(m_{\mathrm{in}})$, $H_{\mathrm{in}}:=\nabla^2V(m_{\mathrm{in}})$, and set
\begin{equation}
\label{eq:rescue}
    \boxed{\;
    m_0:=m_{\mathrm{in}}-H_{\mathrm{in}}^{-1}G_{\mathrm{in}},
    \qquad
    C_0:=H_{\mathrm{in}}^{-1}.
    \;}
\end{equation}
\end{definition}

\begin{lemma}[Curvature-scale covariance rescue]
\label{lem:rescue-band}
Under \Cref{assump:logconcave-smooth}, \eqref{eq:rescue} gives $\tfrac1\beta\Id\preceq C_0\preceq\tfrac1\alpha\Id$.
\end{lemma}

\begin{proof}
$\alpha\Id\preceq H_{\mathrm{in}}\preceq\beta\Id$ and the matrix inverse is operator antitone on
$\SPD(N_\theta)$, so $\tfrac1\beta\Id\preceq H_{\mathrm{in}}^{-1}\preceq\tfrac1\alpha\Id$, i.e.\
$\tfrac1\beta\Id\preceq C_0\preceq\tfrac1\alpha\Id$.
\end{proof}

\begin{theorem}[Exact one-query recovery of Gaussian targets]
\label{thm:gauss-recovery}
Suppose $V(x)=\tfrac12(x-\mu)^\top H(x-\mu)+c$ with $H\succ0$. Then \eqref{eq:rescue} gives $m_0=\mu$,
$C_0=H^{-1}$, so $a_0=a_\star$ and $\Delta(a_0)=0$, for both the Riemannian and KL schemes. Over the
class of unknown Gaussian targets, the oracle complexity is exactly one.
\end{theorem}

\begin{proof}
For a quadratic potential $\nabla V(m_{\mathrm{in}})=H(m_{\mathrm{in}}-\mu)$ and $\nabla^2V(m_{\mathrm{in}})=H$,
so $m_0=m_{\mathrm{in}}-H^{-1}H(m_{\mathrm{in}}-\mu)=\mu$ and $C_0=H^{-1}$. The target belongs to the
Gaussian family, so $a_\star=(\mu,H^{-1})=a_0$. At $a_\star$ every STL score sample vanishes
pathwise, $-\nabla V(X)+C_\star^{-1}(X-m_\star)=-H(X-\mu)+H(X-\mu)=0$, and $\nabla^2V(X)=H$ for all
$X$, so both covariance updates preserve $H^{-1}$. Necessity of one query: an algorithm making zero
queries is independent of $(\mu,H)$; one query at any $x$ returns $H(x-\mu)$ and $H$, from which
$\mu=x-H^{-1}G$, $C_\star=H^{-1}$.
\end{proof}

\begin{theorem}[Quadratic-rescued deterministic natural gradient; no covariance burn-in]
\label{thm:disc-rescue-to-FR}
Assume \Cref{assump:logconcave-smooth}, and initialize by the quadratic rescue of
\Cref{def:rescue}. By \Cref{lem:rescue-band} then $\tfrac1\beta\Id\preceq C_0\preceq\tfrac1\alpha\Id$,
so $\lambda_0:=\lambda_{0,\min}\geq\tfrac1\beta$. Consequently, in the covariance envelopes of
\Cref{lem:Riem-cov,lem:KL-cov}, $L_0:=\min\{\lambda_0,\tfrac1{2\beta}\}=\tfrac1{2\beta}$ and
$y_0:=(\beta L_0)^{-1}=2$, so the covariance floor $\tfrac1{2\beta}\Id$ holds from $n=0$ and the
burn-in count is $N_{\mathrm{cov},-}^{\Riem}=N_{\mathrm{cov},-}^{\KL}=0$. Running \eqref{eq:Riem-det}
(resp.\ \eqref{eq:KL-det}) with
$0<\Delta t\leq\min\{1/\LB,\,2/(4+\Gamma),\,1/(\beta_\star\bar\Lambda_{\delta_\bullet})\}$, for small
$\epsilon>0$,
\[
    \boxed{\;
    N_{\mathrm{rescue}\to\bullet}(\epsilon)
    =1
    +O\!\left(\frac\kappa{\Delta t}\logp\frac{\Delta_0}{\delta_\bullet}
    +\frac{c_\bullet}{\Delta t}\log\frac1\epsilon\right),
    \qquad
    c_{\Riem}=4+\Gamma,\quad c_{\KL}=4+2\Delta t+\Gamma,
    \;}
\]
where the additive $1$ counts the single deterministic Hessian query, $\delta_\bullet$ is a level
admissible for \Cref{thm:disc-local}, and $\Delta_0:=\Delta(a_0)$ is finite (bounded in
\Cref{lem:rescue-energy} under \Cref{assump:global-Hess-Lip}, and $\Delta_0=0$ for Gaussian targets
by \Cref{thm:gauss-recovery}). In particular the $\logp\tfrac1{\beta\lambda_0}$ burn-in stage of
\Cref{thm:det-three-stage} is eliminated.
\end{theorem}

\begin{proof}
\emph{No burn-in.} By \Cref{lem:rescue-band}, $C_0\succeq\tfrac1\beta\Id\succ\tfrac1{2\beta}\Id$, so
$\lambda_0\geq\tfrac1\beta$ and, in the envelope variable $y_n:=(\beta L_n)^{-1}$ of
\Cref{lem:Riem-cov,lem:KL-cov} with $L_0=\min\{\lambda_0,\tfrac1{2\beta}\}=\tfrac1{2\beta}$, one has
$y_0=2$. For the Riemannian scheme the exact envelope solution \eqref{eq:Riem-y} is
$y_n=y_\infty+(y_0-y_\infty)e^{-n\Delta t}$ with $y_\infty\in[1,2)$; since $y_0=2>y_\infty$ and
$e^{-n\Delta t}\leq1$, the sequence decreases monotonically from $2$ to $y_\infty$, so $y_n\leq2$ for
all $n$. For the KL scheme \Cref{lem:KL-cov} gives $y_n=1+(y_0-1)(1+\Delta t)^{-n}=1+(1+\Delta t)^{-n}$,
which decreases from $2$ to $1$, so again $y_n\leq2$. In both cases $y_n\leq2$ means
$C_n\succeq\tfrac1{2\beta}\Id$ for all $n\geq0$; equivalently, in \eqref{eq:Ncov-minus} the count is
$N_{\mathrm{cov},-}=\lceil\Delta t^{-1}\logp\tfrac{y_0-y_\infty}{2-y_\infty}\rceil=\lceil\Delta t^{-1}\logp1\rceil=0$
(Riemannian) and likewise $0$ for KL. Hence the covariance floor is already attained and
$N_{\mathrm{cov},-}=0$.

\emph{Global localization.} With the covariance frozen at floor $\tfrac1{2\beta}$,
\Cref{thm:Riem-global}(2) (resp.\ \Cref{thm:KL-global}(2)) gives the geometric per-step contraction
$\Delta_{n+1}\leq(1-\tfrac{\Delta t}{4\kappa})\Delta_n$ (resp.\ $1-\tfrac{\Delta t}{16\kappa}$), so
reaching $\Delta(a)\leq\delta_\bullet$ from $\Delta_0$ costs
$O(\tfrac\kappa{\Delta t}\logp\tfrac{\Delta_0}{\delta_\bullet})$ steps.

\emph{Local phase.} Once $\Delta(a)\leq\delta_\bullet$ we have $a\in\Ureg_{\delta_\bullet}$, and
\Cref{thm:disc-local} contracts $\norm{a-a_\star}_{\star,\bullet}$ at per-step rate
$1-\tfrac12\Delta t\gamma_\bullet$ with $\gamma_{\Riem}=\gund=1/(4+\Gamma)$,
$\gamma_{\KL}=1/(4+2\Delta t+\Gamma)$, so $O(\tfrac{c_\bullet}{\Delta t}\log\tfrac1\epsilon)$ further
steps give $\norm{a-a_\star}_\star^2\leq\epsilon$, using $-\log(1-x)\geq x$. Finiteness of $\Delta_0$
is \Cref{lem:rescue-energy}. The single Hessian query is the additive $1$.
\end{proof}

\begin{remark}[Two initializers, one oracle distinction]
\label{rem:two-initializers}
Both warm-starts delete the $\logp\tfrac1{\beta\lambda_0}$ stage but sit in different oracle models
and geometries. The transport warm-start (\Cref{thm:disc-W-to-FR}) is \emph{expectation-level}: it
uses $g(a)$ and $A(a)=\E_{\rho_a}[\nabla^2V]$ and inflates the covariance additively along the
Bures--Wasserstein forward--backward step, so it never queries a pointwise Hessian and only guarantees
a covariance \emph{floor} $C_0\succeq\tfrac c\beta\Id$ (the lower band edge; $c=1$ gives $\tfrac1\beta\Id$),
without matching the pointwise curvature. The quadratic rescue (\Cref{thm:disc-rescue-to-FR}) is a single
\emph{pointwise-Hessian} (local-curvature/Newton) query and lands directly on the full curvature
band $[\tfrac1\beta,\tfrac1\alpha]$, exact for Gaussian targets. Thus transport trades a weaker
oracle for a coarser landing scale, while the rescue trades one Hessian evaluation for exact
Gaussian recovery; the ensuing Fisher--Rao iterations are identical.
\end{remark}

\section{Deterministic three-stage complexity}
\label{sec:three-stage}

\begin{theorem}[Deterministic three-stage complexity]
\label{thm:det-three-stage}
Assume \Cref{assump:logconcave-smooth} and $\lambda_{0,\min}\Id\preceq C_0\preceq\lambda_{0,\max}\Id$.
Let $\delta_{\mathrm c}>0$ be admissible for \Cref{thm:core-flow} with
$\rho_{\mathrm{core}}(\delta_{\mathrm c})\geq\gund$, and $\delta_{\Riem},\delta_{\KL}>0$ admissible for
\Cref{thm:disc-local}; such levels exist by \Cref{cor:core-universal} and \Cref{thm:disc-local}.
For sufficiently small $\epsilon>0$:
\begin{enumerate}
\item \textup{(Continuous flow.)}
$T_{\mathrm c}(\epsilon)=O\bigl(\logp\tfrac1{\beta\lambda_0}+\kappa\logp\tfrac{\Delta_0}{\delta_{\mathrm c}}+(4+\Gamma)\log\tfrac1\epsilon\bigr)$
gives $\norm{a_T-a_\star}_\star^2\leq\epsilon$.
\item \textup{(Riemannian, $0<\Delta t\leq\min\{1/\LB,2/(4+\Gamma),1/(\beta_\star\bar\Lambda_{\delta_{\Riem}})\}$.)}
$N_{\Riem}(\epsilon)=O\bigl(\tfrac1{\Delta t}\logp\tfrac1{\beta\lambda_0}+\tfrac\kappa{\Delta t}\logp\tfrac{\Delta_0}{\delta_{\Riem}}+\tfrac{4+\Gamma}{\Delta t}\log\tfrac1\epsilon\bigr)$.
\item \textup{(KL, same window.)}
$N_{\KL}(\epsilon)=O\bigl(\tfrac1{\Delta t}\logp\tfrac1{\beta\lambda_0}+\tfrac\kappa{\Delta t}\logp\tfrac{\Delta_0}{\delta_{\KL}}+\tfrac{4+2\Delta t+\Gamma}{\Delta t}\log\tfrac1\epsilon\bigr)$.
\end{enumerate}
\end{theorem}

\begin{proof}
By \Cref{thm:cont-global} (resp.\ \Cref{thm:Riem-global}, \Cref{thm:KL-global}) with $\delta=\delta_\bullet$,
one reaches $\Delta(a)\leq\delta_\bullet$, hence $a\in\Ureg_{\delta_\bullet}$ and
$\norm{a-a_\star}_\star\leq\mathfrak b(\delta_\bullet)\sqrt{\delta_\bullet}$; \Cref{thm:core-flow} resp.\
\Cref{thm:disc-local}, restarted, then needs $O((4+\Gamma)\log\tfrac1\epsilon)$ further \emph{continuous
time} (\Cref{thm:core-flow}) or, discretely (\Cref{thm:disc-local}, per-step factor $1-\tfrac{\Delta t}{2(4+\Gamma)}$
Riemannian), $O\bigl(\tfrac{4+\Gamma}{\Delta t}\log\tfrac1\epsilon\bigr)$ further \emph{iterations}
(and $\tfrac{4+2\Delta t+\Gamma}{\Delta t}$ for KL), using $-\log(1-x)\geq x$.
\end{proof}

\begin{corollary}[Burn-in-free three-stage complexity]
\label{cor:burnin-free}
Under \Cref{assump:logconcave-smooth}, initialize by either the quadratic rescue
(\Cref{thm:disc-rescue-to-FR}) or the transport warm-start (\Cref{thm:disc-W-to-FR}) \emph{with target
$c=1$} (i.e.\ $\eta=1/\beta$, at most one forward--backward step). Both land $C_0\succeq\tfrac1\beta\Id$,
so $\lambda_0\geq\tfrac1\beta$, and the first stage of \Cref{thm:det-three-stage} is removed: for both
discretizations
\[
    N_\bullet(\epsilon)=n_{\mathrm{init}}
    +O\!\left(\frac\kappa{\Delta t}\logp\frac{\Delta_0}{\delta_\bullet}
    +\frac{c_\bullet}{\Delta t}\log\frac1\epsilon\right),
    \qquad c_{\Riem}=4+\Gamma,\quad c_{\KL}=4+2\Delta t+\Gamma,
\]
where $n_{\mathrm{init}}\in\{0,1\}$ is the warm-start ($1$ Hessian query for the rescue, or
$\mathbf 1_{\{\lambda_0<1/\beta\}}\leq1$ forward--backward steps for transport) and $\Delta_0=\Delta(a_0)$
is the energy of the warm-started initializer; for Gaussian targets the rescue gives $\Delta_0=0$ and
the localization stage vanishes as well.
\end{corollary}

\begin{proof}
For the rescue, \Cref{lem:rescue-band} gives $C_0\succeq\tfrac1\beta\Id$; for transport with $c=1$,
\Cref{lem:one-step-W} with $\eta=1/\beta$ gives $C^+\succeq\tfrac1\beta\Id$ in one step (skipped when
$\lambda_0\geq1/\beta$ already). In both cases $\lambda_0\geq\tfrac1\beta>\tfrac1{2\beta}$, so
$L_0=\min\{\lambda_0,\tfrac1{2\beta}\}=\tfrac1{2\beta}$ and $y_0=(\beta L_0)^{-1}=2$. By the exact
envelope solutions \eqref{eq:Riem-y} (Riemannian, $y_n=y_\infty+(2-y_\infty)e^{-n\Delta t}\leq2$,
$y_\infty<2$) and \Cref{lem:KL-cov} (KL, $y_n=1+(1+\Delta t)^{-n}\leq2$), one has $C_n\succeq\tfrac1{2\beta}\Id$
for all $n\geq0$, i.e.\ $N_{\mathrm{cov},-}=0$ from \eqref{eq:Ncov-minus} (as in
\Cref{thm:disc-rescue-to-FR}). Dropping the now-empty burn-in stage from
\Cref{thm:det-three-stage}(2)--(3) leaves the displayed bound; the Gaussian case is
\Cref{thm:gauss-recovery}.
\end{proof}

\begin{remark}[Deterministic local restart without the descent stepsize]
\label{rem:local-restart}
The window in \Cref{thm:det-three-stage}(2)--(3) still carries the energy-descent constraint
$\Delta t\leq1/(\beta_\star\bar\Lambda_{\delta_\bullet})$, which can be $O(\kappa_\star^{-1})$. This is
removable in the \emph{local} phase alone, provided the certified norm ball is contained in the energy
sublevel. Concretely, fix a level $\delta_\bullet$ admissible for \Cref{lem:local-map-contraction} and
choose a radius $r_\bullet>0$ with
\begin{equation}
\label{eq:ball-in-sublevel}
    \mathbb B_{\star,\bullet}(a_\star,r_\bullet)
    :=\bigl\{a:\norm{a-a_\star}_{\star,\bullet}\leq r_\bullet\bigr\}
    \subseteq\Ureg_{\delta_\bullet};
\end{equation}
such $r_\bullet$ exists because $\Delta$ is continuous with $\Delta(a_\star)=0<\delta_\bullet$, so
$a_\star$ lies in the interior of the sublevel set $\Ureg_{\delta_\bullet}=\{\Delta\leq\delta_\bullet\}$.
(The containment \eqref{eq:ball-in-sublevel} is the same condition the stochastic containment corollary
imposes; it is what a norm ball does \emph{not} satisfy automatically.) With this choice, if the
restart iterate $a_{n_0}\in\mathbb B_{\star,\bullet}(a_\star,r_\bullet)$ then, by
\eqref{eq:ball-in-sublevel}, $a_{n_0}\in\Ureg_{\delta_\bullet}$, so \Cref{lem:local-map-contraction}
applies at $a_{n_0}$ and gives, for $\Delta t_{\mathrm{loc}}\leq2/(4+\Gamma)$,
\[
    \norm{a_{n_0+1}-a_\star}_{\star,\bullet}
    \leq\bigl(1-\tfrac12\Delta t_{\mathrm{loc}}\gamma_\bullet\bigr)\norm{a_{n_0}-a_\star}_{\star,\bullet}
    \leq r_\bullet,
\]
whence $a_{n_0+1}\in\mathbb B_{\star,\bullet}(a_\star,r_\bullet)\subseteq\Ureg_{\delta_\bullet}$ and the
lemma reapplies; induction closes. Ball invariance thus follows from contraction \emph{together with}
\eqref{eq:ball-in-sublevel}, not from contraction alone---the ball, not the sublevel, is the invariant
set, and \eqref{eq:ball-in-sublevel} is what keeps it inside the region where the one-step lemma is
valid. No energy-descent bound is used. Restarting the local iteration at $\Delta t_{\mathrm{loc}}$
therefore replaces the third-stage cost $\tfrac{4+\Gamma}{\Delta t}\log\tfrac1\epsilon$ (throttled by
the global $\Delta t$) with $\tfrac{(4+\Gamma)^2}2\log\tfrac1\epsilon$ at the local step, matching the
stochastic accounting of \Cref{rem:local-condnum}. Only the global stages then require the descent
stepsize.

To reach the ball from the global phase (rather than assume $a_{n_0}\in\mathbb B_{\star,\bullet}$),
introduce an inner level $\delta_{\mathrm{in}}\leq\delta_\bullet$ with
\begin{equation}
\label{eq:inner-level-reach}
    \mathfrak b_{\star,\bullet}(\delta_{\mathrm{in}})\sqrt{\delta_{\mathrm{in}}}\leq r_\bullet ,
\end{equation}
which by the norm conversion of \Cref{cor:band} gives the nesting
$\Ureg_{\delta_{\mathrm{in}}}\subseteq\mathbb B_{\star,\bullet}(a_\star,r_\bullet)\subseteq\Ureg_{\delta_\bullet}$
(such $\delta_{\mathrm{in}}$ exists since $\mathfrak b_{\star,\bullet}(\delta)\sqrt\delta\to0$ as
$\delta\downarrow0$). The complete strategy is then: run the descent-controlled global/local iteration
of \Cref{thm:det-three-stage} until $\Delta(a_{n_0})\leq\delta_{\mathrm{in}}$, i.e.\ until
$a_{n_0}\in\Ureg_{\delta_{\mathrm{in}}}\subseteq\mathbb B_{\star,\bullet}(a_\star,r_\bullet)$, then
raise the stepsize to $\Delta t_{\mathrm{loc}}$ and iterate as above. The entrance cost of reaching
$\Ureg_{\delta_{\mathrm{in}}}$ is independent of the final accuracy $\epsilon$ (it depends only on
$\Delta_0/\delta_{\mathrm{in}}$), so the leading local term is unchanged at
$\tfrac{(4+\Gamma)^2}2\log\tfrac1\epsilon$.

The restart also admits a \emph{target-adaptive} form via \eqref{eq:linear-gap-bounds}, which
however differs by scheme because the two linearized maps are self-adjoint in different norms
(\Cref{prop:linearized-maps}): the Riemannian map in $\norm\cdot_\star$ and the KL map in the weighted
$\norm\cdot_{\star,\Delta t_{\mathrm{loc}}}$, whereas the spectral pair $(\gstar,\Lambda_\star)$ of
\eqref{eq:linear-gap-bounds} is defined in $\norm\cdot_\star$. Fix any certified pair
$0<\gamma_0\leq\gstar$, $\Lambda_0\geq\Lambda_\star$.
\emph{Riemannian.} Since $DT_{\Riem,h}(a_\star)=\Id-h\Lstar$ is self-adjoint in $\norm\cdot_\star$ with
$Q(\xi)\geq\gamma_0\norm\xi_\star^2$, it contracts by $1-h\gamma_0$ under $h\leq1/\Lambda_0$; with the
radius condition $\Kdisc(\delta_\bullet)r_\bullet\leq\gamma_0$ and \eqref{eq:ball-in-sublevel} the
per-step factor is $1-\tfrac12h\gamma_0$ and the cost $O\bigl((\Lambda_0/\gamma_0)\log\tfrac1\epsilon\bigr)$.
\emph{KL.} Here $M_{\KL,h}$ is self-adjoint in $\norm\cdot_{\star,h}$; converting the $\norm\cdot_\star$
bounds through $\norm\xi_\star^2\leq\norm\xi_{\star,h}^2\leq(1+h)\norm\xi_\star^2$ gives
$Q(\xi)\geq\tfrac{\gamma_0}{1+h}\norm\xi_{\star,h}^2$ and $Q(\xi)\leq\Lambda_0\norm\xi_{\star,h}^2$, so
under $h\leq1/\Lambda_0$ the KL map contracts by $1-\tfrac{h\gamma_0}{1+h}$; with the correspondingly
tightened radius condition $\Kdisc(\delta_\bullet)r_\bullet\leq\gamma_0/(1+h)$ the per-step factor is
$1-\tfrac{h\gamma_0}{2(1+h)}$ and the cost $O\bigl(\tfrac{(1+h)\Lambda_0}{\gamma_0}\log\tfrac1\epsilon\bigr)$.
Since $h\leq1/\Lambda_0\leq1$ makes $1+h\leq2$, both costs are $O\bigl((\Lambda_0/\gamma_0)\log\tfrac1\epsilon\bigr)$
in order---the $(1+h)$ factor changes only the exact KL constant---which is
$O\bigl((4+\Gamma)^2\log\tfrac1\epsilon\bigr)$ under the universal bounds but
$O(\kappa_\star\log\tfrac1\epsilon)$ (via $\gamma_0=\alpha_\star$, $\Lambda_0=\beta_\star$) for a
well-conditioned target with a pessimistic $\Gamma$.
\end{remark}

\section{The Price/Hessian oracle and the stochastic schemes}
\label{sec:stoch-schemes}

\begin{definition}[Price/Hessian oracle]
\label{def:oracle}
Given $X\in\R^{N_\theta}$, the oracle returns the pair $\bigl(\nabla V(X),\,\nabla^2V(X)\bigr)$.
\end{definition}

Fix a batch size $B\geq1$. Let $\Fil_n:=\sigma(a_0,\ldots,a_n)$, $\xi_{n,s}\overset{\text{iid}}{\sim}\N(0,\Id)$,
$X_{n,s}:=m_n+C_n^{1/2}\xi_{n,s}$, $H_{n,s}:=\nabla^2V(X_{n,s})$.

\paragraph{Mean/score estimator: STL.}
\begin{equation}
\label{eq:STL-score}
    \overline b_n:=\frac1B\sum_{s=1}^B\bigl(-\nabla V(X_{n,s})+C_n^{-1}(X_{n,s}-m_n)\bigr)
    =\frac1B\sum_{s=1}^B\bigl(-\nabla V(X_{n,s})+C_n^{-1/2}\xi_{n,s}\bigr).
\end{equation}

\paragraph{Covariance estimator: Price/Hessian.}
$\overline H_n:=\tfrac1B\sum_{s=1}^BH_{n,s}$, $\overline K_n:=C_n^{-1}-\overline H_n$. Both are
conditionally unbiased,
\begin{equation}
\label{eq:STL-unbiased}
    \E[\overline b_n\mid\Fil_n]=g(a_n),
    \qquad
    \E[\overline K_n\mid\Fil_n]=R(a_n).
\end{equation}

\paragraph{The global update.}
The \emph{minibatch Price/Hessian Riemannian scheme} is
\begin{equation}
\label{eq:Riem-STL}
    m_{n+1}=m_n+\Delta t\,C_n\overline b_n,
    \qquad
    C_{n+1}=C_n^{1/2}\exp\bigl(\Delta t(\Id-C_n^{1/2}\overline H_nC_n^{1/2})\bigr)C_n^{1/2},
\end{equation}
and the \emph{minibatch Price/Hessian KL scheme} is
\begin{equation}
\label{eq:KL-STL}
    m_{n+1}=m_n+\Delta t\,C_n\overline b_n,
    \qquad
    C_{n+1}=(1+\Delta t)\bigl(C_n^{-1}+\Delta t\,\overline H_n\bigr)^{-1}.
\end{equation}
STL enters through the score $\overline b_n$; the covariance part is Price/Hessian. Since each
$H_{n,s}$ has spectrum in $[\alpha,\beta]$ a.s., so does the average:
\begin{equation}
\label{eq:pathwise-Hess}
    \alpha\Id\preceq\overline H_n\preceq\beta\Id\qquad\text{almost surely.}
\end{equation}
This almost-sure spectral bound is the structural reason the Price/Hessian oracle admits an
unconditional global theory.

\begin{remark}[Price's identity and the choice of covariance oracle]
\label{rem:price}
By Gaussian integration by parts, for $X=m+C^{1/2}\xi$, $\E[\xi\,\nabla V(X)^\top]=C^{1/2}A(a)$, so
$A(a)$ admits both a sampled-Hessian and a gradient-only representation. We use the sampled Hessian:
it is the representation with the almost-sure spectral bound \eqref{eq:pathwise-Hess}, which is
exactly what the pathwise confinement arguments require. The gradient-only representation has
unbounded samplewise spectrum; none of the pathwise statements below apply to it.
\end{remark}

\subsection{The quadratic rescue}
\label{subsec:rescue}

The stochastic pipeline is initialized by the single deterministic Price/Hessian query of the
\emph{quadratic rescue} (\Cref{def:rescue}), developed for the deterministic theory in
\Cref{subsec:rescue-det}: it sets $m_0=m_{\mathrm{in}}-H_{\mathrm{in}}^{-1}G_{\mathrm{in}}$,
$C_0=H_{\mathrm{in}}^{-1}$, places the covariance on the curvature band
$\tfrac1\beta\Id\preceq C_0\preceq\tfrac1\alpha\Id$ (\Cref{lem:rescue-band}), and recovers the exact
optimizer in one query for Gaussian targets (\Cref{thm:gauss-recovery}). We record here the energy
quality of the rescued initializer for non-Gaussian targets, used by the global stochastic
complexity bounds.

\begin{lemma}[Rescue energy for nonquadratic targets]
\label{lem:rescue-energy}
Under \Cref{assump:logconcave-smooth}, \Cref{lem:rescue-band} places $C_0$ on the band
$[\tfrac1\beta,\tfrac1\alpha]$ pathwise, and $\Delta(a_0)$ is a deterministic finite quantity. Under
the additional Hessian-Lipschitz \Cref{assump:global-Hess-Lip}, for the point-Newton rescue
$H_{\mathrm{in}}:=\nabla^2V(m_{\mathrm{in}})$, $C_0=H_{\mathrm{in}}^{-1}$,
$m_0=m_{\mathrm{in}}-H_{\mathrm{in}}^{-1}\nabla V(m_{\mathrm{in}})$ of Phase~0
(\Cref{alg:unified}), writing $r_{\mathrm{in}}:=\norm{m_{\mathrm{in}}-m_\star}$,
\[
    \norm{m_0-m_\star}\leq\frac{\LH}{2\alpha}\Bigl(r_{\mathrm{in}}^2+\frac{N_\theta}\alpha\Bigr),
    \qquad
    \norm{C_0-C_\star}_{\Frob}\leq\frac{\sqrt{N_\theta}\,\LH}{\alpha^2}\Bigl(r_{\mathrm{in}}+\sqrt{N_\theta/\alpha}\Bigr),
\]
and hence
\[
    \Delta(a_0)
    \;\leq\;
    \frac{\beta\LH^2}{8\alpha^2}\Bigl(r_{\mathrm{in}}^2+\frac{N_\theta}\alpha\Bigr)^2
    +\frac{\beta^2N_\theta\LH^2}{4\alpha^4}\Bigl(r_{\mathrm{in}}+\sqrt{N_\theta/\alpha}\Bigr)^2,
\]
which is \emph{quadratic} in $\LH$ and vanishes when $\LH=0$. (One may also retain the first-order
covariance-segment estimate
$\Delta(a_0)\leq\tfrac\beta2\norm{m_0-m_\star}^2+\tfrac\beta2\sqrt{N_\theta}\norm{C_0-C_\star}_{\Frob}$
and take the minimum of the two; the quadratic form is the one with the correct near-Gaussian
$O(\LH^2)$ scaling, matching the $O(\LHs^2)$ stochastic local floors.)
\end{lemma}

\begin{proof}
Write $d:=m_{\mathrm{in}}-m_\star$ and $\bar H:=\int_0^1\nabla^2V(m_\star+td)\dd t$, so that
$\nabla V(m_{\mathrm{in}})=\nabla V(m_\star)+\bar H d$ by the integral Taylor form. Note $m_\star$ is the
optimal variational \emph{mean}, not the target mode: the stationarity condition is
$\E_{\rho_\star}[\nabla V]=0$, which need not force $\nabla V(m_\star)=0$.

\emph{Bound on $\nabla V(m_\star)$.} With $X=m_\star+C_\star^{1/2}Z$ and Taylor expansion,
$\nabla V(X)=\nabla V(m_\star)+\nabla^2V(m_\star)C_\star^{1/2}Z+r(X)$ with $\norm{r(X)}\leq\tfrac{\LH}2\norm{C_\star^{1/2}Z}^2$.
Taking $\E_{\rho_\star}$, the linear term vanishes ($\E Z=0$) and $\E_{\rho_\star}[\nabla V]=0$, so
\begin{equation}
\label{eq:gradV-mstar}
    \norm{\nabla V(m_\star)}\leq\tfrac{\LH}2\E\norm{C_\star^{1/2}Z}^2=\tfrac{\LH}2\Tr C_\star\leq\frac{\LH N_\theta}{2\alpha}.
\end{equation}

\emph{Mean.} Since $C_0^{-1}=H_{\mathrm{in}}$,
\[
    m_0-m_\star=d-H_{\mathrm{in}}^{-1}\bigl(\nabla V(m_\star)+\bar H d\bigr)
    =H_{\mathrm{in}}^{-1}(H_{\mathrm{in}}-\bar H)d-H_{\mathrm{in}}^{-1}\nabla V(m_\star).
\]
Hessian-Lipschitzness gives $\norm{H_{\mathrm{in}}-\bar H}\leq\int_0^1\LH(1-t)r_{\mathrm{in}}\dd t=\tfrac{\LH}2r_{\mathrm{in}}$,
and $\norm{H_{\mathrm{in}}^{-1}}\leq1/\alpha$, so with \eqref{eq:gradV-mstar},
$\norm{m_0-m_\star}\leq\tfrac1\alpha\cdot\tfrac{\LH}2r_{\mathrm{in}}\cdot r_{\mathrm{in}}+\tfrac1\alpha\cdot\tfrac{\LH N_\theta}{2\alpha}=\tfrac{\LH}{2\alpha}(r_{\mathrm{in}}^2+\tfrac{N_\theta}\alpha)$.

\emph{Covariance.} With $C_\star^{-1}=A(a_\star)=\E_{\rho_\star}[\nabla^2V]$,
\[
    \norm{H_{\mathrm{in}}-C_\star^{-1}}\leq\E_{\rho_\star}\norm{\nabla^2V(m_{\mathrm{in}})-\nabla^2V(X)}
    \leq\LH\,\E_{\rho_\star}\norm{m_{\mathrm{in}}-X}\leq\LH\bigl(r_{\mathrm{in}}+\sqrt{N_\theta/\alpha}\bigr),
\]
using $\E_{\rho_\star}\norm{m_{\mathrm{in}}-X}\leq r_{\mathrm{in}}+\sqrt{\Tr C_\star}$ and
$\Tr C_\star\leq N_\theta/\alpha$. From $C_0-C_\star=C_0(C_\star^{-1}-C_0^{-1})C_\star=H_{\mathrm{in}}^{-1}(C_\star^{-1}-H_{\mathrm{in}})C_\star$
and $\norm{X}_{\Frob}\leq\sqrt{N_\theta}\norm{X}_{\op}$,
$\norm{C_0-C_\star}_{\Frob}\leq\norm{H_{\mathrm{in}}^{-1}}\norm{C_\star}\sqrt{N_\theta}\norm{H_{\mathrm{in}}-C_\star^{-1}}\leq\tfrac{\sqrt{N_\theta}\LH}{\alpha^2}(r_{\mathrm{in}}+\sqrt{N_\theta/\alpha})$.

\emph{Gap (second order in the square-root parameter).} Since $a_\star$ is a stationary point, the
gap must be second order in the parameter error; a first-order segment estimate would be structurally
loose. Parameterize $f(m,B):=\E_Z[V(m+BZ)]-\log\det B$ over $\R^{N_\theta}\times\SPD(N_\theta)$, so
that $f(m,B)=\calE(m,BB^\top)+\text{const}$ and $\Delta(a_0)=f(m_0,B_0)-f(m_\star,B_\star)$ with
$B_0:=C_0^{1/2}$, $B_\star:=C_\star^{1/2}$. Along the affine path $m_s=m_\star+s\,\delta m$,
$B_s=B_\star+s\,\delta B$ ($\delta m=m_0-m_\star$, $\delta B=B_0-B_\star$ symmetric), the first
derivative vanishes at $s=0$: $\partial_mf=\E[\nabla V]=0$ at $a_\star$, and by Stein's identity
$\tfrac{\dd}{\dd s}\E[V(m_\star+B_sZ)]\big|_0=\Tr\bigl(\delta B\,\E[\nabla^2V(X)]B_\star\bigr)
=\Tr(\delta B\,C_\star^{-1}B_\star)=\Tr(\delta B\,B_\star^{-1})
=\tfrac{\dd}{\dd s}\log\det B_s\big|_0$. For the second derivatives,
\[
    \frac{\dd^2}{\dd s^2}\E[V(m_s+B_sZ)]
    =\E\bigl[(\delta m+\delta BZ)^\top\nabla^2V(X_s)(\delta m+\delta BZ)\bigr]
    \leq\beta\,\E\norm{\delta m+\delta BZ}^2
    =\beta\bigl(\norm{\delta m}^2+\norm{\delta B}_{\Frob}^2\bigr)
\]
(the cross term vanishes by $\E Z=0$), and, with $M_s:=B_s^{-1}\delta B$,
$\tfrac{\dd^2}{\dd s^2}\bigl[-\log\det B_s\bigr]=\Tr(M_s^2)\leq\norm{M_s}_{\Frob}^2\leq\beta\norm{\delta B}_{\Frob}^2$,
using $\Tr(M^2)=\sum_{ij}M_{ij}M_{ji}\leq\norm M_{\Frob}^2$ (Cauchy--Schwarz) and
$\norm{B_s^{-1}}_{\op}\leq\sqrt\beta$ (both $C_0$ and $C_\star$ lie on the band
$\succeq\tfrac1\beta\Id$, so $B_s\succeq\beta^{-1/2}\Id$ along the path). Taylor with vanishing first
derivative gives
\[
    \Delta(a_0)\leq\frac\beta2\norm{m_0-m_\star}^2+\beta\norm{C_0^{1/2}-C_\star^{1/2}}_{\Frob}^2 .
\]
Finally the square root is Frobenius-Lipschitz on the band: eigenvalue divided differences of
$\sqrt{\cdot}$ on $[\tfrac1\beta,\infty)$ are at most $\tfrac{\sqrt\beta}2$, so
$\norm{C_0^{1/2}-C_\star^{1/2}}_{\Frob}\leq\tfrac{\sqrt\beta}2\norm{C_0-C_\star}_{\Frob}$. Substituting
the two displayed parameter bounds gives the stated estimate; every term carries $\LH^2$, so
$\Delta(a_0)=0$ when $\LH=0$ (quadratic targets, recovering \Cref{thm:gauss-recovery}). The
first-order variant in the statement is the previous covariance-segment argument
($\norm{A(m_0,C_s)-C_s^{-1}}_{\op}\leq\beta$ along $C_s=(1-s)C_\star+sC_0$).
\end{proof}
\subsection{Pathwise covariance bands}
\label{subsec:bands}

\begin{lemma}[Covariance bands for the rescued schemes]
\label{lem:stoch-band}
Assume \Cref{assump:logconcave-smooth} and initialize with \eqref{eq:rescue}, so
$\tfrac1\beta\Id\preceq C_0\preceq\tfrac1\alpha\Id$.
\begin{enumerate}
\item \textup{(Riemannian.)} If $0<\Delta t\leq1/\kappa$, then almost surely
$\tfrac1{2\beta}\Id\preceq C_n\preceq\tfrac1\alpha\Id$ for all $n\geq0$.
\item \textup{(KL/Bregman.)} For every $\Delta t>0$, almost surely
$\tfrac1\beta\Id\preceq C_n\preceq\tfrac1\alpha\Id$ for all $n\geq0$.
\end{enumerate}
No stopping time, clipping, or containment hypothesis is used.
\end{lemma}

\begin{proof}
By \eqref{eq:pathwise-Hess}, $\alpha\Id\preceq\overline H_n\preceq\beta\Id$ a.s.

\emph{KL.} If $\tfrac1\beta\Id\preceq C_n\preceq\tfrac1\alpha\Id$ then $\alpha\Id\preceq C_n^{-1}\preceq\beta\Id$,
so $(1+\Delta t)\alpha\Id\preceq C_n^{-1}+\Delta t\overline H_n\preceq(1+\Delta t)\beta\Id$; inverting
and multiplying by $(1+\Delta t)$ gives $\tfrac1\beta\Id\preceq C_{n+1}\preceq\tfrac1\alpha\Id$. No
stepsize restriction is used (operator antitonicity of the inverse).

\emph{Riemannian.} \emph{Upper:} $C_{n+1}^{-1}\succeq(1-\Delta t)C_n^{-1}+\Delta t\overline H_n\succeq\alpha\Id$
when $C_n\preceq\tfrac1\alpha\Id$ and $\Delta t\leq1$, so $C_{n+1}\preceq\tfrac1\alpha\Id$.
\emph{Lower envelope:} with upper bound $\tfrac1\alpha\Id$, $\overline B_n:=C_n^{1/2}\overline H_nC_n^{1/2}\preceq\tfrac\beta\alpha\Id=\kappa\Id$,
so $\Delta t\overline B_n\preceq\Id$ needs $\Delta t\leq1/\kappa$; the lower envelope of
\Cref{lem:Riem-cov}, started at $L_0=\min\{\lambda_0,\tfrac1{2\beta}\}=\tfrac1{2\beta}$ (since the
rescue gives $\lambda_0\geq\tfrac1\beta>\tfrac1{2\beta}$), has $y_0=(\beta L_0)^{-1}=2$; since
$y_\infty<2$ and the recursion $y_{n+1}=e^{-\Delta t}(y_n+(e-1)\Delta t)$ moves $y_n$ monotonically
toward $y_\infty$, we have $y_n\leq2$, i.e.\ $L_n\geq\tfrac1{2\beta}$. The pathwise sampled-Hessian
bounds \eqref{eq:pathwise-Hess} let the deterministic envelope apply sample by sample.
\end{proof}

\begin{algorithm}[t]
\caption{Quadratic-rescued minibatch Price/Hessian STL Fisher--Rao scheme}
\label{alg:unified}
\begin{algorithmic}[1]
\Require initial mean $m_{\mathrm{in}}$; scheme selector $\mathsf S\in\{\Riem,\KL\}$; batch size $B$;
stepsizes $\Delta t_n>0$; iteration count $N$
\State \textbf{Phase 0 (quadratic rescue; one query).} Query the oracle at $m_{\mathrm{in}}$:
$G_{\mathrm{in}}\gets\nabla V(m_{\mathrm{in}})$, $H_{\mathrm{in}}\gets\nabla^2V(m_{\mathrm{in}})$;
set $m_0\gets m_{\mathrm{in}}-H_{\mathrm{in}}^{-1}G_{\mathrm{in}}$, $C_0\gets H_{\mathrm{in}}^{-1}$
\State \textbf{Phase I (Price/Hessian STL iterations).}
\For{$n=0,\dots,N-1$}
    \State Draw $\xi_{n,s}\overset{\text{iid}}{\sim}\N(0,\Id)$, $X_{n,s}\gets m_n+C_n^{1/2}\xi_{n,s}$, $s=1,\dots,B$; query the oracle at each $X_{n,s}$
    \State $\overline b_n\gets\frac1B\sum_s(-\nabla V(X_{n,s})+C_n^{-1/2}\xi_{n,s})$ \Comment{STL score}
    \State $\overline H_n\gets\frac1B\sum_s\nabla^2V(X_{n,s})$ \Comment{Price/Hessian covariance}
    \State $m_{n+1}\gets m_n+\Delta t_n\,C_n\overline b_n$
    \If{$\mathsf S=\Riem$}
        \State $C_{n+1}\gets C_n^{1/2}\exp\bigl(\Delta t_n(\Id-C_n^{1/2}\overline H_nC_n^{1/2})\bigr)C_n^{1/2}$
    \Else
        \State $C_{n+1}\gets(1+\Delta t_n)\bigl(C_n^{-1}+\Delta t_n\overline H_n\bigr)^{-1}$
    \EndIf
\EndFor
\end{algorithmic}
\end{algorithm}

\section{Intrinsic sticking-the-landing variance}
\label{sec:intrinsic-STL}

For a single sample $X\sim\rho_a$, the single-sample tangent noise is
$\xi(a;X):=(-C\varepsilon_b,\,-C\varepsilon_RC)$ with $\varepsilon_b:=\widehat b-g(a)$,
$\varepsilon_R:=-(\nabla^2V(X)-A(a))$, and the minibatch noise is
$\overline\xi_n:=\widehat\grad\calE(a_n)-\grad\calE(a_n)=\tfrac1B\sum_s\xi(a_n;X_{n,s})$.

\begin{definition}[Intrinsic Hessian fluctuation]
\label{def:Psi}
$\Psi(a):=\E_{X\sim\rho_a}\bigl[\norm{C^{1/2}(\nabla^2V(X)-A(a))C^{1/2}}_{\Frob}^2\bigr]$.
\end{definition}

\begin{lemma}[Intrinsic STL variance bound]
\label{lem:intrinsic-STL}
Under \Cref{assump:logconcave-smooth}, the STL noises are conditionally centered, and, with
$\norm{\xi(a;X)}_{a,\FR}^2=\varepsilon_b^\top C\varepsilon_b+\tfrac12\norm{C^{1/2}\varepsilon_RC^{1/2}}_{\Frob}^2$,
\begin{equation}
\label{eq:intrinsic-mean}
    \E\bigl[\varepsilon_b^\top C\varepsilon_b\bigr]\leq\norm{C^{1/2}R(a)C^{1/2}}_{\Frob}^2+\Psi(a),
    \qquad
    \E\bigl[\norm{\xi(a;X)}_{a,\FR}^2\bigr]\leq2\,\Gradnorm(a)^2+\frac32\,\Psi(a).
\end{equation}
Consequently, for every $B\geq1$,
$\E[\norm{\overline\xi_n}_{a_n,\FR}^2\mid\Fil_n]\leq B^{-1}(2\Gradnorm(a_n)^2+\tfrac32\Psi(a_n))$.
\end{lemma}

\begin{proof}
$\E[\widehat b]=g(a)$ and $\E[\widehat K]=R(a)$ give centering; $\E\norm{C^{1/2}\varepsilon_RC^{1/2}}_{\Frob}^2=\Psi(a)$.
For the mean, let $F(X):=-\nabla V(X)+C^{-1}(X-m)$, so $\varepsilon_b=F(X)-\E F(X)$ and $\nabla F(X)=C^{-1}-\nabla^2V(X)$
is symmetric. The Gaussian Poincar\'e inequality applied to $\phi_i(X):=(C^{1/2}F(X))_i$ and summed gives
$\E[\varepsilon_b^\top C\varepsilon_b]\leq\E\norm{C^{1/2}\nabla F(X)C^{1/2}}_{\Frob}^2$; expanding
$C^{1/2}\nabla F(X)C^{1/2}=C^{1/2}R(a)C^{1/2}-C^{1/2}(\nabla^2V(X)-A(a))C^{1/2}$ (the cross term
vanishes in expectation) gives \eqref{eq:intrinsic-mean}(left). Then
$\E\norm\xi_{a,\FR}^2\leq\norm{C^{1/2}RC^{1/2}}_{\Frob}^2+\tfrac32\Psi\leq2\Gradnorm^2+\tfrac32\Psi$,
using $\norm{C^{1/2}RC^{1/2}}_{\Frob}^2\leq2\Gradnorm^2$. The minibatch bound follows by independence.
\end{proof}

\begin{remark}[Regularity]
\label{rem:regularity}
The differentiations above transfer to Gaussian score polynomials for observables of at most
polynomial growth, which holds under \Cref{assump:logconcave-smooth} since $\nabla V$ grows at most
linearly; no third or fourth derivatives of $V$ are used. When $\nabla^2V$ is only Lipschitz it is
differentiable a.e., and the Gaussian Poincar\'e step is applied to the Lipschitz map
$F(X)=-\nabla V(X)+C^{-1}(X-m)$ using its weak derivative $\nabla F=C^{-1}-\nabla^2V$ (equivalently by
mollification), which is all the inequality requires.
\end{remark}

\subsection{Assumption-minimal regime: the self-bounding estimate}
\label{sec:stoch-selfbound}

\begin{lemma}[Self-bounding estimate for the intrinsic fluctuation]
\label{lem:Psi-self-bound}
Under \Cref{assump:logconcave-smooth}, if $C\preceq\Lambda\Id$ and $\LB=\beta\Lambda$, then
\begin{equation}
\label{eq:Psi-self-bound}
    \boxed{\;
    \Psi(a)\leq2\LB N_\theta+\LB\norm{C^{1/2}R(a)C^{1/2}}_{\Frob}^2
    \leq2\LB N_\theta+2\LB\,\Gradnorm(a)^2 .
    \;}
\end{equation}
\end{lemma}

\begin{proof}
Write $B_X:=C^{1/2}\nabla^2V(X)C^{1/2}$, $S_A:=C^{1/2}A(a)C^{1/2}=\E[B_X]$. \emph{Step 1:}
$\Psi(a)=\E\Tr(B_X^2)-\Tr(S_A^2)\leq\E\Tr(B_X^2)$. \emph{Step 2:} $0\preceq B_X\preceq\beta C\preceq\LB\Id$,
and for $0\preceq M\preceq L\Id$, $M^2\preceq LM$; so $B_X^2\preceq\LB B_X$ a.s.\ and
$\E\Tr(B_X^2)\leq\LB\Tr(S_A)$. \emph{Step 3:} $s\leq2+(s-1)^2$ for all real $s$, so
$\Tr(S_A)\leq2N_\theta+\norm{S_A-\Id}_{\Frob}^2$; since $S_A-\Id=-C^{1/2}R(a)C^{1/2}$, combining gives
the first inequality; the second uses $\norm{C^{1/2}RC^{1/2}}_{\Frob}^2\leq2\Gradnorm^2$.
\end{proof}

\begin{remark}[Sharper self-bounding constants]
\label{rem:selfbound-sharp}
Retaining the centering term $-\Tr(S_A^2)$ discarded in Step~1 improves \eqref{eq:Psi-self-bound}:
with $s_i$ the eigenvalues of $S_A$, $\Psi\leq\sum_is_i(\LB-s_i)$, and the scalar bound
$s(\LB-s)\leq\tfrac{\LB(5\LB-4)}{4(\LB+1)}+\LB(s-1)^2$ (a one-dimensional maximization) gives
$\Psi\leq\tfrac54\LB N_\theta+\LB\norm{S_A-\Id}_{\Frob}^2$, reducing the dimension coefficient from
$2\LB$ to at most $\tfrac54\LB$. On the rescued band one also has the target-conditioned bound
$\Psi\leq N_\theta(\kappa-1)^2$ (from $\norm{\nabla^2V(X)-A(a)}_\op\leq\beta-\alpha$ and
$C\preceq\tfrac1\alpha\Id$), which vanishes at $\kappa=1$; the minimum of the estimates may be used.
The downstream theorems retain the conservative constants of \eqref{eq:Psi-self-bound}, since only
their orders enter the conclusions.
\end{remark}

\begin{corollary}[Self-bounded second moment]
\label{cor:second-moment}
If $C\preceq\Lambda\Id$, then for every $B\geq1$,
\begin{equation}
\label{eq:second-moment}
    \E\bigl[\norm{\widehat\grad\calE(a)}_{a,\FR}^2\mid a\bigr]
    \leq\rho_B\,\Gradnorm(a)^2+\frac{3\LB N_\theta}B,
    \qquad
    \rho_B:=1+\frac{2+3\LB}B .
\end{equation}
\end{corollary}

\begin{proof}
$\E\norm{\widehat\grad\calE}^2=\Gradnorm^2+\E\norm{\overline\xi}^2$; by \Cref{lem:intrinsic-STL} and
\Cref{lem:Psi-self-bound},
$\E\norm{\overline\xi}^2\leq B^{-1}(2\Gradnorm^2+\tfrac32(2\LB N_\theta+2\LB\Gradnorm^2))=\tfrac{(2+3\LB)\Gradnorm^2}B+\tfrac{3\LB N_\theta}B$.
\end{proof}

\subsection{Hessian-Lipschitz regime}
\label{sec:stoch-Hlip}

\begin{assumption}[Global Hessian-Lipschitzness]
\label{assump:global-Hess-Lip}
There is $\LH<\infty$ with $\norm{\nabla^2V(x)-\nabla^2V(y)}_\op\leq\LH\norm{x-y}$ for all $x,y$. We
write the dimensionless parameter $\barLH:=\LH/\alpha^{3/2}$; a Gaussian target has $\LH=\barLH=0$.
\end{assumption}

\begin{lemma}[Hessian-Lipschitz control of the intrinsic fluctuation]
\label{lem:Psi-Hg}
Under \Cref{assump:global-Hess-Lip}, the anisotropic (effective-dimension) bound
$\Psi(a)\leq\LH^2\,\Tr(C^2)\,\Tr C\leq\LH^2\norm{C}_\op(\Tr C)^2$ holds. On the band
$C\preceq\tfrac1\alpha\Id$,
\begin{equation}
\label{eq:Psi-Hg}
    \boxed{\;\Psi(a)\leq\Psi_H:=\frac{N_\theta^2\LH^2}{\alpha^3}=N_\theta^2\barLH^2 .\;}
\end{equation}
\end{lemma}

\begin{proof}
Write $X=m+C^{1/2}Z$, $F(Z):=C^{1/2}\nabla^2V(m+C^{1/2}Z)C^{1/2}$. Gaussian Poincar\'e (componentwise)
gives $\Psi(a)=\E\norm{F(Z)-\E F(Z)}_{\Frob}^2\leq\E\norm{D_ZF(Z)}_{\mathrm{HS}}^2$. For a unit vector
$u$, $D_ZF(Z)[u]=C^{1/2}M_uC^{1/2}$ with $M_u:=D(\nabla^2V)(X)[C^{1/2}u]$ symmetric and
$\norm{M_u}_\op\leq\LH\norm{C^{1/2}u}=\LH(u^\top Cu)^{1/2}$. In the eigenbasis of $C$ (eigenvalues
$c_i$), Cauchy--Schwarz on the weights gives
\[
    \norm{C^{1/2}M_uC^{1/2}}_{\Frob}^2=\sum_{ij}c_ic_j(M_u)_{ij}^2
    \leq\sum_{ij}c_i^2(M_u)_{ij}^2=\Tr(C^2M_u^2)\leq\norm{M_u}_\op^2\,\Tr(C^2),
\]
so $\norm{D_ZF(Z)[u]}_{\Frob}^2\leq\LH^2\,(u^\top Cu)\,\Tr(C^2)$; summing over an orthonormal basis
gives $\norm{D_ZF(Z)}_{\mathrm{HS}}^2\leq\LH^2\,\Tr C\,\Tr(C^2)$, and $\Tr(C^2)\leq\norm{C}_\op\Tr C$
recovers the coarser form. On $C\preceq\tfrac1\alpha\Id$, $\Tr(C^2)\leq N_\theta/\alpha^2$,
$\Tr C\leq N_\theta/\alpha$. The worst case over the band remains $N_\theta^2$; the anisotropic form
is smaller whenever the covariance spectrum decays (effective dimension
$\Tr(C^2)\Tr C/\norm C_\op^3$ in place of $N_\theta^2$).
\end{proof}
\section{Stochastic global convergence: Riemannian variant}
\label{sec:stoch-Riem}

On the rescued bands, \Cref{lem:FR-W2} gives the Fisher--Rao PL inequalities
\begin{equation}
\label{eq:PL-bands}
    \Gradnorm(a_n)^2\geq\frac1{2\kappa}\Delta_n\ \ (\text{Riemannian band }C_n\succeq\tfrac1{2\beta}\Id),
    \qquad
    \Gradnorm(a_n)^2\geq\frac1\kappa\Delta_n\ \ (\text{KL band }C_n\succeq\tfrac1\beta\Id).
\end{equation}

\begin{lemma}[Riemannian stochastic one-step descent]
\label{lem:Riem-onestep-stoch}
Assume \Cref{assump:logconcave-smooth} and $\tfrac1{2\beta}\Id\preceq C_n\preceq\tfrac1\alpha\Id$. For
one step of \eqref{eq:Riem-STL} with $0<\Delta t\leq1/\kappa$,
\begin{equation}
\label{eq:Riem-onestep-Psi}
    \E[\Delta_{n+1}\mid\Fil_n]
    \leq\Delta_n-\Delta t\Bigl(1-\kappa\Delta t\bigl(1+\tfrac2B\bigr)\Bigr)\Gradnorm(a_n)^2
    +\frac{3\kappa\Delta t^2}{2B}\Psi(a_n).
\end{equation}
\end{lemma}

\begin{proof}
By \Cref{lem:stoch-band}(1), $C_n\preceq\tfrac1\alpha\Id$, so the scope condition \eqref{eq:scope}
holds pathwise for $\zeta_n=-\Delta t\widehat\grad\calE(a_n)$ (as in \Cref{lem:retr-smooth}, since
$\Delta t\kappa\leq1$), and \Cref{lem:retr-smooth} with $L=2\LB=2\kappa$ gives pathwise
\[
    \calE(a_{n+1})\leq\calE(a_n)+\inner{\grad\calE(a_n)}{\zeta_n}_{a_n}+\kappa\Delta t^2\norm{\widehat\grad\calE(a_n)}_{a_n}^2 .
\]
Taking $\E[\cdot\mid\Fil_n]$, \eqref{eq:STL-unbiased} gives
$\E[\inner{\grad\calE(a_n)}{\widehat\grad\calE(a_n)}_{a_n}\mid\Fil_n]=\Gradnorm(a_n)^2$, and
\Cref{lem:intrinsic-STL} gives
$\E[\norm{\widehat\grad\calE(a_n)}_{a_n}^2\mid\Fil_n]\leq(1+\tfrac2B)\Gradnorm(a_n)^2+\tfrac{3}{2B}\Psi(a_n)$.
Collecting terms gives \eqref{eq:Riem-onestep-Psi}.
\end{proof}

\begin{theorem}[Riemannian, Hessian-Lipschitz]
\label{thm:Riem-Hlip}
Assume \Cref{assump:logconcave-smooth,assump:global-Hess-Lip}, initialize with \eqref{eq:rescue}, and
run \eqref{eq:Riem-STL} with $B\geq1$ and $\Delta t=\tfrac1{2\kappa(1+2/B)}$. Then for every $N\geq0$,
\begin{equation}
\label{eq:Riem-Hlip}
    \boxed{\;
    \E\Delta_N\leq\Bigl(1-\frac1{8\kappa^2(1+2/B)}\Bigr)^N\Delta(a_0)+\frac{3\kappa}{B+2}\,N_\theta^2\barLH^2 .
    \;}
\end{equation}
For $B=1$: $\Delta t=\tfrac1{6\kappa}$, $\E\Delta_N\leq(1-\tfrac1{24\kappa^2})^N\Delta(a_0)+\kappa N_\theta^2\barLH^2$,
so the single-sample iteration and oracle complexities are $N_{\Riem}=\widetilde O(\kappa^2)$ and
$Q_{\Riem}=1+\widetilde O(\kappa^2)$.
\end{theorem}

\begin{proof}
\Cref{lem:stoch-band}(1) gives the band pathwise. At $\Delta t=\tfrac1{2\kappa(1+2/B)}$ the bracket in
\eqref{eq:Riem-onestep-Psi} equals $\tfrac12$, so by \Cref{lem:Psi-Hg} ($\Psi\leq\Psi_H=N_\theta^2\barLH^2$)
and the PL bound \eqref{eq:PL-bands},
$\E[\Delta_{n+1}\mid\Fil_n]\leq(1-\tfrac{\Delta t}{4\kappa})\Delta_n+\tfrac{3\kappa\Delta t^2}{2B}\Psi_H$.
Unrolling, the floor is $\tfrac{3\kappa\Delta t^2\Psi_H/(2B)}{\Delta t/(4\kappa)}=\tfrac{6\kappa^2\Delta t}B\Psi_H=\tfrac{3\kappa}{B+2}\Psi_H$,
and the contraction is $1-\tfrac{\Delta t}{4\kappa}=1-\tfrac1{8\kappa^2(1+2/B)}$.
\end{proof}

\begin{theorem}[Riemannian, assumption-minimal self-bounding]
\label{thm:Riem-selfbound}
Assume \Cref{assump:logconcave-smooth}, initialize with \eqref{eq:rescue}, and run \eqref{eq:Riem-STL}
with $B\geq1$ and $\Delta t=\tfrac1{2\kappa\rho_B}$, $\rho_B:=1+\tfrac{2+3\kappa}B$. Then for every
$N\geq0$,
\begin{equation}
\label{eq:Riem-selfbound}
    \boxed{\;
    \E\Delta_N\leq\Bigl(1-\frac1{8\kappa^2\rho_B}\Bigr)^N\Delta(a_0)+\frac{6\kappa^2N_\theta}{B+2+3\kappa} .
    \;}
\end{equation}
The floor is $O(\kappa N_\theta)$ for every $B\lesssim\kappa$ and at most $2\kappa N_\theta$ for every
$B\geq1$. For $B=1$, $\rho_1=3+3\kappa\asymp\kappa$, giving $N_{\Riem}=\widetilde O(\kappa^3)$; for
$B=\Theta(\kappa)$, $\rho_B=O(1)$, giving $N_{\Riem}=\widetilde O(\kappa^2)$ rounds and
$Q_{\Riem}=BN_{\Riem}=\widetilde O(\kappa^3)$ oracle pairs. No Hessian-Lipschitz assumption is used.
\end{theorem}

\begin{proof}
On the band $C_n\preceq\tfrac1\alpha\Id$, $\LB=\kappa$, so \Cref{lem:Psi-self-bound} gives
$\tfrac32\Psi(a_n)\leq3\kappa N_\theta+3\kappa\Gradnorm(a_n)^2$. Substituting into
\eqref{eq:Riem-onestep-Psi} and collecting the $\Gradnorm^2$ terms,
\[
    \E[\Delta_{n+1}\mid\Fil_n]
    \leq\Delta_n-\Delta t\bigl(1-\kappa\Delta t\rho_B\bigr)\Gradnorm(a_n)^2+\frac{3\kappa^2\Delta t^2N_\theta}B,
    \qquad
    \rho_B=1+\frac{2+3\kappa}B,
\]
since $\kappa\Delta t^2(1+\tfrac2B)+\tfrac{3\kappa\Delta t^2}{2B}\cdot2\kappa=\kappa\Delta t^2(1+\tfrac{2+3\kappa}B)=\kappa\Delta t^2\rho_B$.
At $\Delta t=\tfrac1{2\kappa\rho_B}$ the bracket is $\tfrac12$, and PL \eqref{eq:PL-bands} gives
$\E[\Delta_{n+1}\mid\Fil_n]\leq(1-\tfrac{\Delta t}{4\kappa})\Delta_n+\tfrac{3\kappa^2\Delta t^2N_\theta}B$.
The floor is $\tfrac{3\kappa^2\Delta t^2N_\theta/B}{\Delta t/(4\kappa)}=\tfrac{12\kappa^3\Delta tN_\theta}B=\tfrac{6\kappa^2N_\theta}{B\rho_B}=\tfrac{6\kappa^2N_\theta}{B+2+3\kappa}$,
and the contraction is $1-\tfrac{\Delta t}{4\kappa}=1-\tfrac1{8\kappa^2\rho_B}$. For $B+2+3\kappa\geq3\kappa$,
the floor is $\leq2\kappa N_\theta$.
\end{proof}
\section{Stochastic global convergence: KL/Bregman variant}
\label{sec:stoch-KL}

The KL/Bregman scheme \eqref{eq:KL-STL} enjoys the tighter pathwise band of
\Cref{lem:stoch-band}(2), $\tfrac1\beta\Id\preceq C_n\preceq\tfrac1\alpha\Id$ for \emph{every}
$\Delta t>0$, and the corresponding Fisher--Rao PL inequality from \eqref{eq:PL-bands},
$\Gradnorm(a_n)^2\geq\tfrac1\kappa\Delta_n$. The key observation is that the implicit resolvent
covariance update is exactly an affine-invariant exponential retraction, so the one-step estimate
follows from the arbitrary-direction retraction smoothness of \Cref{lem:retr-smooth} without any
principal-square-root coupling.

\begin{lemma}[Logarithmic tangent of the KL resolvent update]
\label{lem:KL-logtangent}
Fix $a=(m,C)$ with $\tfrac1\beta\Id\preceq C\preceq\tfrac1\alpha\Id$ and $\alpha\Id\preceq\overline H\preceq\beta\Id$,
and set $P:=C^{1/2}\overline HC^{1/2}$, so that $\tfrac1\kappa\Id\preceq P\preceq\kappa\Id$. Write
$\overline S:=\Id-P$ and define the logarithmic tangent
\begin{equation}
\label{eq:Wlog}
    W(P):=\log(1+\Delta t)\,\Id-\log(\Id+\Delta t\,P).
\end{equation}
Then the KL covariance update $C^+=(1+\Delta t)(C^{-1}+\Delta t\,\overline H)^{-1}$ satisfies
\begin{equation}
\label{eq:KL-retr}
    C^+=C^{1/2}e^{W(P)}C^{1/2},
    \qquad\text{equivalently}\qquad
    a^+=\Retr_a\bigl(\Delta t\,C\overline b,\;C^{1/2}W(P)C^{1/2}\bigr),
\end{equation}
where $\overline b$ is the minibatch STL score. Moreover, with $\mathcal R(P):=W(P)-\Delta t\,\overline S$,
\begin{equation}
\label{eq:Wlog-bounds}
    \norm{W(P)}_\op\leq\Delta t\,\kappa,
    \qquad
    \norm{W(P)}_{\Frob}\leq\Delta t\,\norm{\overline S}_{\Frob},
    \qquad
    \norm{\mathcal R(P)}_{\Frob}\leq\kappa\,\Delta t^2\,\norm{\overline S}_{\Frob}.
\end{equation}
In particular, for $\Delta t\leq1/(8\kappa)$, $\norm{W(P)}_\op\leq\tfrac18<1$, so the tangent
$\zeta:=(\Delta t\,C\overline b,\,C^{1/2}W(P)C^{1/2})$ satisfies the scope condition \eqref{eq:scope}.
\end{lemma}

\begin{proof}
Since $C^{-1}+\Delta t\,\overline H=C^{-1/2}(\Id+\Delta t\,P)C^{-1/2}$, we have
$(C^{-1}+\Delta t\,\overline H)^{-1}=C^{1/2}(\Id+\Delta t\,P)^{-1}C^{1/2}$ and hence
$C^+=C^{1/2}[(1+\Delta t)(\Id+\Delta t\,P)^{-1}]C^{1/2}$. As $\log(1+\Delta t)\Id$ and $\log(\Id+\Delta t\,P)$
are commuting functions of $P$, $e^{W(P)}=(1+\Delta t)(\Id+\Delta t\,P)^{-1}$, giving the first identity
in \eqref{eq:KL-retr}; the retraction form follows from
$\Retr_a(u,U)=(m+u,\,C^{1/2}\exp(C^{-1/2}UC^{-1/2})C^{1/2})$ with $U=C^{1/2}W(P)C^{1/2}$, so that
$C^{-1/2}UC^{-1/2}=W(P)$. The band $\tfrac1\kappa\Id\preceq P\preceq\kappa\Id$ is
$\alpha C\preceq P\preceq\beta C$ combined with $\tfrac1\beta\Id\preceq C\preceq\tfrac1\alpha\Id$.
For the scalar bounds, each eigenvalue $p\in[\tfrac1\kappa,\kappa]$ of $P$ contributes
$w(p):=\log(1+\Delta t)-\log(1+\Delta t\,p)=\int_p^1\frac{\Delta t}{1+\Delta t\,r}\dd r$. Since
$0\leq\tfrac{\Delta t}{1+\Delta t\,r}\leq\Delta t$ for $r\geq0$, $|w(p)|\leq\Delta t\,|1-p|$, which gives both
$\norm{W(P)}_\op\leq\Delta t\max_p|1-p|\leq\Delta t\,\kappa$ and, summing over eigenvalues,
$\norm{W(P)}_{\Frob}\leq\Delta t\,\norm{\Id-P}_{\Frob}=\Delta t\,\norm{\overline S}_{\Frob}$. For the
remainder, $w(p)-\Delta t(1-p)=\int_p^1(\tfrac{\Delta t}{1+\Delta t\,r}-\Delta t)\dd r=-\Delta t^2\int_p^1\tfrac{r}{1+\Delta t\,r}\dd r$,
and $0\leq\tfrac{r}{1+\Delta t\,r}\leq r\leq\kappa$ on the integration range gives
$|w(p)-\Delta t(1-p)|\leq\kappa\,\Delta t^2|1-p|$; summing over eigenvalues yields
$\norm{\mathcal R(P)}_{\Frob}\leq\kappa\,\Delta t^2\norm{\Id-P}_{\Frob}=\kappa\,\Delta t^2\norm{\overline S}_{\Frob}$.
\end{proof}

\begin{lemma}[KL stochastic one-step descent]
\label{lem:KL-onestep-stoch}
Assume \Cref{assump:logconcave-smooth} and $\tfrac1\beta\Id\preceq C_n\preceq\tfrac1\alpha\Id$. For
one step of \eqref{eq:KL-STL} with $0<\Delta t\leq1/(8\kappa)$,
\begin{equation}
\label{eq:KL-onestep-Psi}
    \E[\Delta_{n+1}\mid\Fil_n]
    \leq\Delta_n-\frac{\Delta t}2\,\Gradnorm(a_n)^2
    +\frac{7\kappa\,\Delta t^2}{4B}\,\Psi(a_n).
\end{equation}
\end{lemma}

\begin{proof}
Suppress $n$ and condition on $\Fil_n$. Write $A:=A(a)$, $R:=C^{-1}-A$, $S:=C^{1/2}RC^{1/2}$,
$P_0:=C^{1/2}AC^{1/2}=\Id-S$, and $D:=P-P_0$, so $\overline S=S-D$ with $\E[D\mid\Fil]=0$ and
$\E[\norm D_{\Frob}^2\mid\Fil]=\Psi(a)/B$ by \Cref{def:Psi} and minibatch independence. Set
\[
    M:=g(a)^\top C\,g(a),\qquad V:=\norm S_{\Frob}^2,\qquad\sigma^2:=\frac{\Psi(a)}B,
    \qquad\text{so}\qquad\Gradnorm(a)^2=M+\tfrac12V .
\]

\emph{Step 1 (retraction smoothness).} On the band $C\preceq\tfrac1\alpha\Id$, \Cref{lem:retr-smooth}
applies with $\Lambda=\tfrac1\alpha$, $L=2\beta\Lambda=2\kappa$; by \Cref{lem:KL-logtangent} the tangent
$\zeta=(\Delta t\,C\overline b,\,C^{1/2}W(P)C^{1/2})$ meets the scope condition, so pathwise
\[
    \calE(a^+)\leq\calE(a)+\inner{\grad\calE(a)}\zeta_a+\kappa\norm\zeta_a^2 .
\]
The first-order term evaluates to
$\inner{\grad\calE(a)}\zeta_a=-\Delta t\,g^\top C\overline b-\tfrac12\Tr(S\,W(P))$: the mean block uses
$\grad_m\calE=-Cg$ in the metric $v^\top C^{-1}v$, giving $-\Delta t\,g^\top C\overline b$; the covariance
block uses $\grad_C\calE=-CRC$ in the metric $\tfrac12\Tr(C^{-1}\cdot C^{-1}\cdot)$, giving
$\tfrac12\Tr(C^{-1}(-CRC)C^{-1}\,C^{1/2}W(P)C^{1/2})=-\tfrac12\Tr(S\,W(P))$. Likewise
$\norm\zeta_a^2=\Delta t^2\,\overline b^\top C\overline b+\tfrac12\norm{W(P)}_{\Frob}^2$. Hence, pathwise,
\begin{equation}
\label{eq:KL-pathwise}
    \calE(a^+)-\calE(a)
    \leq-\Delta t\,g^\top C\overline b-\tfrac12\Tr(S\,W(P))
    +\kappa\,\Delta t^2\,\overline b^\top C\overline b+\tfrac\kappa2\norm{W(P)}_{\Frob}^2 .
\end{equation}

\emph{Step 2 (mean terms).} Write $\overline b=g+\varepsilon_b$ with $\E[\varepsilon_b\mid\Fil]=0$.
By \eqref{eq:intrinsic-mean}(left) and independence,
$\E[\varepsilon_b^\top C\varepsilon_b\mid\Fil]\leq(V+\Psi(a))/B=\tfrac VB+\sigma^2$, so
$\E[\overline b^\top C\overline b\mid\Fil]\leq M+\tfrac VB+\sigma^2$. Therefore
\begin{equation}
\label{eq:KL-mean}
    \E\bigl[-\Delta t\,g^\top C\overline b+\kappa\,\Delta t^2\,\overline b^\top C\overline b\mid\Fil\bigr]
    \leq-\Delta t(1-\kappa\Delta t)M+\frac{\kappa\,\Delta t^2}BV+\kappa\,\Delta t^2\sigma^2 .
\end{equation}

\emph{Step 3 (first-order covariance drift).} With $W(P)=\Delta t\,\overline S+\mathcal R(P)$,
$-\tfrac12\Tr(S\,W(P))=-\tfrac{\Delta t}2\Tr(S\overline S)-\tfrac12\Tr(S\,\mathcal R(P))$. Since
$\E[\overline S\mid\Fil]=S$, $\E[\Tr(S\overline S)\mid\Fil]=\norm S_{\Frob}^2=V$. For the remainder,
Cauchy--Schwarz, \eqref{eq:Wlog-bounds}, and $\E[\norm{\overline S}_{\Frob}^2\mid\Fil]=V+\sigma^2$
(as $\overline S=S-D$) give, via Jensen and $x\sqrt{x^2+y^2}\leq x^2+\tfrac12y^2$ with $x=\sqrt V,\,y=\sigma$,
\[
    \bigl|\E[\Tr(S\,\mathcal R(P))\mid\Fil]\bigr|
    \leq\sqrt V\,\kappa\,\Delta t^2\,\E[\norm{\overline S}_{\Frob}\mid\Fil]
    \leq\kappa\,\Delta t^2\sqrt V\sqrt{V+\sigma^2}\leq\kappa\,\Delta t^2\bigl(V+\tfrac12\sigma^2\bigr).
\]
Hence
\begin{equation}
\label{eq:KL-covdrift}
    \E\bigl[-\tfrac12\Tr(S\,W(P))\mid\Fil\bigr]
    \leq-\frac{\Delta t}2V+\frac{\kappa\,\Delta t^2}2V+\frac{\kappa\,\Delta t^2}4\sigma^2 .
\end{equation}

\emph{Step 4 (quadratic covariance term).} From $\norm{W(P)}_{\Frob}\leq\Delta t\norm{\overline S}_{\Frob}$,
$\E[\norm{W(P)}_{\Frob}^2\mid\Fil]\leq\Delta t^2(V+\sigma^2)$, so
\begin{equation}
\label{eq:KL-covquad}
    \frac\kappa2\E[\norm{W(P)}_{\Frob}^2\mid\Fil]\leq\frac{\kappa\,\Delta t^2}2V+\frac{\kappa\,\Delta t^2}2\sigma^2 .
\end{equation}

\emph{Step 5 (collect).} Adding \eqref{eq:KL-mean}, \eqref{eq:KL-covdrift}, \eqref{eq:KL-covquad} to
$\Delta_n$ (using $\Delta_{n+1}-\Delta_n=\calE(a^+)-\calE(a)$),
\[
    \E[\Delta_{n+1}\mid\Fil]
    \leq\Delta_n-\Delta t(1-\kappa\Delta t)M
    -\Delta t\Bigl[\tfrac12-\kappa\Delta t\bigl(1+\tfrac1B\bigr)\Bigr]V
    +\Bigl(1+\tfrac14+\tfrac12\Bigr)\kappa\,\Delta t^2\sigma^2 ,
\]
where the $\sigma^2$ coefficient is $\tfrac74$. For $\Delta t\leq1/(8\kappa)$, $1-\kappa\Delta t\geq\tfrac78\geq\tfrac12$
and, since $B\geq1$, $\tfrac12-\kappa\Delta t(1+\tfrac1B)\geq\tfrac12-\tfrac18\cdot2=\tfrac14$. Thus
\[
    \E[\Delta_{n+1}\mid\Fil]\leq\Delta_n-\frac{\Delta t}2M-\frac{\Delta t}4V+\frac{7\kappa\,\Delta t^2}{4B}\Psi(a)
    =\Delta_n-\frac{\Delta t}2\Bigl(M+\tfrac12V\Bigr)+\frac{7\kappa\,\Delta t^2}{4B}\Psi(a),
\]
which is \eqref{eq:KL-onestep-Psi} since $\mathsf G(a)^2=M+\tfrac12V$.
\end{proof}

\begin{theorem}[KL, Hessian-Lipschitz]
\label{thm:KL-Hlip}
Assume \Cref{assump:logconcave-smooth,assump:global-Hess-Lip}, initialize with \eqref{eq:rescue},
and run \eqref{eq:KL-STL} with $B\geq1$ and $\Delta t=\tfrac1{8\kappa}$. Then for every $N\geq0$,
\begin{equation}
\label{eq:KL-Hlip}
    \boxed{\;
    \E\Delta_N\leq\Bigl(1-\frac1{16\kappa^2}\Bigr)^N\Delta(a_0)+\frac{7\kappa}{16B}\,N_\theta^2\barLH^2 .
    \;}
\end{equation}
For $B=1$ the single-sample iteration and oracle complexities are $N_{\KL}=\widetilde O(\kappa^2)$ and
$Q_{\KL}=1+\widetilde O(\kappa^2)$.
\end{theorem}

\begin{proof}
\Cref{lem:stoch-band}(2) gives the band pathwise for every $\Delta t>0$. By \Cref{lem:Psi-Hg},
$\Psi\leq\Psi_H=N_\theta^2\barLH^2$; substituting into \eqref{eq:KL-onestep-Psi} and using the KL
PL bound $\Gradnorm(a_n)^2\geq\tfrac1\kappa\Delta_n$ from \eqref{eq:PL-bands},
$\E[\Delta_{n+1}\mid\Fil_n]\leq(1-\tfrac{\Delta t}{2\kappa})\Delta_n+\tfrac{7\kappa\Delta t^2}{4B}\Psi_H$.
At $\Delta t=\tfrac1{8\kappa}$ the contraction is $1-\tfrac1{16\kappa^2}$ and the floor is
$\tfrac{7\kappa\Delta t^2\Psi_H/(4B)}{\Delta t/(2\kappa)}=\tfrac{7\kappa}{16B}\Psi_H$; unrolling gives
\eqref{eq:KL-Hlip}. Reaching $\E\Delta_N\leq\epsilon+$floor needs $N=\widetilde O(\kappa^2)$.
\end{proof}

\begin{remark}[KL batch-dependent stepsize window and rate--floor tradeoff]
\label{rem:KL-batch}
The proof of \Cref{lem:KL-onestep-stoch} holds verbatim under the weaker batch-dependent condition
$\Delta t\leq\tfrac1{4\kappa(1+1/B)}$ (then $1-\kappa\Delta t\geq\tfrac34$ and
$\tfrac12-\kappa\Delta t(1+\tfrac1B)\geq\tfrac14$; the scope bound
$\norm{W(P)}_\op\leq\kappa\Delta t\leq\tfrac14<1$ of \Cref{lem:KL-logtangent} persists), of which
$\Delta t\leq\tfrac1{8\kappa}$ is exactly the $B=1$ case. Before fixing the step, the
Hessian-Lipschitz recursion unrolls to
\[
    \E\Delta_N\leq\Bigl(1-\frac{\Delta t}{2\kappa}\Bigr)^N\Delta(a_0)+\frac{7\kappa^2\Delta t}{2B}\Psi_H,
    \qquad 0<\Delta t\leq\frac1{4\kappa(1+1/B)},
\]
making the rate--floor tradeoff explicit: for large $B$ the certified step nearly doubles, halving
the outer-round constant at the cost of a proportionally larger fixed-step floor. The analogous
self-bounding window is $\Delta t\leq\min\{\tfrac1{4\kappa(1+1/B)},\tfrac B{14\kappa^2}\}$. The
theorem statements retain the conservative $\Delta t=\tfrac1{8\kappa}$.
\end{remark}

\begin{theorem}[KL, assumption-minimal self-bounding]
\label{thm:KL-selfbound}
Assume \Cref{assump:logconcave-smooth}, initialize with \eqref{eq:rescue}, and run \eqref{eq:KL-STL}
with $B\geq1$ and $\Delta t=\min\{\tfrac1{8\kappa},\tfrac B{14\kappa^2}\}$. Then for every $N\geq0$,
\begin{equation}
\label{eq:KL-selfbound}
    \boxed{\;
    \E\Delta_N\leq\Bigl(1-\frac{\Delta t}{4\kappa}\Bigr)^N\Delta(a_0)+14\kappa^3\Delta t\,\frac{N_\theta}B .
    \;}
\end{equation}
For $B\leq\tfrac74\kappa$, $\Delta t=\tfrac B{14\kappa^2}$, the floor is exactly $\kappa N_\theta$ and
the contraction is $1-\tfrac B{56\kappa^3}$; for $B=1$ this gives $N_{\KL}=\widetilde O(\kappa^3)$, and
for $B=\Theta(\kappa)$, $N_{\KL}=\widetilde O(\kappa^2)$ rounds and $Q_{\KL}=\widetilde O(\kappa^3)$
oracle pairs. No Hessian-Lipschitz assumption is used.
\end{theorem}

\begin{proof}
On the band $C_n\preceq\tfrac1\alpha\Id$, $\LB=\kappa$, so \Cref{lem:Psi-self-bound} gives
$\Psi(a_n)\leq2\kappa N_\theta+2\kappa\Gradnorm(a_n)^2$. Substituting into \eqref{eq:KL-onestep-Psi},
\[
    \E[\Delta_{n+1}\mid\Fil_n]
    \leq\Delta_n-\frac{\Delta t}2\Bigl(1-\frac{7\kappa^2\Delta t}B\Bigr)\Gradnorm(a_n)^2+\frac{7\kappa^2\Delta t^2}{2B}N_\theta .
\]
The choice $\Delta t\leq\tfrac B{14\kappa^2}$ makes $7\kappa^2\Delta t/B\leq\tfrac12$, so the bracket is
$\geq\tfrac12$; using the KL PL bound $\Gradnorm(a_n)^2\geq\tfrac1\kappa\Delta_n$,
\[
    \E[\Delta_{n+1}\mid\Fil_n]\leq\Bigl(1-\frac{\Delta t}{4\kappa}\Bigr)\Delta_n+\frac{7\kappa^2\Delta t^2}{2B}N_\theta .
\]
Unrolling, the floor is $\tfrac{7\kappa^2\Delta t^2N_\theta/(2B)}{\Delta t/(4\kappa)}=\tfrac{14\kappa^3\Delta t N_\theta}B$.
When $\Delta t=\tfrac B{14\kappa^2}$ this equals $\kappa N_\theta$, and the contraction is
$1-\tfrac{\Delta t}{4\kappa}=1-\tfrac B{56\kappa^3}$.
\end{proof}

\begin{remark}[The two schemes side by side]
\label{rem:KL-vs-Riem}
On the rescued bands the two schemes attain the same leading complexity in each regime:
$\widetilde O(\kappa^2)$ single-sample iterations under \Cref{assump:global-Hess-Lip} and
$\widetilde O(\kappa^3)$ under \Cref{assump:logconcave-smooth} alone, both unconditionally. The KL band
$\tfrac1\beta\Id\preceq C_n\preceq\tfrac1\alpha\Id$ holds for every $\Delta t>0$ and gives the
stronger PL constant $\tfrac1\kappa$ (versus $\tfrac1{2\kappa}$ for the Riemannian band
$\tfrac1{2\beta}\Id\preceq C_n$), while the Riemannian scheme requires $\Delta t\leq1/\kappa$ to
maintain its band. The structural reason the KL analysis matches the Riemannian one is
\Cref{lem:KL-logtangent}: the implicit resolvent update is the exponential retraction
$C^{1/2}e^{W(P)}C^{1/2}$ with $W(P)$ a spectral function of the sampled relative Hessian $P$, so its
first-order stochastic part is unbiased and only a second-order fluctuation survives. Both floors are
proportional to $\Psi_H=N_\theta^2\barLH^2$ (Hessian-Lipschitz) or $\kappa N_\theta$ (self-bounding),
and both vanish at the estimator level for Gaussian targets (\Cref{prop:gauss-nofloor}).
\end{remark}
\section{Stochastic local convergence: Gaussian-core Price/Hessian--STL theory}
\label{sec:stoch-local}

The global stochastic theorems (\Cref{sec:stoch-Riem,sec:stoch-KL}) drive $\E\Delta_N$ to a
$\kappa$-dependent floor at the global $\kappa$-dependent rate. Once an iterate has entered a
certified local region $\Ureg_{\delta_\bullet}$, however, the \emph{deterministic} discrete local
map contracts at the rate of \Cref{thm:disc-local}. The local linearized spectral scale is
$4+\Gamma=O(1+\sqrt{N_\theta}\log\kappa)$ (after whitening, independent of the global covariance
conditioning); together with the universal certified stepsize $\Delta t=\Theta((4+\Gamma)^{-1})$ in the
large-batch branch $B\gtrsim V_1$ ($V_1=24+6N_\theta\LHs^2$) this yields the
$O((4+\Gamma)^2\log\tfrac1\epsilon)$ iteration bound derived below (\Cref{rem:local-condnum}),
quadratic in the scale (for general $B$, with the extra factor $\max\{1,V_1/B\}$). This section gives the stochastic analogue of \Cref{thm:disc-local}: a mean-square
contraction on the non-exit event, i.e.\ a bound on $\E[\norm{a_N-a_\star}_{\star,\bullet}^2;\tau>N]$,
with a \emph{target-adaptive} additive floor proportional to the \emph{squared} whitened
Hessian-Lipschitz modulus $\LHs^2$, together with a finite-horizon high-probability containment bound. This is a
\emph{killed-process} quantity (trajectories that have exited by time $N$ contribute zero), not the
mean square of the stopped process $a_{N\wedge\tau}$; the containment proposition supplies the exit
control that makes it operative. As already noted in \Cref{prop:entry-residual}, a gate alone yields
neither pathwise nor expected containment for the stochastic schemes; the results here are the
promised separate stability argument. The standing assumptions do not guarantee bounded support for
the STL mean-score noise: even for a Gaussian target, away from $C=\Id$ it equals
$(C^{-1/2}-C^{1/2})Z$ (see the Gaussian example below), which is unbounded---this single counterexample
already shows that pathwise invariance of a bounded local energy sublevel cannot be guaranteed under
the standing assumptions at fixed stepsize without projection, clipping, or safeguarding, so the
honest statements are non-exit contraction plus exit-probability control rather than infinite-horizon
containment. For Gaussian targets $\LHs=0$, the
\emph{additive} floor vanishes (\Cref{cor:stoch-local-gauss}); the noise is pathwise zero (deterministic
contraction, no exit control) whenever $C_0=C_\star$, even if $m_0\neq m_\star$, with quadratic rescue
the special case $a_0=a_\star$, while generic $C_0\neq C_\star$ retains unbounded mean noise.

\subsection{Whitened reduction and the local modulus}
\label{subsec:local-whiten}

The Fisher--Rao and KL natural-gradient schemes are equivariant under invertible affine
reparametrization of $\R^{N_\theta}$. Whitening by $C_\star^{1/2}$ maps $a_\star$ to $(0,\Id)$ and
$V$ to the whitened potential $\widehat V(z)=V(m_\star+C_\star^{1/2}z)$ of \Cref{sec:local}, under
which $\alpha_\star\Id\preceq\nabla^2\widehat V\preceq\beta_\star\Id$, $A(a_\star)=\Id$, and the
minibatch updates \eqref{eq:Riem-STL}--\eqref{eq:KL-STL} take the same form with $V\to\widehat V$,
$\nabla^2 V\to\nabla^2\widehat V$. We work in these coordinates throughout the section; $C$ denotes
the (whitened) covariance and $\overline H_n=\tfrac1B\sum_s\nabla^2\widehat V(X_{n,s})$, so that
$\alpha_\star\Id\preceq\overline H_n\preceq\beta_\star\Id$ almost surely by \eqref{eq:pathwise-Hess}.
Write
\[
    \overline\varepsilon_{b,n}:=\overline b_n-g(a_n),\qquad
    \overline\varepsilon_{H,n}:=\overline H_n-A(a_n),
\]
which are conditionally centered by \eqref{eq:STL-unbiased}; at the single-sample level
$\varepsilon_b:=\bigl(-\nabla\widehat V(X)+C^{-1}(X-m)\bigr)-g(a)$ and
$\varepsilon_H:=\nabla^2\widehat V(X)-A(a)=-\varepsilon_R$, so that
$\E\norm{\varepsilon_H}_{\Frob}^2$ is the (un-whitened-weighted) form of the intrinsic fluctuation
$\Psi$ of \Cref{def:Psi}.

\begin{lemma}[Whitened local Hessian-Lipschitz modulus]
\label{lem:whitened-HessLip}
Under \Cref{assump:global-Hess-Lip} (global modulus $\LH$), the whitened potential $\widehat V$ has a
Hessian-Lipschitz modulus $\LHs$, i.e.\
$\norm{\nabla^2\widehat V(x)-\nabla^2\widehat V(y)}_{\op}\leq\LHs\norm{x-y}$ for all $x,y$, with
\[
    \LHs\leq\LH\,\norm{C_\star}_{\op}^{3/2}\leq\LH\,\alpha^{-3/2}=\barLH .
\]
\end{lemma}

\begin{proof}
By the chain rule $\nabla^2\widehat V(z)=C_\star^{1/2}\nabla^2 V(m_\star+C_\star^{1/2}z)C_\star^{1/2}$.
Hence, for all $x,y$,
\[
    \norm{\nabla^2\widehat V(x)-\nabla^2\widehat V(y)}_{\op}
    =\norm{C_\star^{1/2}\bigl[\nabla^2V(m_\star+C_\star^{1/2}x)-\nabla^2V(m_\star+C_\star^{1/2}y)\bigr]C_\star^{1/2}}_{\op}.
\]
Using $\norm{C_\star^{1/2}MC_\star^{1/2}}_{\op}\leq\norm{C_\star^{1/2}}_{\op}^2\norm{M}_{\op}=\norm{C_\star}_{\op}\norm{M}_{\op}$
and then global Hessian-Lipschitzness with
$\norm{C_\star^{1/2}(x-y)}\leq\norm{C_\star^{1/2}}_{\op}\norm{x-y}=\norm{C_\star}_{\op}^{1/2}\norm{x-y}$,
\[
    \norm{\nabla^2\widehat V(x)-\nabla^2\widehat V(y)}_{\op}
    \leq\norm{C_\star}_{\op}\cdot\LH\norm{C_\star}_{\op}^{1/2}\norm{x-y}
    =\LH\norm{C_\star}_{\op}^{3/2}\norm{x-y}.
\]
Finally $C_\star=A(a_\star)^{-1}\preceq\alpha^{-1}\Id$ (as $A(a_\star)=\E_{\rho_\star}[\nabla^2V]\succeq\alpha\Id$),
so $\norm{C_\star}_{\op}^{3/2}\leq\alpha^{-3/2}$ and $\LHs\leq\LH\alpha^{-3/2}=\barLH$.
\end{proof}

\begin{definition}[Stochastically admissible level]
\label{def:stoch-admissible}
For a fixed scheme $\bullet\in\{\Riem,\KL\}$ and a fixed stepsize $\Delta t$, a level $\delta_\bullet>0$
is \emph{stochastically admissible} (for that $\Delta t$) if the discrete radius condition
$\Kdisc(\delta_\bullet)\,\mathfrak b_{\star,\bullet}(\delta_\bullet)\sqrt{\delta_\bullet}\leq\gamma_\bullet$
of \Cref{thm:disc-local} holds and additionally
\[
    \ell_{\delta_\bullet}\geq\tfrac12,\qquad U_{\delta_\bullet}\leq\tfrac32,\qquad
    \mathfrak b_{\star,\bullet}(\delta_\bullet)\sqrt{\delta_\bullet}\leq1 .
\]
This is a joint condition on the pair $(\Delta t,\delta_\bullet)$: $\Kdisc$, $\gamma_\bullet$ (for KL),
and $\mathfrak b_{\star,\bullet}$ all depend on $\Delta t$, so it is understood that $\Delta t$ is fixed
first---satisfying the linearized-gap and stochastic-absorption bounds of
\Cref{thm:stoch-local-Riem,thm:stoch-local-KL}---and $\delta_\bullet$ chosen afterward.
\end{definition}

\begin{remark}[Existence]
\label{rem:stoch-admissible-exists}
As $\delta\downarrow0$, the band roots satisfy $\ell_\delta\uparrow1$ and $U_\delta\downarrow1$
(\Cref{cor:band}), $\mathfrak b_{\star,\bullet}(\delta)\sqrt\delta\to0$, and, \emph{for every fixed
admissible stepsize $\Delta t>0$}, $\Kdisc(\delta)\to\Kdisc(0)<\infty$ as $\delta\downarrow0$ (the
supremum of the continuous $\Delta t^{-1}D^2T_{\bullet,\Delta t}$ over the shrinking hull
$\Hreg_\delta\to\{a_\star\}$, finite by \Cref{rem:Kdisc-finite}), so the discrete radius condition
$\Kdisc(\delta)\mathfrak b_{\star,\bullet}(\delta)\sqrt\delta\leq\gamma_\bullet$ holds for all small
$\delta$. Hence, for every fixed admissible $\Delta t$, sufficiently small stochastically admissible
levels $\delta_\bullet$ exist. On $\Ureg_{\delta_\bullet}$ one has
$\tfrac12\Id\preceq C\preceq\tfrac32\Id$ by \Cref{cor:band}(i), which is the band hypothesis of the
variance lemma below. This is an \emph{existence} statement: the manuscript does not give a
closed-form value of the switching level $\delta_\bullet$ in terms of $\alpha,\beta,N_\theta,\LH$
alone, since $\Kdisc$ is defined through a supremum of $D^2T_{\bullet,\Delta t}$ over the local hull
and is not bounded in closed form here; an explicit upper bound on $\Kdisc$ (a natural later
strengthening) would make the discrete switching threshold precomputable.
\end{remark}

\subsection{Local Price/Hessian--STL variance}
\label{subsec:local-var}

\begin{lemma}[Local variance decomposition]
\label{lem:local-variance}
Assume \Cref{assump:logconcave-smooth,assump:global-Hess-Lip}, work in whitened coordinates, and let
$a=(m,C)$ with $\tfrac12\Id\preceq C\preceq\tfrac32\Id$. Then for a single sample $X\sim\rho_a$,
\begin{equation}
\label{eq:local-var}
    \E\Bigl[\norm{\varepsilon_b}^2+\tfrac12\norm{\varepsilon_H}_{\Frob}^2\Bigr]
    \leq V_0+V_1\,\norm{a-a_\star}_\star^2,
    \qquad
    V_0:=3N_\theta^2\LHs^2,\quad V_1:=24+6N_\theta\LHs^2 .
\end{equation}
By conditional independence, the minibatch noise obeys
$\E[\norm{\overline\varepsilon_b}^2+\tfrac12\norm{\overline\varepsilon_H}_{\Frob}^2\mid\Fil]\leq B^{-1}(V_0+V_1\norm{a-a_\star}_\star^2)$.
At the optimizer, $\E[\norm{\varepsilon_b(a_\star)}^2+\tfrac12\norm{\varepsilon_H(a_\star)}_{\Frob}^2]\leq\sigma_\star^2:=\tfrac32N_\theta^2\LHs^2$,
which is $0$ for Gaussian targets.
\end{lemma}

\begin{proof}
Write $H(x):=\nabla^2\widehat V(x)$, so $\varepsilon_H=H(X)-A(a)$ and $A(a)=\E_{\rho_a}[H(X)]$.

\emph{Step 1 (covariance noise).} Gaussian Poincar\'e applied componentwise to $X\mapsto H_{ij}(X)$,
$X\sim\N(m,C)$, gives $\Var(H_{ij}(X))\leq\E\bigl[\nabla H_{ij}(X)^\top C\,\nabla H_{ij}(X)\bigr]\leq\norm{C}_{\op}\E\norm{\nabla H_{ij}(X)}^2$.
Summing over $i,j$,
\[
    \E\norm{\varepsilon_H}_{\Frob}^2=\sum_{i,j}\Var(H_{ij}(X))
    \leq\norm{C}_{\op}\,\E\sum_{i,j,k}\bigl(\partial_kH_{ij}(X)\bigr)^2 .
\]
For each $k$, $\sum_{i,j}(\partial_kH_{ij}(x))^2=\norm{DH(x)[e_k]}_{\Frob}^2\leq N_\theta\norm{DH(x)[e_k]}_{\op}^2\leq N_\theta\LHs^2$,
where $DH(x)[u]$ is the (symmetric-matrix-valued) directional derivative of the Hessian and
$\norm{DH(x)[u]}_{\op}\leq\LHs\norm u$ is exactly \Cref{lem:whitened-HessLip}. Summing over the
$N_\theta$ coordinates $k$ gives $\sum_{i,j,k}(\partial_kH_{ij})^2\leq N_\theta^2\LHs^2$, so with
$\norm{C}_{\op}\leq\tfrac32$,
\begin{equation}
\label{eq:covnoise}
    \E\norm{\varepsilon_H}_{\Frob}^2\leq\tfrac32N_\theta^2\LHs^2 .
\end{equation}

\emph{Step 2 (mean noise).} Let $F(x):=-\nabla\widehat V(x)+C^{-1}(x-m)$, so $\varepsilon_b=F(X)-\E F(X)$
and $DF(x)=C^{-1}-H(x)$ is symmetric. Gaussian Poincar\'e applied componentwise to $X\mapsto F_i(X)$
and summed gives
$\E\norm{\varepsilon_b}^2=\sum_i\Var(F_i(X))\leq\norm{C}_{\op}\E\sum_i\norm{\nabla F_i(X)}^2=\norm{C}_{\op}\E\norm{DF(X)}_{\Frob}^2=\norm{C}_{\op}\E\norm{C^{-1}-H(X)}_{\Frob}^2$.
Since $\E H(X)=A(a)$, the orthogonal (bias--variance) decomposition
$C^{-1}-H(X)=(C^{-1}-A(a))-\varepsilon_H$ has vanishing cross expectation, so
$\E\norm{C^{-1}-H(X)}_{\Frob}^2=\norm{C^{-1}-A(a)}_{\Frob}^2+\E\norm{\varepsilon_H}_{\Frob}^2$, whence
\begin{equation}
\label{eq:meannoise}
    \E\norm{\varepsilon_b}^2\leq\norm{C}_{\op}\Bigl(\norm{C^{-1}-A(a)}_{\Frob}^2+\E\norm{\varepsilon_H}_{\Frob}^2\Bigr)
    \leq\tfrac32\Bigl(\norm{C^{-1}-A(a)}_{\Frob}^2+\E\norm{\varepsilon_H}_{\Frob}^2\Bigr).
\end{equation}

\emph{Step 3 (combine).} Adding $\tfrac12\E\norm{\varepsilon_H}_{\Frob}^2$ to \eqref{eq:meannoise},
\[
    \E\Bigl[\norm{\varepsilon_b}^2+\tfrac12\norm{\varepsilon_H}_{\Frob}^2\Bigr]
    \leq 2\,\E\norm{\varepsilon_H}_{\Frob}^2+\tfrac32\norm{C^{-1}-A(a)}_{\Frob}^2 .
\]
By \eqref{eq:covnoise}, $2\E\norm{\varepsilon_H}_{\Frob}^2\leq3N_\theta^2\LHs^2=V_0$. For the second term,
$\norm{C^{-1}-A(a)}_{\Frob}^2\leq2\norm{C^{-1}-\Id}_{\Frob}^2+2\norm{A(a)-\Id}_{\Frob}^2$. Since
$C^{-1}-\Id=-C^{-1}(C-\Id)$ and $\norm{C^{-1}}_{\op}\leq2$ (from $C\succeq\tfrac12\Id$),
$\norm{C^{-1}-\Id}_{\Frob}\leq2\norm{C-\Id}_{\Frob}$, so $2\norm{C^{-1}-\Id}_{\Frob}^2\leq8\norm{C-\Id}_{\Frob}^2$.
For the Hessian term, $A(a)-\Id=\E_Z[H(m+C^{1/2}Z)-H(Z)]$ (using $A(a_\star)=\E_Z H(Z)=\Id$), so by
Jensen and \Cref{lem:whitened-HessLip},
\[
\begin{split}
    \norm{A(a)-\Id}_{\Frob}^2&\leq\E\norm{H(m+C^{1/2}Z)-H(Z)}_{\Frob}^2
    \leq N_\theta\LHs^2\,\E\norm{m+(C^{1/2}-\Id)Z}^2\\
    &=N_\theta\LHs^2\bigl(\norm m^2+\norm{C^{1/2}-\Id}_{\Frob}^2\bigr),
\end{split}
\]
where the last equality holds because the cross term vanishes by $\E Z=0$ (the $\norm m^2$ term is
deterministic). Moreover $\norm{C^{1/2}-\Id}_{\Frob}^2\leq\norm{C-\Id}_{\Frob}^2$ (eigenvalue-wise
$(\sqrt c-1)^2\leq(c-1)^2$). Hence
$\tfrac32\norm{C^{-1}-A(a)}_{\Frob}^2\leq12\norm{C-\Id}_{\Frob}^2+3N_\theta\LHs^2(\norm m^2+\norm{C-\Id}_{\Frob}^2)$.
Converting through $\norm{a-a_\star}_\star^2=\norm m^2+\tfrac12\norm{C-\Id}_{\Frob}^2$ gives
$\norm{C-\Id}_{\Frob}^2\leq2\norm{a-a_\star}_\star^2$ and
$\norm m^2+\norm{C-\Id}_{\Frob}^2=\norm{a-a_\star}_\star^2+\tfrac12\norm{C-\Id}_{\Frob}^2\leq2\norm{a-a_\star}_\star^2$,
so the state-dependent part is $\leq(24+6N_\theta\LHs^2)\norm{a-a_\star}_\star^2=V_1\norm{a-a_\star}_\star^2$.
This is \eqref{eq:local-var}. The minibatch bound is $\E\norm{\overline\varepsilon}^2=B^{-1}\E\norm\varepsilon^2$
for centered i.i.d.\ averages. At $a_\star$ we have $C=\Id$ ($\norm{C}_{\op}=1$) and
$C^{-1}-A(a_\star)=0$, so \eqref{eq:covnoise} tightens to $\E\norm{\varepsilon_H}_{\Frob}^2\leq N_\theta^2\LHs^2$
and \eqref{eq:meannoise} to $\E\norm{\varepsilon_b}^2\leq N_\theta^2\LHs^2$, giving
$\sigma_\star^2\leq\tfrac32N_\theta^2\LHs^2$; for Gaussian targets $\LHs=0$ and $\sigma_\star^2=0$.
\end{proof}

\subsection{One-step stochastic perturbation of the local maps}
\label{subsec:local-perturb}

We compare one stochastic step to the deterministic local step $T_\bullet$ of \Cref{thm:disc-local}.
The following elementary bounds on the matrix exponential are the only nonlinear ingredient.

\begin{lemma}[Exponential perturbation bounds]
\label{lem:exp-perturb}
Let $S,\widehat S\in\Sym(N_\theta)$ with $S\preceq\Id$, $\widehat S\preceq\Id$, put
$\Delta:=\widehat S-S$, and let $0<\Delta t\leq1$. Then
\begin{align*}
\textup{(i)}\quad & \norm{e^{\Delta t\widehat S}-e^{\Delta tS}}_{\Frob}\leq e\,\Delta t\,\norm{\Delta}_{\Frob},\\
\textup{(ii)}\quad & \norm{e^{\Delta t\widehat S}-e^{\Delta tS}-D\exp(\Delta tS)[\Delta t\Delta]}_{\Frob}
\leq\tfrac{e}{2}\,\Delta t^2\,\norm{\Delta}_{\op}\norm{\Delta}_{\Frob}.
\end{align*}
\end{lemma}

\begin{proof}
For symmetric $M$ with eigendecomposition $M=Q\diag(\lambda_i)Q^\top$, the Fr\'echet derivative in
the eigenbasis acts by the Loewner divided-difference matrix,
$\bigl(Q^\top D\exp(M)[E]\,Q\bigr)_{ij}=\Phi_{ij}(Q^\top EQ)_{ij}$ with
$\Phi_{ij}=(e^{\lambda_i}-e^{\lambda_j})/(\lambda_i-\lambda_j)$ (and $e^{\lambda_i}$ if $\lambda_i=\lambda_j$);
since $\Phi_{ij}\leq e^{\max(\lambda_i,\lambda_j)}\leq e^{\lambda_{\max}(M)}$, one has
$\norm{D\exp(M)[E]}_{\Frob}\leq e^{\lambda_{\max}(M)}\norm E_{\Frob}$. By the Duhamel identity
$e^{Y}-e^{X}=\int_0^1\!\!\int_0^1 e^{u M_s}(Y-X)e^{(1-u)M_s}\,\dd u\,\dd s$ with $M_s=X+s(Y-X)$; taking
$X=\Delta tS$, $Y=\Delta t\widehat S$ so that $M_s\preceq\Delta t\Id$ (convexity, $S,\widehat S\preceq\Id$)
and $\norm{e^{uM_s}}_{\op}\leq e^{u\Delta t}$ gives
$\norm{e^{\Delta t\widehat S}-e^{\Delta tS}}_{\Frob}\leq e^{\Delta t}\Delta t\norm\Delta_{\Frob}\leq e\Delta t\norm\Delta_{\Frob}$,
which is (i).

For (ii), Taylor's theorem with integral remainder gives
$e^{\Delta t\widehat S}-e^{\Delta tS}-D\exp(\Delta tS)[\Delta t\Delta]=\int_0^1(1-s)\,\Phi''(s)\,\dd s$
with $\Phi(s):=\exp(\Delta tS+s\Delta t\Delta)$ and $\Phi''(s)=D^2\exp(M_s)[\Delta t\Delta,\Delta t\Delta]$,
$M_s=\Delta tS+s\Delta t\Delta\preceq\Delta t\Id$. The second Fr\'echet derivative admits the simplex
representation $D^2\exp(M)[E,E]=2\!\int_{u_0+u_1+u_2=1,\,u_i\geq0}e^{u_0M}Ee^{u_1M}Ee^{u_2M}$, whose
Frobenius norm is bounded, using $\norm{e^{u_iM}}_{\op}\leq e^{u_i\lambda_{\max}(M)}$ and
$\norm{ABC}_{\Frob}\leq\norm A_{\op}\norm B_{\op}\norm C_{\Frob}$, by
\[
    2\cdot(\text{simplex volume }\tfrac12)\cdot e^{\lambda_{\max}(M)}\norm E_{\op}\norm E_{\Frob}
    =e^{\lambda_{\max}(M)}\norm E_{\op}\norm E_{\Frob}.
\]
Hence
$\norm{\Phi''(s)}_{\Frob}\leq e^{\Delta t}\Delta t^2\norm\Delta_{\op}\norm\Delta_{\Frob}$ and
$\int_0^1(1-s)\dd s=\tfrac12$ give (ii) with $e^{\Delta t}\leq e$.
\end{proof}

\begin{lemma}[One-step perturbation, Riemannian]
\label{lem:local-perturb-Riem}
Assume \Cref{assump:logconcave-smooth,assump:global-Hess-Lip}, work in whitened coordinates, and let
$a=(m,C)$ with $\tfrac12\Id\preceq C\preceq\tfrac32\Id$ and $\norm{T_{\Riem}(a)-a_\star}_\star\leq1$,
where $T_{\Riem}$ is the deterministic map of \eqref{eq:Riem-det} and $\widehat a^+$ the one stochastic
step \eqref{eq:Riem-STL} with batch $B$ and $0<\Delta t\leq1$. Then
\[
    \E\bigl[\norm{\widehat a^+-a_\star}_\star^2\mid\Fil\bigr]
    \leq\norm{T_{\Riem}(a)-a_\star}_\star^2+\frac{64\,\Delta t^2}{B}\bigl(V_0+V_1\norm{a-a_\star}_\star^2\bigr).
\]
\end{lemma}

\begin{proof}
Let $a^+:=T_{\Riem}(a)$ and $\widehat a^+-a^+=(\delta_m,\delta_C)$. Expanding the whitened square,
\begin{equation}
\label{eq:perturb-expand}
    \E\bigl[\norm{\widehat a^+-a_\star}_\star^2\mid\Fil\bigr]
    =\norm{a^+-a_\star}_\star^2
    +2\inner{a^+-a_\star}{\E[\widehat a^+-a^+\mid\Fil]}_\star
    +\E\bigl[\norm{\widehat a^+-a^+}_\star^2\mid\Fil\bigr].
\end{equation}
Write $\overline\varepsilon_b,\overline\varepsilon_H$ for the minibatch noises and abbreviate
$Q:=\E[\norm{\overline\varepsilon_b}^2+\tfrac12\norm{\overline\varepsilon_H}_{\Frob}^2\mid\Fil]$.

\emph{Second moment.} The mean components differ by $\delta_m=\Delta t\,C\,\overline\varepsilon_b$, so
$\E\norm{\delta_m}^2=\Delta t^2\E\norm{C\overline\varepsilon_b}^2\leq\tfrac94\Delta t^2\E\norm{\overline\varepsilon_b}^2$.
The covariance components are $a^+$-block $C^{1/2}e^{\Delta tS}C^{1/2}$ and $\widehat a^+$-block
$C^{1/2}e^{\Delta t\widehat S}C^{1/2}$ with $S=\Id-C^{1/2}A(a)C^{1/2}\preceq\Id$ and
$\widehat S=\Id-C^{1/2}\overline H C^{1/2}\preceq\Id$ ($A(a),\overline H\succeq0$), and
$\widehat S-S=-C^{1/2}\overline\varepsilon_H C^{1/2}=:\Delta_S$. By \Cref{lem:exp-perturb}(i) and
$\norm{\Delta_S}_{\Frob}\leq\norm C_{\op}\norm{\overline\varepsilon_H}_{\Frob}\leq\tfrac32\norm{\overline\varepsilon_H}_{\Frob}$,
\[
    \norm{\delta_C}_{\Frob}\leq\norm C_{\op}\norm{e^{\Delta t\widehat S}-e^{\Delta tS}}_{\Frob}
    \leq\tfrac32\cdot e\Delta t\cdot\tfrac32\norm{\overline\varepsilon_H}_{\Frob}
    =\tfrac{9e}{4}\Delta t\norm{\overline\varepsilon_H}_{\Frob}.
\]
Therefore $\E\norm{\widehat a^+-a^+}_\star^2=\E\norm{\delta_m}^2+\tfrac12\E\norm{\delta_C}_{\Frob}^2
\leq\tfrac94\Delta t^2\E\norm{\overline\varepsilon_b}^2+\tfrac{81e^2}{32}\Delta t^2\E\norm{\overline\varepsilon_H}_{\Frob}^2
\leq\tfrac{81e^2}{16}\Delta t^2\,Q$, the last step using $\tfrac94\leq\tfrac{81e^2}{16}$ and matching the
$\tfrac12\norm{\overline\varepsilon_H}_{\Frob}^2$ coefficient exactly.

\emph{Bias.} The mean bias vanishes, $\E[\delta_m\mid\Fil]=\Delta t\,C\,\E[\overline\varepsilon_b\mid\Fil]=0$.
For the covariance, $e^{\Delta t\widehat S}-e^{\Delta tS}=D\exp(\Delta tS)[\Delta t\Delta_S]+\mathcal R$
with $\E[D\exp(\Delta tS)[\Delta t\Delta_S]\mid\Fil]=D\exp(\Delta tS)[\Delta t\,\E\Delta_S]=0$, so
$\E[\delta_C\mid\Fil]=C^{1/2}\E[\mathcal R\mid\Fil]C^{1/2}$ and, by \Cref{lem:exp-perturb}(ii) with
$\norm{\Delta_S}_{\op}\norm{\Delta_S}_{\Frob}\leq\tfrac94\norm{\overline\varepsilon_H}_{\Frob}^2$,
\[
    \norm{\E[\delta_C\mid\Fil]}_{\Frob}\leq\norm C_{\op}\,\E\norm{\mathcal R}_{\Frob}
    \leq\tfrac32\cdot\tfrac e2\Delta t^2\cdot\tfrac94\E\norm{\overline\varepsilon_H}_{\Frob}^2
    =\tfrac{27e}{16}\Delta t^2\,\E\norm{\overline\varepsilon_H}_{\Frob}^2 .
\]
Since $\E[\delta_m\mid\Fil]=0$, $\norm{\E[\widehat a^+-a^+\mid\Fil]}_\star=\tfrac1{\sqrt2}\norm{\E[\delta_C\mid\Fil]}_{\Frob}$,
so with $\norm{a^+-a_\star}_\star\leq1$ and Cauchy--Schwarz the cross term of
\eqref{eq:perturb-expand} obeys
\[
    2\,\bigl|\inner{a^+-a_\star}{\E[\widehat a^+-a^+\mid\Fil]}_\star\bigr|
    \leq\tfrac{2}{\sqrt2}\cdot\tfrac{27e}{16}\Delta t^2\E\norm{\overline\varepsilon_H}_{\Frob}^2
    =\tfrac{27e}{8\sqrt2}\Delta t^2\E\norm{\overline\varepsilon_H}_{\Frob}^2
    \leq\tfrac{27e}{4\sqrt2}\Delta t^2\,Q,
\]
using $\E\norm{\overline\varepsilon_H}_{\Frob}^2=2\E[\tfrac12\norm{\overline\varepsilon_H}_{\Frob}^2]\leq2Q$.

\emph{Aggregate.} Adding the two contributions, the coefficient of $\Delta t^2Q$ in
\eqref{eq:perturb-expand} is $\tfrac{81e^2}{16}+\tfrac{27e}{4\sqrt2}<64$ (numerically $\approx50.4$).
Bounding $Q\leq B^{-1}(V_0+V_1\norm{a-a_\star}_\star^2)$ by \Cref{lem:local-variance} yields the claim.
\end{proof}

\begin{lemma}[One-step perturbation, KL]
\label{lem:local-perturb-KL}
Under the hypotheses of \Cref{lem:local-perturb-Riem} but for the KL maps
\eqref{eq:KL-det}/\eqref{eq:KL-STL} and the norm $\norm\cdot_{\star,\Delta t}$ (with
$\norm{T_{\KL}(a)-a_\star}_{\star,\Delta t}\leq1$),
\[
    \E\bigl[\norm{\widehat a^+-a_\star}_{\star,\Delta t}^2\mid\Fil\bigr]
    \leq\norm{T_{\KL}(a)-a_\star}_{\star,\Delta t}^2+\frac{68\,\Delta t^2}{B}\bigl(V_0+V_1\norm{a-a_\star}_\star^2\bigr).
\]
\end{lemma}

\begin{proof}
Let $a^+:=T_{\KL}(a)$, $P:=C^{-1}+\Delta t\,A(a)$, $\widehat P:=C^{-1}+\Delta t\,\overline H=P+\Delta t\,\overline\varepsilon_H$,
so the covariance blocks are $C^+=(1+\Delta t)P^{-1}$, $\widehat C^+=(1+\Delta t)\widehat P^{-1}$, and
$\delta_C=(1+\Delta t)(\widehat P^{-1}-P^{-1})$. Since $A(a),\overline H\succeq0$, $P\succeq C^{-1}$ and
$\widehat P\succeq C^{-1}$, giving $\norm{P^{-1}}_{\op},\norm{\widehat P^{-1}}_{\op}\leq\norm C_{\op}\leq\tfrac32$
(a.s.). The expansion \eqref{eq:perturb-expand} holds with $\norm\cdot_\star\to\norm\cdot_{\star,\Delta t}$;
its covariance weight is $\tfrac{1+\Delta t}2$.

\emph{Second moment.} $\delta_m=\Delta t\,C\overline\varepsilon_b$ as before, $\E\norm{\delta_m}^2\leq\tfrac94\Delta t^2\E\norm{\overline\varepsilon_b}^2$.
The first-order resolvent identity $\widehat P^{-1}-P^{-1}=-P^{-1}(\Delta t\,\overline\varepsilon_H)\widehat P^{-1}$
gives $\norm{\widehat P^{-1}-P^{-1}}_{\Frob}\leq\norm{P^{-1}}_{\op}\Delta t\norm{\overline\varepsilon_H}_{\Frob}\norm{\widehat P^{-1}}_{\op}\leq\tfrac94\Delta t\norm{\overline\varepsilon_H}_{\Frob}$,
hence $\norm{\delta_C}_{\Frob}\leq(1+\Delta t)\tfrac94\Delta t\norm{\overline\varepsilon_H}_{\Frob}\leq\tfrac92\Delta t\norm{\overline\varepsilon_H}_{\Frob}$.
As $\tfrac{1+\Delta t}2\leq1$, $\tfrac{1+\Delta t}2\E\norm{\delta_C}_{\Frob}^2\leq\tfrac{81}4\Delta t^2\E\norm{\overline\varepsilon_H}_{\Frob}^2$,
so $\E\norm{\widehat a^+-a^+}_{\star,\Delta t}^2\leq\tfrac94\Delta t^2\E\norm{\overline\varepsilon_b}^2+\tfrac{81}4\Delta t^2\E\norm{\overline\varepsilon_H}_{\Frob}^2\leq\tfrac{81}2\Delta t^2Q$.

\emph{Bias.} The second-order resolvent identity
\[
    \widehat P^{-1}-P^{-1}=-P^{-1}(\Delta t\overline\varepsilon_H)P^{-1}+P^{-1}(\Delta t\overline\varepsilon_H)P^{-1}(\Delta t\overline\varepsilon_H)\widehat P^{-1}
\]
has centered first term, so the conditional bias is the second term:
\[
    \norm{\E[\widehat P^{-1}-P^{-1}\mid\Fil]}_{\Frob}
    \leq\norm{P^{-1}}_{\op}^2\norm{\widehat P^{-1}}_{\op}\Delta t^2\E\bigl[\norm{\overline\varepsilon_H}_{\op}\norm{\overline\varepsilon_H}_{\Frob}\bigr]
    \leq(\tfrac32)^3\Delta t^2\E\norm{\overline\varepsilon_H}_{\Frob}^2
    =\tfrac{27}8\Delta t^2\E\norm{\overline\varepsilon_H}_{\Frob}^2 .
\]
Thus $\norm{\E[\delta_C\mid\Fil]}_{\Frob}\leq(1+\Delta t)\tfrac{27}8\Delta t^2\E\norm{\overline\varepsilon_H}_{\Frob}^2\leq\tfrac{27}4\Delta t^2\E\norm{\overline\varepsilon_H}_{\Frob}^2$,
and, the mean bias being zero,
\[
    \norm{\E[\widehat a^+-a^+\mid\Fil]}_{\star,\Delta t}=\sqrt{\tfrac{1+\Delta t}2}\,\norm{\E[\delta_C\mid\Fil]}_{\Frob}\leq\tfrac{27}4\Delta t^2\E\norm{\overline\varepsilon_H}_{\Frob}^2 .
\]
With $\norm{a^+-a_\star}_{\star,\Delta t}\leq1$, the cross term is at most
$2\cdot\tfrac{27}4\Delta t^2\E\norm{\overline\varepsilon_H}_{\Frob}^2\leq27\Delta t^2Q$.

\emph{Aggregate.} The coefficient of $\Delta t^2Q$ is $\tfrac{81}2+27=67.5<68$. Finally
$\norm{a-a_\star}_\star^2\leq\norm{a-a_\star}_{\star,\Delta t}^2$ (as $\tfrac{1+\Delta t}2\geq\tfrac12$), so
$Q\leq B^{-1}(V_0+V_1\norm{a-a_\star}_\star^2)$ of \Cref{lem:local-variance} gives the stated bound.
\end{proof}

\subsection{Localized non-exit mean-square contraction}
\label{subsec:local-contraction}

\begin{theorem}[Stochastic local convergence, Riemannian]
\label{thm:stoch-local-Riem}
Assume \Cref{assump:logconcave-smooth,assump:global-Hess-Lip}, work in whitened coordinates, and set
$A_{\Riem}:=4+\Gamma$, $\mu_{\Riem}:=1/(2A_{\Riem})$, $V_1=24+6N_\theta\LHs^2$. Fix a stepsize
\[
    0<\Delta t\leq\min\Bigl\{\tfrac{2}{4+\Gamma},\ \tfrac{B}{256\,A_{\Riem}\,V_1}\Bigr\},
\]
and let $\delta_{\Riem}$ be stochastically admissible for this $\Delta t$ (\Cref{def:stoch-admissible}).
Run \eqref{eq:Riem-STL} with batch $B$.
Let $\tau_{\Riem}:=\inf\{n\geq0:a_n\notin\Ureg_{\delta_{\Riem}}\}$ and $a_0\in\Ureg_{\delta_{\Riem}}$.
Then for all $N\geq0$,
\[
    \boxed{\;
    \E\bigl[\norm{a_N-a_\star}_\star^2\,\mathbf 1_{\{\tau_{\Riem}>N\}}\bigr]
    \leq\bigl(1-\mu_{\Riem}\Delta t\bigr)^N\norm{a_0-a_\star}_\star^2
    +\frac{384\,\Delta t\,(4+\Gamma)\,N_\theta^2\LHs^2}{B}.
    \;}
\]
The left-hand side is the \emph{non-exit} (killed) second moment: trajectories that leave
$\Ureg_{\delta_{\Riem}}$ by time $N$ contribute zero, so this is not the mean square of the stopped
process $a_{N\wedge\tau_{\Riem}}$; \Cref{prop:finite-horizon} bounds the discarded mass
$\Prob(\tau_{\Riem}\leq N)$, and the two together give an operative guarantee.
\end{theorem}

\begin{proof}
Write $r_n:=\norm{a_n-a_\star}_\star$, $\Fil_n=\sigma(a_0,\dots,a_n)$. On $\{n<\tau_{\Riem}\}$ we have
$a_n\in\Ureg_{\delta_{\Riem}}$, so $\tfrac12\Id\preceq C_n\preceq\tfrac32\Id$ (\Cref{rem:stoch-admissible-exists}),
$r_n\leq\mathfrak b_{\star,\Riem}(\delta_{\Riem})\sqrt{\delta_{\Riem}}\leq1$, and \Cref{lem:local-map-contraction}
(one-step contraction, needing only $\Delta t\leq\tfrac2{4+\Gamma}$, no energy-descent bound)
gives $\norm{T_{\Riem}(a_n)-a_\star}_\star\leq(1-\mu_{\Riem}\Delta t)\,r_n\leq r_n\leq1$
($\tfrac12\gamma_{\Riem}=\mu_{\Riem}$). Thus
\Cref{lem:local-perturb-Riem} applies:
\[
    \E[r_{n+1}^2\mid\Fil_n]\leq(1-\mu_{\Riem}\Delta t)^2r_n^2+\tfrac{64\Delta t^2}B\bigl(V_0+V_1r_n^2\bigr).
\]
Because $\Delta t\leq\tfrac2{4+\Gamma}$, $\mu_{\Riem}\Delta t=\tfrac{\Delta t}{2(4+\Gamma)}\leq\tfrac1{(4+\Gamma)^2}\leq\tfrac1{16}\leq\tfrac12$,
so $(1-\mu_{\Riem}\Delta t)^2\leq1-\tfrac32\mu_{\Riem}\Delta t$ (from $x^2\leq\tfrac12 x$ for $x\leq\tfrac12$).
The stepsize bound $\Delta t\leq\tfrac{B}{256A_{\Riem}V_1}$ gives $\tfrac{64\Delta t^2}BV_1\leq\tfrac12\mu_{\Riem}\Delta t$,
hence
\[
    \E[r_{n+1}^2\mid\Fil_n]\leq(1-\mu_{\Riem}\Delta t)r_n^2+\eta_{\Riem},
    \qquad \eta_{\Riem}:=\tfrac{64\Delta t^2}BV_0=\tfrac{192\Delta t^2N_\theta^2\LHs^2}{B}.
\]
Set $M_n:=\E[r_n^2\mathbf1_{\{\tau_{\Riem}>n\}}]$. As $\{\tau_{\Riem}>n+1\}\subseteq\{\tau_{\Riem}>n\}\in\Fil_n$,
\[
    M_{n+1}\leq\E\bigl[\mathbf1_{\{\tau_{\Riem}>n\}}\E[r_{n+1}^2\mid\Fil_n]\bigr]
    \leq\E\bigl[\mathbf1_{\{\tau_{\Riem}>n\}}\bigl((1-\mu_{\Riem}\Delta t)r_n^2+\eta_{\Riem}\bigr)\bigr]
    \leq(1-\mu_{\Riem}\Delta t)M_n+\eta_{\Riem}.
\]
Iterating from $M_0=r_0^2$ and summing the geometric series $\sum_{j\geq0}(1-\mu_{\Riem}\Delta t)^j=(\mu_{\Riem}\Delta t)^{-1}$,
\[
    M_N\leq(1-\mu_{\Riem}\Delta t)^Nr_0^2+\frac{\eta_{\Riem}}{\mu_{\Riem}\Delta t}
    =(1-\mu_{\Riem}\Delta t)^Nr_0^2+\frac{192\Delta t^2N_\theta^2\LHs^2/B}{\Delta t/(2(4+\Gamma))},
\]
which is the boxed bound.
\end{proof}

\begin{theorem}[Stochastic local convergence, KL]
\label{thm:stoch-local-KL}
Under the analogous hypotheses, set $A_{\KL}:=4+2\Delta t+\Gamma$, $\mu_{\KL}:=1/(2A_{\KL})$. Fix a
stepsize
\[
    0<\Delta t\leq\min\Bigl\{\tfrac2{4+\Gamma},\ \tfrac{B}{272\,A_{\KL}\,V_1}\Bigr\},
\]
let $\delta_{\KL}$ be stochastically admissible for this $\Delta t$, and put $\tau_{\KL}:=\inf\{n:a_n\notin\Ureg_{\delta_{\KL}}\}$;
running \eqref{eq:KL-STL} with batch $B$, one has for $a_0\in\Ureg_{\delta_{\KL}}$ and all $N\geq0$,
\[
    \boxed{\;
    \E\bigl[\norm{a_N-a_\star}_{\star,\Delta t}^2\,\mathbf 1_{\{\tau_{\KL}>N\}}\bigr]
    \leq\bigl(1-\mu_{\KL}\Delta t\bigr)^N\norm{a_0-a_\star}_{\star,\Delta t}^2
    +\frac{408\,\Delta t\,(4+2\Delta t+\Gamma)\,N_\theta^2\LHs^2}{B}.
    \;}
\]
As in the Riemannian case the left-hand side is the non-exit (killed) second moment, paired with the
exit bound of \Cref{prop:finite-horizon}.
\end{theorem}

\begin{proof}
Identical to \Cref{thm:stoch-local-Riem}, replacing $\norm\cdot_\star\to\norm\cdot_{\star,\Delta t}$,
$\gamma_{\Riem}\to\gamma_{\KL}$ (so $\mu_{\Riem}\to\mu_{\KL}$), \Cref{lem:local-perturb-Riem} by
\Cref{lem:local-perturb-KL} (coefficient $64\to68$), and $\eta_{\Riem}$ by
$\eta_{\KL}=\tfrac{68\Delta t^2}BV_0=\tfrac{204\Delta t^2N_\theta^2\LHs^2}B$. The step condition
$\Delta t\leq\tfrac{B}{272A_{\KL}V_1}$ gives $\tfrac{68\Delta t^2}BV_1\leq\tfrac12\mu_{\KL}\Delta t$, and
$\mu_{\KL}\Delta t\leq\tfrac12$ holds as before, so
$\E[r_{n+1}^2\mid\Fil_n]\leq(1-\mu_{\KL}\Delta t)r_n^2+\eta_{\KL}$ with $r_n:=\norm{a_n-a_\star}_{\star,\Delta t}$;
the non-exit recursion gives $M_N\leq(1-\mu_{\KL}\Delta t)^Nr_0^2+\eta_{\KL}/(\mu_{\KL}\Delta t)$, i.e.\ the
boxed bound.
\end{proof}

\subsection{Finite-horizon containment}
\label{subsec:finite-horizon}

The non-exit bounds contract the surviving mass on $\{\tau_\bullet>N\}$; to control the discarded mass
we bound the exit probability. Fix a radius $r_\bullet>0$ with $r_\bullet\leq\tfrac1{2\sqrt2}$ and
$\{a:\norm{a-a_\star}_{\star,\bullet}\leq r_\bullet\}\subseteq\Ureg_{\delta_\bullet}$ (both hold for
$r_\bullet$ small, the first forcing $\tfrac12\Id\preceq C\preceq\tfrac32\Id$ since then
$\norm{C-\Id}_{\Frob}\leq\sqrt2\,r_\bullet\leq\tfrac12$), and let
$\tau_{r_\bullet}:=\inf\{n\geq0:\norm{a_n-a_\star}_{\star,\bullet}>r_\bullet\}$.

\begin{proposition}[Exit probability]
\label{prop:finite-horizon}
In the setting of \Cref{thm:stoch-local-Riem} (resp.\ \Cref{thm:stoch-local-KL}), with
$q_\bullet:=1-\mu_\bullet\Delta t$ and $\eta_\bullet$ the corresponding one-step floor increment
($\eta_{\Riem}=192\Delta t^2N_\theta^2\LHs^2/B$, $\eta_{\KL}=204\Delta t^2N_\theta^2\LHs^2/B$), if
$a_0$ satisfies $\norm{a_0-a_\star}_{\star,\bullet}\leq r_\bullet$ then for every $N\geq1$,
\[
    \boxed{\;
    \Prob\bigl(\tau_{r_\bullet}\leq N\bigr)
    \leq\frac{\norm{a_0-a_\star}_{\star,\bullet}^2+N\,\eta_\bullet}{r_\bullet^2}.
    \;}
\]
(Freezing the process at the exit time removes the factor $(\mu_\bullet\Delta t)^{-1}$ that a
per-exit-time union bound would incur; the denominator $\mu_\bullet\Delta t$ remains genuine in the
\emph{stationary killed-moment floor} $\eta_\bullet/(\mu_\bullet\Delta t)$ of
\Cref{thm:stoch-local-Riem,thm:stoch-local-KL}, which sums repeated noise within the surviving
process, but not in the exit probability.) The two terms are controlled separately: the batch $B$
affects only $\eta_\bullet$, not the initial-condition term
$\norm{a_0-a_\star}_{\star,\bullet}^2/r_\bullet^2$, so containment needs both a \emph{deep-entry}
condition on $a_0$ and a batch condition. Indeed
$\Prob(\tau_{r_\bullet}\leq N)\leq\epsilon'$ holds whenever
\[
    \norm{a_0-a_\star}_{\star,\bullet}^2\leq\frac{\epsilon'}2\,r_\bullet^2
    \qquad\text{and}\qquad
    N\eta_\bullet\leq\frac{\epsilon'}2\,r_\bullet^2,
\]
the latter, on substituting $\eta_\bullet=c'_\bullet\Delta t^2N_\theta^2\LHs^2/B$ ($c'_{\Riem}=192$,
$c'_{\KL}=204$), becoming
\[
    B\geq\frac{2c'_\bullet\,N\,\Delta t^2\,N_\theta^2\LHs^2}{\epsilon'\,r_\bullet^2},
    \qquad 2c'_{\Riem}=384,\quad 2c'_{\KL}=408.
\]
At the certified stepsize $\Delta t=\Theta((4+\Gamma)^{-1})$ (where also
$\mu_\bullet=\Theta((4+\Gamma)^{-1})$), this batch requirement is smaller by a factor of order
$(4+\Gamma)^2$ than the condition $B\gtrsim N\Delta t\,N_\theta^2\LHs^2/(\epsilon'\mu_\bullet
r_\bullet^2)$ that the union-bound argument would demand. The deep-entry requirement motivates two
nested regions $\mathcal U_{\mathrm{in}}\subsetneq\mathcal U_{\mathrm{out}}$: the local phase is
entered on the smaller $\mathcal U_{\mathrm{in}}$ (guaranteeing small
$\norm{a_0-a_\star}_{\star,\bullet}$) while the exit time is measured against the larger
$\mathcal U_{\mathrm{out}}$ (radius $r_\bullet$); \Cref{cor:stoch-three-stage} supplies the
high-probability entry into $\mathcal U_{\mathrm{in}}$.
\end{proposition}

\begin{proof}
Since $\{a:\norm{a-a_\star}_{\star,\bullet}\leq r_\bullet\}\subseteq\Ureg_{\delta_\bullet}$, on
$\{n<\tau_{r_\bullet}\}$ the one-step recursion
$\E[r_{n+1}^2\mid\Fil_n]\leq q_\bullet r_n^2+\eta_\bullet$ of the proofs of
\Cref{thm:stoch-local-Riem,thm:stoch-local-KL} holds (with
$r_n:=\norm{a_n-a_\star}_{\star,\bullet}$). Define the process frozen at exit,
$R_n:=r_{n\wedge\tau_{r_\bullet}}$. On $\{\tau_{r_\bullet}>n\}$ one has $R_{n+1}=r_{n+1}$ and, since
$q_\bullet\leq1$,
$\E[R_{n+1}^2\mid\Fil_n]=\E[r_{n+1}^2\mid\Fil_n]\leq q_\bullet r_n^2+\eta_\bullet\leq R_n^2+\eta_\bullet$;
on $\{\tau_{r_\bullet}\leq n\}$, $R_{n+1}=R_n$. In either case
$\E[R_{n+1}^2\mid\Fil_n]\leq R_n^2+\eta_\bullet$, so iterating,
\[
    \E R_N^2\leq r_0^2+N\eta_\bullet .
\]
On $\{\tau_{r_\bullet}\leq N\}$ the frozen value satisfies $R_N=r_{\tau_{r_\bullet}}>r_\bullet$, so
Markov's inequality gives
$r_\bullet^2\,\Prob(\tau_{r_\bullet}\leq N)\leq\E[R_N^2\mathbf1_{\{\tau_{r_\bullet}\leq N\}}]\leq\E R_N^2\leq r_0^2+N\eta_\bullet$,
which is the claim.
\end{proof}

\subsection{Local rate, condition number, and the Gaussian limit}
\label{subsec:local-condnum}

\begin{remark}[Condition-number dependence and two stepsize regimes]
\label{rem:local-condnum}
By $\beta_\star\leq\kappa$, the definition \eqref{eq:Gamma} has $t_0\leq\max\{1,\log(4\kappa)\}$, so
\[
    \Gamma\leq\sqrt{N_\theta}\bigl(4\sqrt{2t_0}+4t_0+4\bigr)=O\bigl(1+\sqrt{N_\theta}\log\kappa\bigr),
    \qquad A_\bullet=4+O(\Delta t)+\Gamma=O\bigl(1+\sqrt{N_\theta}\log\kappa\bigr).
\]
Moreover, by \eqref{eq:linear-gap-bounds} the same proofs run with any certified pair
$\gamma_0\leq\gstar$, $\Lambda_0\geq\Lambda_\star$ in place of the universal
$(1/(4+\Gamma),(4+\Gamma)/2)$ (for the KL scheme, whose linearized map is self-adjoint in the weighted
$\norm\cdot_{\star,\Delta t}$ rather than $\norm\cdot_\star$, converting the spectral pair costs a
bounded factor $1+\Delta t\leq2$, absorbed into the $O(\cdot)$; see the target-adaptive discussion in
\Cref{rem:local-restart}): the stepsize window becomes
$\Delta t\leq\min\{1/\Lambda_0,\,B\gamma_0/(2c_\bullet V_1)\}$ and the count the target-adaptive
\[
    N_{\mathrm{loc}}=O\!\Bigl(\max\Bigl\{\frac{\Lambda_\star}{\gstar},\,\frac{V_1}{B\gstar^2}\Bigr\}\log\frac{r_0^2}\epsilon\Bigr),
\]
of which the universal displays below are the worst case; for well-conditioned targets
($\alpha_\star=\Omega(1)$) this is $O(\max\{\beta_\star,V_1/B\}\log\tfrac1\epsilon)$ even when the
universal $\Gamma$ estimate is pessimistic.
Crucially, \Cref{lem:local-map-contraction} removed the energy-descent bound
$\Delta t\leq1/(\beta_\star\bar\Lambda_\delta)$ from the stochastic local stepsize condition: since
$\bar\Lambda_\delta\geq1/\alpha_\star$ gives $\beta_\star\bar\Lambda_\delta\geq\kappa_\star$ (up to
$\kappa^2$), retaining that bound would have forced $\Delta t=O(\kappa^{-2})$ and contaminated the
local rate with a spurious polynomial-$\kappa$ factor. Write $A_0:=4+\Gamma$, so $A_{\Riem}=A_0$ and,
since $\Delta t\leq2/A_0\leq\tfrac12$ (as $A_0\geq4$), $A_{\KL}=A_0+2\Delta t\leq A_0+1$. The corrected
windows are
\[
    0<\Delta t\leq\min\Bigl\{\tfrac2{A_0},\ \tfrac{B}{256\,A_0\,V_1}\Bigr\}\ (\Riem),
    \qquad
    0<\Delta t\leq\min\Bigl\{\tfrac2{A_0},\ \tfrac{B}{272\,A_{\KL}\,V_1}\Bigr\}\ (\KL),
\]
with $V_1=24+6N_\theta\LHs^2$; the first (linearized-gap) restriction is $\Delta t\leq2/(4+\Gamma)$ for
both schemes. The Riemannian choice below is the largest value permitted by the two displayed
sufficient conditions; the KL choice is an explicit \emph{conservative} certified stepsize that avoids
the implicit dependence of $A_{\KL}$ on $\Delta t$ by replacing $A_{\KL}$ with $A_0+1$:
\[
    \Delta t_{\Riem}^\star:=\min\Bigl\{\tfrac2{A_0},\ \tfrac{B}{256\,A_0\,V_1}\Bigr\},
    \qquad
    \Delta t_{\KL}^\star:=\min\Bigl\{\tfrac2{A_0},\ \tfrac{B}{272\,(A_0+1)\,V_1}\Bigr\}\leq\tfrac{B}{272\,A_{\KL}\,V_1}.
\]
(The exact largest KL stepsize solves $2\Delta t^2+A_0\Delta t-\tfrac{B}{272V_1}\leq0$, i.e.\
$\Delta t\leq\min\{\tfrac2{A_0},\,\tfrac14(\sqrt{A_0^2+B/(34V_1)}-A_0)\}$; the conservative choice above is
simpler and of the same order.)
At $\Delta t_\bullet^\star$ the per-step non-exit factor is
$q_\bullet=1-\mu_\bullet\Delta t_\bullet^\star=1-\Theta\!\bigl(A_0^{-2}\min\{1,B/V_1\}\bigr)$ for both
schemes (using $\mu_{\KL}\Delta t_{\KL}^\star\geq\Delta t_{\KL}^\star/(2(A_0+1))$). The floor
$F_\bullet=c'_\bullet\Delta t\,A_\bullet N_\theta^2\LHs^2/B$ ($c'_{\Riem}=384$, $c'_{\KL}=408$)
evaluates, via $\Delta t_\bullet^\star A_\bullet=\Theta(\min\{2,B/(c_\bullet V_1)\})$ ($c_{\Riem}=256$,
$c_{\KL}=272$), to
\[
    F_\bullet=O\!\Bigl(N_\theta^2\LHs^2\min\{\tfrac1B,\tfrac1{V_1}\}\Bigr)
    =\begin{cases}O(N_\theta^2\LHs^2/B), & B\gtrsim V_1,\\[1mm] O(N_\theta^2\LHs^2/V_1), & B\lesssim V_1,\end{cases}
\]
so the simplified floor $O(N_\theta^2\LHs^2/B)$ holds precisely in the large-batch branch. Hence, for
both schemes,
\[
    N_{\mathrm{loc}}=\frac1{1-q_\bullet}\log\frac{r_0^2}\epsilon
    =O\!\Bigl((4+\Gamma)^2\max\{1,V_1/B\}\log\tfrac{r_0^2}\epsilon\Bigr).
\]
Two regimes result.
\emph{(i) Small-$\Gamma$ regime} (which includes sufficiently near-Gaussian targets). If $\Gamma=O(1)$
(e.g.\ $\beta_\star-1=O(1)$) and $V_1/B=O(1)$ (batch $B\gtrsim V_1$), then $A_\bullet=O(1)$,
$\Delta t_\bullet^\star=\Theta(1)$, and $N_{\mathrm{loc}}=O(\log\tfrac1\epsilon)$, \emph{independent of
$\kappa$}. (This no longer requires
$\beta_\star\bar\Lambda_\delta=O(1)$, which the removed constraint would have demanded; $\Gamma=O(1)$
alone does not bound $V_1$ or the admissible radius, so the batch condition $B\gtrsim V_1$ is stated
separately.)
\emph{(ii) Universal window.} When $\Gamma$ grows, the worst case
$\Gamma=\Theta(\sqrt{N_\theta}\log\kappa)$ gives, in the large-batch branch $B\gtrsim V_1$,
\[
    N_{\mathrm{loc}}=O\bigl((4+\Gamma)^2\log\tfrac1\epsilon\bigr)
    =O\bigl((1+\sqrt{N_\theta}\log\kappa)^2\log\tfrac1\epsilon\bigr)
    =O\bigl((1+N_\theta\log^2\kappa)\log\tfrac1\epsilon\bigr);
\]
for general $B$ this carries the extra factor $\max\{1,V_1/B\}$.
Thus, suppressing dimension and the non-Gaussianity factor $V_1/B$, the universal local complexity is
$O((\log\kappa)^2\log\tfrac1\epsilon)$. A genuinely intermediate
$O(\sqrt{N_\theta}\log\kappa\cdot\log\tfrac1\epsilon)$ bound, linear rather than quadratic in the
spectral scale $4+\Gamma$, is \emph{not} proved here: it
would require a sharper $\Gamma$-independent linearized-gap stepsize theorem. This is the same
limitation flagged in the deterministic three-stage analysis and is separate from the
polynomial-$\kappa$ global floor.
\end{remark}

\begin{corollary}[Gaussian targets: zero additive floor; deterministic under matched covariance; exact recovery under rescue]
\label{cor:stoch-local-gauss}
If $V(x)=\tfrac12(x-\mu)^\top H(x-\mu)+c$ with $H\succ0$, then $\LHs=0$, so $V_0=\sigma_\star^2=0$ and
$\eta_\bullet=0$. Under the stepsize hypotheses of the corresponding stochastic-local theorem
(\Cref{thm:stoch-local-Riem,thm:stoch-local-KL})---in particular $0<\Delta t\leq2/(4+\Gamma)=\tfrac12$
for Gaussian targets, since $\Gamma=0$, which gives $|1-\Delta t|<1$---three cases are distinguished by
the covariance initialization (whitened, so $C_\star=\Id$).
\begin{enumerate}
\item[(a)] \emph{Generic $C_0\neq\Id$.} The \emph{additive} floors in
\Cref{thm:stoch-local-Riem,thm:stoch-local-KL} vanish, and the exit bound of
\Cref{prop:finite-horizon} loses its $N$-growing term:
\[
    \begin{aligned}
    \E[\norm{a_N-a_\star}_{\star,\bullet}^2;\tau_\bullet>N]&\leq(1-\mu_\bullet\Delta t)^N\norm{a_0-a_\star}_{\star,\bullet}^2,\\
    \Prob(\tau_{r_\bullet}\leq N)&\leq\frac{\norm{a_0-a_\star}_{\star,\bullet}^2}{r_\bullet^2},
    \end{aligned}
\]
both uniformly in $N$.
These remain mean-square / in-probability statements: the STL mean-score noise
$\varepsilon_b=(C^{-1/2}-C^{1/2})Z$ is nonzero whenever $C\neq\Id$, so contraction is \emph{not}
pathwise and the exit probability is \emph{not} zero. The state-dependent constant $V_1=24$ survives
($V_1$ does not contain $\LHs$), so the mean noise persists along any trajectory with $C\neq\Id$.
\item[(b)] \emph{Matched covariance $C_0=\Id$ (i.e.\ $C_0=C_\star$), $m_0\neq0$.} Then
$\varepsilon_b=(C^{-1/2}-C^{1/2})Z=0$ and $\varepsilon_H=\nabla^2\widehat V-A(a)=0$ pathwise, both
covariance updates fix $C_n\equiv\Id$, and the mean update is \emph{deterministic},
$m_{n+1}=(1-\Delta t)m_n$ for both schemes. The iterate is therefore pathwise deterministic and
contracts to $a_\star$ with no stochastic noise---so no exit-probability control is needed---even
though $m_0\neq m_\star$. Thus zero-noise pathwise evolution requires only $C_0=C_\star$, not
$a_0=a_\star$.
\item[(c)] \emph{Quadratic-rescue initialization.} The rescue sets $a_0=a_\star$ exactly
(\Cref{thm:gauss-recovery}), the special case of (b) with $m_0=0$; then $a_n\equiv a_\star$ for all $n$
and the pipeline returns $a_\star$ after the single query. No local phase is needed.
\end{enumerate}
Part~(a) is the local mean-square form of the sticking-the-landing zero-\emph{additive}-floor property
(\Cref{prop:gauss-nofloor}); part~(b) is zero-noise pathwise contraction under matched covariance; and
exactness at $a_\star$ is special to (c).
\end{corollary}

\begin{proof}
For a quadratic potential $\nabla^2\widehat V\equiv C_\star^{1/2}HC_\star^{1/2}$ is constant, so
$\LHs=0$; then $V_0=3N_\theta^2\LHs^2=0$ and $\sigma_\star^2=\tfrac32N_\theta^2\LHs^2=0$ in
\Cref{lem:local-variance}, whence $\eta_\bullet\propto V_0=0$. Setting $\eta_\bullet=0$ in
\Cref{thm:stoch-local-Riem,thm:stoch-local-KL,prop:finite-horizon} gives the displayed non-exit and
exit bounds of (a); these are expectation/probability statements, and the whitened single-sample mean
noise is $\varepsilon_b=\bigl(-\nabla\widehat V(X)+C^{-1}(X-m)\bigr)-g(a)$, which for
$\widehat V(x)=\tfrac12\norm x^2$ and $X=m+C^{1/2}Z$ equals $-C^{1/2}Z+C^{-1/2}Z=(C^{-1/2}-C^{1/2})Z$,
nonzero for $C\neq\Id$, so no pathwise or zero-exit conclusion follows. For (b), $C_0=\Id$ gives
$\varepsilon_b=(\Id-\Id)Z=0$ and, since $\nabla^2\widehat V\equiv\Id$, $\overline H_n=\Id=A(a)$ so
$\varepsilon_H=0$, pathwise; the Riemannian covariance update returns
$C_1=C_0^{1/2}\exp(\Delta t(\Id-C_0^{1/2}\overline H_0C_0^{1/2}))C_0^{1/2}=\exp(\Delta t\cdot0)=\Id$ and
the KL update $C_1=(1+\Delta t)(\Id+\Delta t\Id)^{-1}=\Id$, so $C_n\equiv\Id$; the mean update is then
$m_{n+1}=m_n+\Delta t\,C_n\overline b_n=m_n+\Delta t(-m_n)=(1-\Delta t)m_n$ (using
$\overline b_n=g(a_n)=-m_n$), deterministic. For (c), \Cref{thm:gauss-recovery} gives
$a_0=a_\star=(0,\Id)$ (the $m_0=0$ instance of (b)), so $m_n\equiv0$, $C_n\equiv\Id$, i.e.\
$a_n\equiv a_\star$.
\end{proof}

\subsection{Global-to-local entry}
\label{subsec:global-to-local}

The local theorems assume entry into the certified region. The following corollary supplies that
entry from the global stochastic guarantees, completing the stochastic three-stage story with a
single union bound. It uses the two nested regions of \Cref{prop:finite-horizon}: entry is certified
on the inner level $\delta_{\mathrm{in}}$, while the exit time is measured against the outer ball of
radius $r_\bullet$.

\begin{corollary}[Stochastic global-to-local pipeline]
\label{cor:stoch-three-stage}
Fix a scheme $\bullet$, a stochastically admissible outer level $\delta_{\mathrm{out}}$ with radius
$r_\bullet\leq\tfrac1{2\sqrt2}$ satisfying $\{a:\norm{a-a_\star}_{\star,\bullet}\leq r_\bullet\}\subseteq\Ureg_{\delta_{\mathrm{out}}}$,
and stepsize/batch $(\Delta t,B)$ satisfying the hypotheses of \Cref{thm:stoch-local-Riem} (resp.\
\Cref{thm:stoch-local-KL}) at level $\delta_{\mathrm{out}}$. Fix failure budgets
$\epsilon_{\mathrm{ent}},\epsilon'\in(0,1)$ and an inner level $\delta_{\mathrm{in}}\leq\delta_{\mathrm{out}}$
small enough for the deep-entry condition
\[
    \mathfrak b_{\star,\bullet}(\delta_{\mathrm{in}})^2\,\delta_{\mathrm{in}}
    \leq\frac{\epsilon'}2\,r_\bullet^2 .
\]
Run a stochastic global scheme (\Cref{thm:Riem-Hlip,thm:Riem-selfbound,thm:KL-Hlip,thm:KL-selfbound}
at constant step, or \Cref{thm:decreasing}) for $N_1$ steps until $\E\Delta_{N_1}\leq\epsilon_{\mathrm{ent}}\delta_{\mathrm{in}}$
(achievable when the global floor is below $\epsilon_{\mathrm{ent}}\delta_{\mathrm{in}}$ at constant step,
or for $N_1$ large under the decreasing schedule), then switch to the local scheme with batch
$B\geq c_\bullet N_2\Delta t^2\,N_\theta^2\LHs^2/(\epsilon'r_\bullet^2)$ ($c_{\Riem}=384$,
$c_{\KL}=408$). Write $E_{\mathrm{in}}:=\{a_{N_1}\in\Ureg_{\delta_{\mathrm{in}}}\}$ and the shifted exit
time $\tau^{(N_1)}:=\inf\{n\geq0:\norm{a_{N_1+n}-a_\star}_{\star,\bullet}>r_\bullet\}$. Then for any local
horizon $N_2$,
\[
    \Prob\bigl(E_{\mathrm{in}}\cap\{\tau^{(N_1)}>N_2\}\bigr)\geq 1-\epsilon_{\mathrm{ent}}-\epsilon',
\]
and the entry-restricted non-exit second moment obeys
\[
    \E\bigl[\norm{a_{N_1+N_2}-a_\star}_{\star,\bullet}^2\,;\,E_{\mathrm{in}},\,\tau^{(N_1)}>N_2\bigr]
    \leq(1-\mu_\bullet\Delta t)^{N_2}\,\mathfrak b_{\star,\bullet}(\delta_{\mathrm{in}})^2\delta_{\mathrm{in}}
    +F_\bullet,
\]
with the explicit floors
$F_{\Riem}=\dfrac{384\,\Delta t\,(4+\Gamma)\,N_\theta^2\LHs^2}{B}$ and
$F_{\KL}=\dfrac{408\,\Delta t\,(4+2\Delta t+\Gamma)\,N_\theta^2\LHs^2}{B}$. The restriction to
$E_{\mathrm{in}}$ is necessary: on $E_{\mathrm{in}}^c$ the iterate may lie in the outer ball with
$\norm{a_{N_1}-a_\star}_{\star,\bullet}$ only bounded by $r_\bullet$ (not $r_0$), so the local
contraction does not control that mass.
\end{corollary}

\begin{proof}
By Markov's inequality applied to $\Delta_{N_1}\geq0$ and $\Ureg_{\delta_{\mathrm{in}}}=\{\Delta\leq\delta_{\mathrm{in}}\}$,
\[
    \Prob\bigl(a_{N_1}\notin\Ureg_{\delta_{\mathrm{in}}}\bigr)=\Prob(\Delta_{N_1}>\delta_{\mathrm{in}})
    \leq\frac{\E\Delta_{N_1}}{\delta_{\mathrm{in}}}\leq\epsilon_{\mathrm{ent}} .
\]
On $E_{\mathrm{in}}$, \Cref{cor:band} gives
$\norm{a_{N_1}-a_\star}_{\star,\bullet}\leq\mathfrak b_{\star,\bullet}(\delta_{\mathrm{in}})\sqrt{\delta_{\mathrm{in}}}=:r_0$,
and the deep-entry choice makes $r_0^2\leq\tfrac{\epsilon'}2r_\bullet^2\leq r_\bullet^2$, so
$a_{N_1}$ lies in the outer ball. Since
$E_{\mathrm{in}}\in\Fil_{N_1}$, apply \Cref{prop:finite-horizon} to the process
$(a_{N_1+n})_{n\geq0}$ started at $a_{N_1}$, on the event $E_{\mathrm{in}}$: the deep-entry condition
controls the initial-condition term ($\leq\epsilon'/2$) and the batch condition controls the
$N_2\eta_\bullet$ term ($\leq\epsilon'/2$), giving
$\Prob(\{\tau^{(N_1)}\leq N_2\}\cap E_{\mathrm{in}})\leq\epsilon'$. Hence
$\Prob(E_{\mathrm{in}}\cap\{\tau^{(N_1)}>N_2\})=\Prob(E_{\mathrm{in}})-\Prob(E_{\mathrm{in}}\cap\{\tau^{(N_1)}\leq N_2\})
\geq(1-\epsilon_{\mathrm{ent}})-\epsilon'$, the stated survival bound.

For the second moment, write $\rho_n:=\norm{a_{N_1+n}-a_\star}_{\star,\bullet}$. On
$\{n<\tau^{(N_1)}\}\cap E_{\mathrm{in}}$ the outer ball contains $a_{N_1+n}$, so the one-step recursion
$\E[\rho_{n+1}^2\mid\Fil_{N_1+n}]\leq q_\bullet\rho_n^2+\eta_\bullet$ of
\Cref{thm:stoch-local-Riem,thm:stoch-local-KL} holds there; setting
$M_n:=\E[\rho_n^2\mathbf1_{E_{\mathrm{in}}}\mathbf1_{\{\tau^{(N_1)}>n\}}]$ and using
$\{\tau^{(N_1)}>n\}\cap E_{\mathrm{in}}\in\Fil_{N_1+n}$ gives
$M_{n+1}\leq q_\bullet M_n+\eta_\bullet\Prob(E_{\mathrm{in}})\leq q_\bullet M_n+\eta_\bullet$, so
$M_{N_2}\leq q_\bullet^{N_2}\E[\rho_0^2\mathbf1_{E_{\mathrm{in}}}]+\eta_\bullet/(\mu_\bullet\Delta t)
\leq q_\bullet^{N_2}r_0^2+F_\bullet$, with $F_\bullet=\eta_\bullet/(\mu_\bullet\Delta t)$ equal to the
displayed $F_{\Riem},F_{\KL}$ (using $\rho_0\leq r_0$ on $E_{\mathrm{in}}$). This is the entry-restricted
non-exit bound.
\end{proof}

\begin{remark}[What is and is not proved]
\label{rem:stoch-local-scope}
\Cref{thm:stoch-local-Riem,thm:stoch-local-KL} are \emph{non-exit} (killed-process) statements: they
bound the mean-square deviation on the survival event $\{\tau_\bullet>N\}$, and do not claim a
pathwise or stopped-process bound. They are localized because a bounded local energy sublevel cannot
be pathwise invariant at fixed stepsize: the standing assumptions do not guarantee bounded support for
the STL mean-score noise. Even for a standard Gaussian target, away from $C=\Id$ it equals
$(C^{-1/2}-C^{1/2})Z$ and so has unbounded support---this counterexample alone suffices---so a
projection, clipping, or safeguarding mechanism would be needed for pathwise containment. We do \emph{not} assert almost-sure eventual exit; the present assumptions simply do not
imply infinite-horizon almost-sure containment in a bounded region. Note that the sampled-Hessian
covariance noise $\varepsilon_H$ is, by contrast, bounded a.s.\ ($\alpha_\star\Id\preceq\overline H_n\preceq\beta_\star\Id$),
so it is specifically the mean/score component that obstructs pathwise invariance.
\Cref{prop:finite-horizon} converts the non-exit bound into the operative guarantee: with a deep-entry
condition on $a_0$ and a batch $B=\Omega(N\Delta t^2\,N_\theta^2\LHs^2/(\epsilon'r_\bullet^2))$, the
iterate stays in the band over horizon $N$ with high probability (the frozen-at-exit argument of
\Cref{prop:finite-horizon} keeps this requirement free of the $(\mu_\bullet\Delta t)^{-1}$ factor a
per-exit-time union bound would incur), and on that event the floor is
$O(\Delta t\,A_\bullet N_\theta^2\LHs^2/B)$. The floor
constants ($V_0,V_1$; the aggregate coefficients $64,68$; the exponential-remainder constant $e/2$)
are explicit and conservative, proved in \Cref{lem:local-variance,lem:exp-perturb,lem:local-perturb-Riem,lem:local-perturb-KL};
they are not optimized, and only their orders in $(\Delta t,N_\theta,\LHs,B,A_\bullet)$ enter the
conclusions.
\end{remark}

\section{Floor-free convergence via decreasing stepsizes}
\label{sec:decreasing}

At constant stepsize the proved bounds generally retain an additive residual term.
\emph{Under Hessian-Lipschitzness} (\Cref{assump:global-Hess-Lip}), a common decreasing schedule removes
this certified residual and gives floor-free
$O(1/N)$ convergence for non-Gaussian targets, unconditionally for both schemes. It is stated only in
the Hessian-Lipschitz regime: in the assumption-minimal self-bounding regime the stability threshold
is $\Delta t=\Theta(\kappa^{-2})$ rather than $1/(8\kappa)$, so the schedule
\eqref{eq:decreasing-schedule} does not apply as written and a separate schedule and recursion would
be required.

\begin{theorem}[Decreasing-stepsize floor-free convergence]
\label{thm:decreasing}
Assume \Cref{assump:logconcave-smooth,assump:global-Hess-Lip}, initialize with \eqref{eq:rescue},
and run either \eqref{eq:Riem-STL} or \eqref{eq:KL-STL} with the schedule
\begin{equation}
\label{eq:decreasing-schedule}
    \Delta t_n=\frac{8\kappa}{n+n_0},
    \qquad
    n_0\geq64\kappa^2 .
\end{equation}
Then $\Delta t_n\leq1/(8\kappa)$ for all $n\geq0$, and with $c_{\mathsf S}:=\tfrac74$,
$\Psi_H=N_\theta^2\barLH^2$,
\begin{equation}
\label{eq:decreasing-rate}
    \boxed{\;
    \E\Delta_N\leq\frac{K}{N+n_0},
    \qquad
    K:=\max\Bigl\{\frac{64\,c_{\mathsf S}\,\kappa^3\Psi_H}B,\;n_0\,\Delta(a_0)\Bigr\}.
    \;}
\end{equation}
In particular, with the canonical offset $n_0=\lceil64\kappa^2\rceil$,
\begin{equation}
\label{eq:decreasing-rate-canonical}
    K=\widetilde O\!\Bigl(\frac{\kappa^3N_\theta^2\barLH^2}B+\kappa^2\Delta(a_0)\Bigr);
\end{equation}
for a larger offset $n_0\gg\kappa^2$ the initialization contribution $n_0\Delta(a_0)$ scales with the
chosen $n_0$ rather than $\kappa^2$.
\end{theorem}

\begin{proof}
Since $n+n_0\geq n_0\geq64\kappa^2$, $\Delta t_n=8\kappa/(n+n_0)\leq8\kappa/(64\kappa^2)=1/(8\kappa)$,
and the covariance bands persist for this varying schedule. For KL this follows at each step from the
same induction as in \Cref{lem:stoch-band}(2), which allows every positive stepsize. For the Riemannian
scheme, suppose
$\tfrac1{2\beta}\Id\preceq C_n\preceq\tfrac1\alpha\Id$. Since
$\Delta t_n\kappa\leq1$ and $\Delta t_n\leq1$, the upper-band argument gives
\[
    C_{n+1}^{-1}\succeq(1-\Delta t_n)C_n^{-1}+\Delta t_n\overline H_n\succeq\alpha\Id .
\]
The lower-envelope argument, now used with the current stepsize, gives
\[
    C_{n+1}^{-1}
    \preceq e^{-\Delta t_n}\bigl[C_n^{-1}+(e-1)\Delta t_n\overline H_n\bigr]
    \preceq\beta e^{-\Delta t_n}\bigl[2+(e-1)\Delta t_n\bigr]\Id
    \preceq2\beta\Id,
\]
where $e^{-h}[2+(e-1)h]\leq2$ for $0\leq h\leq1$. The rescued initializer lies in both bands, so
induction proves their persistence. Hence the constant-stepsize one-step estimates apply at each $n$.
For the Riemannian scheme,
\Cref{lem:Riem-onestep-stoch} with $\Delta t_n\leq1/(8\kappa)$ gives bracket
$1-\kappa\Delta t_n(1+2/B)\geq1-\tfrac38\geq\tfrac12$, so the drift is $\geq\tfrac{\Delta t_n}2\Gradnorm^2$;
the Riemannian PL bound $\Gradnorm^2\geq\tfrac1{2\kappa}\Delta$ gives drift $\geq\tfrac{\Delta t_n}{4\kappa}\Delta_n$,
and the floor coefficient is $\tfrac32\leq c_{\mathsf S}$. For the KL scheme,
\Cref{lem:KL-onestep-stoch} gives drift $\tfrac{\Delta t_n}2\Gradnorm^2\geq\tfrac{\Delta t_n}{2\kappa}\Delta_n\geq\tfrac{\Delta t_n}{4\kappa}\Delta_n$
(KL PL $\Gradnorm^2\geq\tfrac1\kappa\Delta$) and floor coefficient $\tfrac74=c_{\mathsf S}$. In both
cases, with $\Psi\leq\Psi_H$ (\Cref{lem:Psi-Hg}),
\begin{equation}
\label{eq:decreasing-recursion}
    \E\Delta_{n+1}\leq\Bigl(1-\frac{\Delta t_n}{4\kappa}\Bigr)\E\Delta_n+\frac{c_{\mathsf S}\kappa\Delta t_n^2}B\Psi_H .
\end{equation}
Write $t:=n+n_0$, so $\tfrac{\Delta t_n}{4\kappa}=\tfrac2t$ and
$\tfrac{c_{\mathsf S}\kappa\Delta t_n^2}B\Psi_H=\tfrac{64c_{\mathsf S}\kappa^3\Psi_H/B}{t^2}=\tfrac S{t^2}$
with $S:=64c_{\mathsf S}\kappa^3\Psi_H/B$. \eqref{eq:decreasing-recursion} becomes
$\E\Delta_{n+1}\leq(1-\tfrac2t)\E\Delta_n+\tfrac S{t^2}$. We show $\E\Delta_n\leq K/t$ by induction
with $K:=\max\{S,n_0\Delta(a_0)\}$. Base case $n=0$: $\E\Delta_0=\Delta(a_0)\leq K/n_0$. Inductive
step: assuming $\E\Delta_n\leq K/t$,
\[
    \E\Delta_{n+1}\leq\Bigl(1-\frac2t\Bigr)\frac Kt+\frac S{t^2}
    =\frac Kt-\frac{2K-S}{t^2}
    \leq\frac Kt-\frac K{t^2}
    =\frac{K(t-1)}{t^2}
    \leq\frac K{t+1},
\]
using $2K-S\geq K$ (since $K\geq S$) and $(t-1)(t+1)=t^2-1\leq t^2$. Since $t+1=(n+1)+n_0$, the
induction closes, giving \eqref{eq:decreasing-rate}.
\end{proof}

\begin{remark}[Reading the floor-free rate]
\label{rem:decreasing-reading}
The bound \eqref{eq:decreasing-rate} is a genuine $O(1/N)$ decay to zero, not to a floor: the
$\Psi_H$-dependent term is carried inside $K$ and divided by $N+n_0$. For a Gaussian target
$\Psi_H=0$ and $\Delta(a_0)=0$ (\Cref{thm:gauss-recovery}), so $K=0$ and convergence is immediate.
The schedule requires only $\barLH<\infty$ to define $K$; the iteration itself does not use $\LH$.
\end{remark}

\section{Gaussian targets: the vanishing floor}
\label{sec:gauss-nofloor}

Every stochastic floor in \Cref{sec:stoch-Riem,sec:stoch-KL} originates from the intrinsic Hessian
fluctuation $\Psi$ entering the one-step estimates \eqref{eq:Riem-onestep-Psi} and
\eqref{eq:KL-onestep-Psi}. The Hessian-Lipschitz floors
(\Cref{thm:Riem-Hlip,thm:KL-Hlip}) display this directly, being proportional to
$\Psi_H=N_\theta^2\barLH^2$; the self-bounding floors (\Cref{thm:Riem-selfbound,thm:KL-selfbound})
no longer show $\Psi$ explicitly only because the coarse universal estimate of
\Cref{lem:Psi-self-bound} has already been substituted. Since $\Psi\equiv0$ for Gaussian targets at
the estimator level, all of them vanish once one returns to the intrinsic one-step; this is the
sticking-the-landing zero-floor property in the intrinsic metric.

\begin{proposition}[Zero intrinsic fluctuation for Gaussian targets]
\label{prop:gauss-nofloor}
If $V(x)=\tfrac12(x-\mu)^\top H(x-\mu)+c$ with $H\succ0$, then $\nabla^2V(X)=H$ for all $X$, so
$A(a)=H$ and $\Psi(a)\equiv0$ for every $a\in\Aspace$; the additive stochastic term in the
\emph{intrinsic} $\Psi$-dependent one-step inequalities of \Cref{sec:stoch-Riem,sec:stoch-KL} (the
Riemannian and KL intrinsic estimates) therefore vanishes \emph{unconditionally} (the coarse
self-bounding recursions, in which the universal bound on $\Psi$ has already been substituted, retain
explicit residual terms and do not visibly vanish).
Consequently, \emph{under the covariance-band and stepsize hypotheses of the corresponding global
one-step theorem}---the band $\tfrac1{2\beta}\Id\preceq C_n\preceq\tfrac1\alpha\Id$ (Riemannian) resp.\
$\tfrac1\beta\Id\preceq C_n\preceq\tfrac1\alpha\Id$ (KL) and the stepsize of \Cref{thm:Riem-Hlip}
resp.\ \Cref{thm:KL-Hlip}, with pathwise band persistence given by the stochastic covariance-band
\Cref{lem:stoch-band} and satisfied from $n=0$ by the quadratic-rescued pipeline
(\Cref{thm:disc-rescue-to-FR})---both schemes contract geometrically with no floor:
\[
    \E\Delta_N\leq(1-c\kappa^{-2})^N\Delta(a_0),
\]
with $c=\tfrac1{24}$ (Riemannian, $B=1$) or $c=\tfrac1{16}$ (KL). Under quadratic rescue $a_0=a_\star$
(\Cref{thm:gauss-recovery}), so $\Delta(a_0)=0$ and the iterates remain exactly at $a_\star$; the
geometric statement is then trivially satisfied but exact. (The factors $c\kappa^{-2}$ are
band-dependent; $\Psi\equiv0$ alone gives the pure descent but not the $\kappa^{-2}$ rate.)
\end{proposition}

\begin{proof}
$\nabla^2V\equiv H$ gives $\nabla^2V(X)-A(a)=H-H=0$ pathwise, so the integrand of \Cref{def:Psi}
vanishes and $\Psi\equiv0$. The floor must be read off the \emph{intrinsic} one-step inequalities,
in which $\Psi$ still appears explicitly, rather than the self-bounding boxed bounds
\eqref{eq:Riem-selfbound}, \eqref{eq:KL-selfbound}, whose floors no longer contain $\Psi$ after the
coarse universal estimate of \Cref{lem:Psi-self-bound} has been applied. Setting $\Psi=0$ in the
Riemannian one-step \eqref{eq:Riem-onestep-Psi} (and in \Cref{lem:KL-onestep-stoch}, at the KL level)
removes the noise term identically, leaving the pure descent
$\E[\Delta_{n+1}\mid\Fil_n]\leq\Delta_n-\tfrac{\Delta t}2\Gradnorm(a_n)^2$. The Fisher--Rao
gradient-domination inequality $\Gradnorm(a)^2\geq\alpha\lambda_{\min}(C)\Delta(a)$ is
\emph{scheme-specific through the covariance floor}: on the Riemannian band $C_n\succeq\tfrac1{2\beta}\Id$
it gives $\Gradnorm(a_n)^2\geq\tfrac1{2\kappa}\Delta_n$, whereas on the KL band $C_n\succeq\tfrac1\beta\Id$
it gives $\Gradnorm(a_n)^2\geq\tfrac1\kappa\Delta_n$ (these bands are the stochastic ones of
\Cref{lem:stoch-band}, which persist pathwise for the Price/Hessian schemes). Hence, with $\Psi=0$:
for Riemannian at $\Delta t=\tfrac1{6\kappa}$,
$\E[\Delta_{n+1}\mid\Fil_n]\leq(1-\tfrac{\Delta t}{4\kappa})\Delta_n=(1-\tfrac1{24\kappa^2})\Delta_n$; for
KL at $\Delta t=\tfrac1{8\kappa}$,
$\E[\Delta_{n+1}\mid\Fil_n]\leq(1-\tfrac{\Delta t}{2\kappa})\Delta_n=(1-\tfrac1{16\kappa^2})\Delta_n$.
Equivalently, $\Psi_H=N_\theta^2\barLH^2=0$ removes the Hessian-Lipschitz
floors \eqref{eq:Riem-Hlip}, \eqref{eq:KL-Hlip} directly, and the factors are those of
\Cref{thm:Riem-Hlip} ($B=1$: $1-\tfrac1{24\kappa^2}$) and \Cref{thm:KL-Hlip} ($1-\tfrac1{16\kappa^2}$).
Finally \Cref{thm:gauss-recovery} gives $a_0=a_\star$, hence $\Delta(a_0)=0$.
\end{proof}

\begin{remark}[Scope of the zero-floor property]
\label{rem:gauss-nofloor-scope}
The zero-floor property is specific to Gaussian targets: it requires the sampled Hessian
$\nabla^2V(X)$ to equal its mean $A(a)$ pathwise, i.e.\ a constant Hessian. For a non-Gaussian
strongly log-concave target, $\Psi$ is generally positive where $\nabla^2V$ varies, so the certified
constant-stepsize recursion generally retains a nonzero additive term. This upper bound is not a
matching lower bound on the algorithm's asymptotic error and does not establish that a decreasing
stepsize is necessary; \Cref{thm:decreasing} instead supplies a proved floor-free guarantee.
\end{remark}

\section{Complexity summary and regime comparison}
\label{sec:regime-comparison}

\Cref{tab:complexity} collects the stochastic complexities. Throughout, the rescue contributes one
deterministic query; the reported iteration counts are those of Phase~I. The four constant-stepsize
global rows and the decreasing-stepsize row are unconditional expectation bounds at single sample
($B=1$); the two local rows are localized non-exit bounds, conditional on certified entry into
$\Ureg_{\delta_\bullet}$, shown with general batch $B$ (minibatching is central to the local survival
result). Here $\widetilde O$ suppresses factors logarithmic in $\kappa$, $N_\theta$,
$\Delta(a_0)/\epsilon$, and $\barLH$.

\begin{table}[t]
\centering
\small
\begin{tabular}{llccc}
\hline
Scheme / regime & controlled qty & floor & contraction & $N$ (to $\epsilon$ above floor) \\
\hline
Riem., self-bounding & $\E\Delta_N$ & $O(\kappa N_\theta)$ & $1-\Theta(\kappa^{-3})$ & $\widetilde O(\kappa^3)$ \\
Riem., Hess.-Lip. & $\E\Delta_N$ & $O(\kappa N_\theta^2\barLH^2)$ & $1-\Theta(\kappa^{-2})$ & $\widetilde O(\kappa^2)$ \\
KL, self-bounding & $\E\Delta_N$ & $O(\kappa N_\theta)$ & $1-\Theta(\kappa^{-3})$ & $\widetilde O(\kappa^3)$ \\
KL, Hess.-Lip. & $\E\Delta_N$ & $O(\kappa N_\theta^2\barLH^2)$ & $1-\Theta(\kappa^{-2})$ & $\widetilde O(\kappa^2)$ \\
\hline
Riem., local (non-exit, $B{\gtrsim}V_1$) & $\E[r_N^2;\tau{>}N]$ & $O(N_\theta^2\LHs^2/B)$ & $1-\Theta((4+\Gamma)^{-2})$ & $\widetilde O((4+\Gamma)^2)$ \\
KL, local (non-exit, $B{\gtrsim}V_1$) & $\E[r_N^2;\tau{>}N]$ & $O(N_\theta^2\LHs^2/B)$ & $1-\Theta((4+\Gamma)^{-2})$ & $\widetilde O((4+\Gamma)^2)$ \\
\hline
Decreasing step & $\E\Delta_N$ & $0$ & --- & $\widetilde O\!\bigl(\tfrac{\kappa^3N_\theta^2\barLH^2+\kappa^2\Delta(a_0)}{\epsilon}\bigr)$ \\
Gaussian target & $\E\Delta_N$ & $0$ & --- & $0$ \\
\hline
\end{tabular}
\caption{Stochastic complexity across regimes; global rows are displayed at $B=1$, local rows retain
general batch $B$. The
iteration counts $N$ are to accuracy $\epsilon$ \emph{above the stated floor} (the constant-stepsize
rows do not converge below their floor). The
self-bounding rows use \Cref{assump:logconcave-smooth}; the Hessian-Lipschitz and decreasing-step
rows additionally use \Cref{assump:global-Hess-Lip}. Minibatching $B=\Theta(\kappa)$ trades the
self-bounding regime to $\widetilde O(\kappa^2)$ rounds at $\widetilde O(\kappa^3)$ total oracle
pairs. The decreasing-step row is at $B=1$ (general $B$ carries a $1/B$ on the stochastic term) and
uses the canonical offset $n_0=\lceil64\kappa^2\rceil$ of \eqref{eq:decreasing-rate-canonical}, under
which the initialization term is $\kappa^2\Delta(a_0)$; it has zero asymptotic floor. The Gaussian
row is exact after the one rescue query. The two \emph{local (non-exit)} rows
(\Cref{sec:stoch-local}) apply once the iterate has entered the certified region
$\Ureg_{\delta_\bullet}$: they control the non-exit (killed) second moment
$\E[\norm{a_N-a_\star}_{\star,\bullet}^2;\tau{>}N]$ (not $\E\Delta_N$; on $\Ureg_{\delta_\bullet}$ these
are related by the coercivity bound $\norm{a-a_\star}_\star^2\leq\mathfrak b(\delta_\bullet)^2\Delta(a)$
of \Cref{cor:band}, the direction used for the entry conversion) at the
target-adaptive local rate (with $\Delta t=\Theta((4+\Gamma)^{-1})$ under the universal window),
paired with the finite-horizon containment of \Cref{prop:finite-horizon} and the entry corollary
\Cref{cor:stoch-three-stage}; their floor is proportional to $\LHs^2\leq\barLH^2$ and its additive part
vanishes for Gaussian targets. The displayed local contraction $1-\Theta((4+\Gamma)^{-2})$ and count
$\widetilde O((4+\Gamma)^2)$ are the large-batch branch $B\gtrsim V_1$; for general $B$ they read
$1-\Theta\bigl((4+\Gamma)^{-2}\min\{1,B/V_1\}\bigr)$ and $\widetilde O\bigl((4+\Gamma)^2\max\{1,V_1/B\}\bigr)$
(\Cref{rem:local-condnum}). In the small-$\Gamma$ regime ($\Gamma=O(1)$, $B\gtrsim V_1$; including
sufficiently near-Gaussian targets) the count is $\widetilde O((4+\Gamma)^2)=O(1)$ in its
$\kappa$-dependence, i.e.\ $O(\log\tfrac1\epsilon)$ (only the squared form follows from the theorem,
though $4+\Gamma=O(1)$ makes $O(4+\Gamma)$ and $O((4+\Gamma)^2)$ coincide). The local $N$-column is a
\emph{round} complexity for contraction of the killed second moment; enforcing the finite-horizon
survival probability of \Cref{prop:finite-horizon} requires the additional horizon-dependent batch
$B=\Omega\bigl(N\,\Delta t^2N_\theta^2\LHs^2/(\epsilon' r_\bullet^2)\bigr)$ and may raise the total oracle
complexity, so these rows are not high-probability oracle-complexity results.}
\label{tab:complexity}
\end{table}

\paragraph{The $\kappa^3\to\kappa^2$ improvement.}
The single-sample gap between the two constant-stepsize regimes is exactly one factor of $\kappa$.
Under \Cref{assump:logconcave-smooth} alone, the self-bounding estimate
(\Cref{lem:Psi-self-bound}) charges the fluctuation $\Psi$ against the gradient norm, forcing the
stepsize down to $\Delta t=\Theta(\kappa^{-2})$ (through $\rho_B\asymp\kappa$ at $B=1$) and giving
contraction $1-\Theta(\kappa^{-3})$. Under the additional \Cref{assump:global-Hess-Lip}, the
fluctuation is controlled by the target's nonquadraticity directly (\Cref{lem:Psi-Hg}), decoupling
it from the gradient norm; the stepsize rises to $\Delta t=\Theta(\kappa^{-1})$ and the contraction
improves to $1-\Theta(\kappa^{-2})$. The floor changes correspondingly from $O(\kappa N_\theta)$,
which is dimension-linear and target-agnostic, to $O(\kappa N_\theta^2\barLH^2)$, which is
dimension-quadratic but proportional to the squared non-Gaussianity $\barLH^2$ and hence far
smaller for near-Gaussian targets.

\paragraph{Scope and limitations.}
The stochastic theory presupposes a Price/Hessian oracle (\Cref{def:oracle}); the almost-sure
spectral bound \eqref{eq:pathwise-Hess} it provides is what makes the covariance band forward
invariant pathwise (\Cref{lem:stoch-band}) and the global theorems unconditional. A gradient-only
covariance estimator does not enjoy this bound (\Cref{rem:price}), and the pathwise arguments do not
transfer to it. The Hessian-Lipschitz floor scales as $N_\theta^2$; reducing this to $N_\theta$ would
require a Frobenius-directional third-derivative assumption
$\norm{D(\nabla^2V)(x)[u]}_{\Frob}\leq L_{H,\Frob}\norm u$ (in place of the operator-norm
\Cref{assump:global-Hess-Lip}). Repeating the Poincar\'e proof under this stronger assumption gives
\[
    \Psi(a)\leq L_{H,\Frob}^2\norm C_\op^2\Tr C
    \leq \frac{N_\theta L_{H,\Frob}^2}{\alpha^3}
    \qquad\text{on }C\preceq\alpha^{-1}\Id,
\]
which is linear in $N_\theta$. By contrast, the present operator-directional assumption gives
$\LH^2\Tr(C^2)\Tr C$, which can scale as $N_\theta^2$ even for isotropic $C$. This stronger regularity
is not assumed here. Both the global Riemannian and KL
stochastic rates are unconditional (the KL analysis matches the Riemannian one through the
exponential-retraction identity of \Cref{lem:KL-logtangent}); the local rates are non-exit bounds,
conditional on certified entry and paired with finite-horizon exit control. All constants are stated
explicitly and are those derived in the one-step lemmas; no unspecified absolute constants enter the
boxed bounds.
The stochastic local theory (\Cref{sec:stoch-local}) is deliberately a \emph{non-exit}
(killed-process) statement, not a stopped-process bound: because the STL mean-score noise generally
has unbounded support under Gaussian sampling (the sampled-Hessian covariance noise, by contrast, is
bounded a.s.), and because the standing assumptions do not provide bounded STL mean-score noise,
pathwise invariance of a bounded local sublevel cannot be guaranteed at fixed stepsize without
projection, clipping, or another safeguard (the Gaussian $C\neq\Id$ counterexample $(C^{-1/2}-C^{1/2})Z$
shows no general pathwise-invariance theorem is possible); so the operative guarantee is mean-square
contraction on the survival event plus the
finite-horizon exit-probability bound of \Cref{prop:finite-horizon}. Under a deep-entry condition on
the initialization and a batch
$B=\Omega\bigl(N\Delta t^2N_\theta^2\LHs^2/(\epsilon'r_\bullet^2)\bigr)$ (for a prescribed failure
probability $\epsilon'$ and fixed outer radius $r_\bullet$), a horizon-$N$ run stays in the band with
finite-horizon high probability; \Cref{cor:stoch-three-stage} supplies the high-probability entry from
the global phase by Markov's inequality. We do not claim almost-sure infinite-horizon containment in general; for Gaussian
targets, however, the noise is pathwise zero whenever $C_0=C_\star$, yielding deterministic contraction
and \emph{exact pathwise containment} for any matched-covariance initialization in the certified region
(\Cref{cor:stoch-local-gauss}(b)). Quadratic rescue is the special case $a_0=a_\star$, for which exact
recovery and stationarity ($a_n\equiv a_\star$) occur immediately with no iterative phase. The explicit constants
($V_0,V_1$, the aggregate $64,68$, the exponential remainder $e/2$) are conservative and not
optimized (\Cref{rem:stoch-local-scope}); only their orders enter the conclusions.

\paragraph{Per-iteration linear algebra.}
The complexities above count Price/Hessian oracle pairs, not dense linear algebra. Each iteration
also performs, in the ambient dimension $N_\theta$: a matrix exponential and a symmetric
factorization for the Riemannian covariance update, or a symmetric positive-definite solve/inverse
for the KL resolvent update, each $O(N_\theta^3)$; $B$ oracle pairs; and $O(N_\theta^2)$ vector work
for the STL score. The one-time rescue is one gradient/Hessian pair and one SPD solve. Thus the
wall-clock cost per iteration is $O(N_\theta^3+B\,c_{\mathrm{oracle}})$, and oracle complexity should
not be read as total arithmetic complexity when $N_\theta$ is large.

\section{Conclusion and outlook}
\label{sec:conclusion}

\paragraph{What the geometry provides.}
The affine-invariant Fisher--Rao structure supplies the results that hold under only strong
log-concavity and smoothness (\Cref{assump:logconcave-smooth}), with no third-order regularity and no
pointwise Hessian oracle: the exact expectation $A(a)=\E_{\rho_a}[\nabla^2V]$ that the deterministic
schemes use may equivalently be represented through Gaussian integration by parts using target
gradients (\Cref{rem:price}), which cleanly separates the deterministic exact-expectation theory from
the stochastic pathwise Price/Hessian theory. The continuous covariance bootstrap
(\Cref{thm:cont-global}), its order-optimality (\Cref{prop:burnin-optimal}), the deterministic global
rates of the retraction and KL schemes (\Cref{thm:Riem-global,thm:KL-global}), and the Gaussian-core
local region (\Cref{sec:local}) are all geometry-intrinsic: the Gaussian core is dissipated exactly,
the non-Gaussian remainder is the only Taylor-bounded object, and the certified local region equals the
exact supremal Gaussian threshold in every dimension. The improved local spectral scale
$\Gamma=O(1+\sqrt{N_\theta}\log\kappa)$ (\Cref{lem:linear-diagonal-bounds}) and the exact spectral
sandwich $\max\{\alpha_\star,(4+\Gamma)^{-1}\}\leq\gstar\leq\Lambda_\star\leq\min\{\beta_\star,(4+\Gamma)/2\}$
(\Cref{prop:linearized}) are properties of the linearized Fisher--Rao field alone.

\paragraph{What the Price/Hessian oracle provides.}
The stochastic theory rests on the almost-sure spectral bound of the Price/Hessian oracle
(\Cref{def:oracle}, \eqref{eq:pathwise-Hess}). This is exactly what makes the covariance band forward
invariant pathwise (\Cref{lem:stoch-band}) and the stochastic global rates
(\Cref{thm:Riem-Hlip,thm:Riem-selfbound,thm:KL-Hlip,thm:KL-selfbound}) unconditional, and what lets the
STL estimator remove the additive variance floor entirely for well-specified Gaussian targets
(\Cref{sec:gauss-nofloor}). A gradient-only covariance estimator, whose samplewise spectrum is
unbounded, does not enjoy this bound (\Cref{rem:price}); the present pathwise confinement and floor
bounds do not transfer to it.

\paragraph{Open problems.}
Four strengthenings are left open. First, the local stepsize theorem is quadratic in the spectral
scale, giving the universal count $O((4+\Gamma)^2\log\tfrac1\epsilon)$; a genuinely intermediate
$O(\sqrt{N_\theta}\log\kappa\cdot\log\tfrac1\epsilon)$ bound---linear rather than quadratic in
$4+\Gamma$---would require a $\Gamma$-independent linearized-gap stepsize argument, an implicit or
accelerated local update, or a sharper relation between $\Lambda_\star$ and $\gstar$ than the sandwich
provides (\Cref{rem:local-condnum}). Second, the finite-horizon containment batch grows linearly in the
horizon, $B=\Omega(N\Delta t^2N_\theta^2\LHs^2/(\epsilon'r_\bullet^2))$, because the exit control uses
only accumulated second moments and Markov's inequality (\Cref{prop:finite-horizon}); a maximal
martingale concentration argument for the locally sub-Gaussian STL noise, or an explicit projection or
clipping safeguard, might reduce this to a logarithmic or horizon-free dependence, but this requires a
separate proof and is \emph{not} implied by Gaussian concentration as stated here. Third, the
Hessian-Lipschitz floor scales as $N_\theta^2$; reducing it to $N_\theta$ needs the stronger
Frobenius-directional third-derivative assumption discussed in the scope paragraph, under which the
corrected estimate $\Psi(a)\leq L_{H,\Frob}^2\norm C_\op^2\Tr C$ applies. Fourth, the discrete local admissible level
enters through $\Kdisc(\delta)$ (\Cref{rem:Kdisc-finite}), which is finite and computable once a level
is supplied but is not yet bounded in closed form by the global target parameters; an explicit
$\Kdisc$ bound would make the online global-to-local switch fully precomputable. Extending the
pathwise stochastic guarantees to gradient-only, diagonal, low-rank, or Hessian-vector-product oracles
is the broadest of these directions and is left to future work.
\appendix

\section{Proof of the deterministic KL/Bregman one-step estimate}
\label{app:KL-onestep}

This appendix proves \Cref{prop:KL-onestep} in full, with no imported constants. Throughout, fix
$a=(m,C)$ with $L_n\Id\preceq C\preceq\Lambda\Id$, write $A:=A(a)$, $g:=g(a)$, $R:=C^{-1}-A$,
$\LB=\beta\Lambda\geq1$, and assume $0<\Delta t\leq1/\LB$. Recall $\calE=\calE_1+\calE_2$ with
$\calE_1(a)=-\tfrac12\log\det C+\text{const}$ and $\calE_2(a)=\E_{\rho_a}[V]$, and the Bonnet--Price
identities $\nabla_m\calE_2=\E_{\rho_a}[\nabla V]=-g$, $\nabla_C\calE_2=\tfrac12\E_{\rho_a}[\nabla^2V]=\tfrac12A$.
The descent argument uses only convexity and $\beta$-smoothness of $V$ ($0\preceq\nabla^2V\preceq\beta\Id$),
which hold under \Cref{assump:logconcave-smooth}.

\subsection{Variational characterization and explicit maps}
\label{app:KL-var}

\begin{lemma}
\label{lem:KL-argmin}
The KL/Bregman update \eqref{eq:KL-det} is the forward--backward proximal step
\begin{equation}
\label{eq:KL-argmin}
    a_{n+1}=\arg\min_{a'=(m',C')}\Bigl\{\Delta t\,\inner{\nabla\calE_2(a)}{a'-a}+\Delta t\,\calE_1(a')+\KL(\rho_{a'}\|\rho_a)\Bigr\},
\end{equation}
which linearizes the potential $\calE_2$ explicitly and treats the entropy $\calE_1$ implicitly, and
its minimizer is
\begin{equation}
\label{eq:KL-maps}
    m_{n+1}=m+\Delta t\,C g,
    \qquad
    C_{n+1}=(1+\Delta t)(C^{-1}+\Delta t\,A)^{-1}.
\end{equation}
\end{lemma}

\begin{proof}
The Gaussian KL splits as
$\KL(\rho_{a'}\|\rho_a)=\tfrac12(m-m')^\top C^{-1}(m-m')+\tfrac12[\Tr(C^{-1}C')-N_\theta+\log\det C-\log\det C']$,
so the objective in \eqref{eq:KL-argmin} decouples in $m'$ and $C'$. The $m'$-part is
$-\Delta t\,g^\top(m'-m)+\tfrac12(m'-m)^\top C^{-1}(m'-m)$, minimized at $m'-m=\Delta t\,Cg$. The
$C'$-part, dropping $C'$-independent terms and using $\nabla_C\calE_2=\tfrac12A$ and
$\calE_1(C')=-\tfrac12\log\det C'$, is
\[
    \tfrac{\Delta t}2\Tr(AC')+\tfrac12\Tr(C^{-1}C')-\tfrac{1+\Delta t}2\log\det C'
    =\tfrac12\Tr\bigl((C^{-1}+\Delta t\,A)C'\bigr)-\tfrac{1+\Delta t}2\log\det C',
\]
whose stationarity condition $\tfrac12(C^{-1}+\Delta t\,A)=\tfrac{1+\Delta t}2(C')^{-1}$ gives
$C'=(1+\Delta t)(C^{-1}+\Delta t\,A)^{-1}$. Both objectives are strictly convex in their arguments
($C^{-1}\succ0$; $-\log\det$ convex), so the stationary points are the minimizers.
\end{proof}

\subsection{A scalar inequality for the covariance factor}
\label{app:KL-scalar}

\begin{lemma}
\label{lem:scalar-KL}
Let $0<\Delta t\leq1/\LB$, $\mu\in[0,\LB]$, and $b:=\sqrt{(1+\Delta t)/(1+\Delta t\,\mu)}$. Then
\begin{equation}
\label{eq:scalar-KL}
    \mu(b-1)-\log b+\frac{\LB}2(b-1)^2\leq-\frac1{2\Delta t}\bigl(b^2-1-2\log b\bigr).
\end{equation}
\end{lemma}

\begin{proof}
Since $\LB\leq1/\Delta t$ and $(b-1)^2\geq0$, it suffices to replace $\LB$ by $1/\Delta t$ on the left.
The relation between $b$ and $\mu$ is $\mu=\bigl((1+\Delta t)b^{-2}-1\bigr)/\Delta t$. Multiplying the
resulting inequality by $\Delta t$ and moving all terms to the left, define
\[
    G(b):=\Delta t\,\mu(b-1)-\Delta t\log b+\tfrac12(b-1)^2+\tfrac12(b^2-1-2\log b),
\]
so that $G(1)=0$ and a direct differentiation gives
$G'(b)=\bigl(b-1\bigr)\bigl(2b^3-(1+\Delta t)(b+2)\bigr)/b^3$. Since $0\leq\Delta t\,\mu\leq\Delta t\,\LB\leq1$,
$b\in[\sqrt{(1+\Delta t)/2},\sqrt{1+\Delta t}]$. The polynomial $p(b):=2b^3-(1+\Delta t)(b+2)$ is convex
on $(0,\infty)$, and at the endpoints
\[
    p\!\bigl(\sqrt{(1+\Delta t)/2}\bigr)=-2(1+\Delta t)<0,
    \qquad
    p\!\bigl(\sqrt{1+\Delta t}\bigr)=(1+\Delta t)\bigl(\sqrt{1+\Delta t}-2\bigr)<0
\]
(using $\Delta t\leq1$), so convexity gives $p(b)<0$ throughout the interval. Hence $G'(b)\geq0$ for
$b\leq1$ and $G'(b)\leq0$ for $b\geq1$, so $G$ attains its maximum at $b=1$, where it vanishes,
proving \eqref{eq:scalar-KL}.
\end{proof}

\subsection{Forward-KL descent and Wasserstein-gradient lower bound}
\label{app:KL-descent}

\begin{lemma}
\label{lem:KL-onestep-app}
For $0<\Delta t\leq1/\LB$,
\begin{equation}
\label{eq:KL-descent-app}
    \calE(a_{n+1})\leq\calE(a)-\frac1{\Delta t}\KL(\rho_{a_{n+1}}\|\rho_a),
\end{equation}
and, with $\grad^{W_2}\calE(a)=(-g,-R)$ and
$\norm{\grad^{W_2}\calE(a)}_{W_2}^2=\norm g^2+\Tr(RCR)$ (\Cref{lem:W2-PL}),
\begin{equation}
\label{eq:KL-lower}
    \KL(\rho_{a_{n+1}}\|\rho_a)\geq\frac{L_n\Delta t^2}{4(1+\Delta t\,\LB)^2}\,\norm{\grad^{W_2}\calE(a)}_{W_2}^2 .
\end{equation}
\end{lemma}

\begin{proof}
Set $M:=C^{1/2}AC^{1/2}$ and $B:=\sqrt{1+\Delta t}\,(\Id+\Delta t\,M)^{-1/2}$, a function of $M$ with
$0\preceq M\preceq\LB\Id$ and $C_{n+1}=C^{1/2}B^2C^{1/2}$. With $\xi\sim\N(0,\Id)$, couple
$X:=m+C^{1/2}\xi\sim\rho_a$ and $Y:=m+\Delta t\,Cg+C^{1/2}B\xi$; since $B$ is symmetric, $Y$ is Gaussian
with mean $m_{n+1}$ and covariance $C^{1/2}B^2C^{1/2}=C_{n+1}$, i.e.\ $Y\sim\rho_{a_{n+1}}$.

\emph{Descent.} The entropy change is $\calE_1(a_{n+1})-\calE_1(a)=-\log\det B$. Writing $E:=B-\Id$,
the $\beta$-smoothness estimate of \Cref{lem:det-descent}(Step~3) applies verbatim with $E$ in place of
$e^{\Delta t S/2}-\Id$ (both are symmetric functions of $M$, and the computation there uses only
$V(Y)\leq V(X)+\nabla V(X)^\top(Y-X)+\tfrac\beta2\norm{Y-X}^2$, $\E[\nabla V(X)]=-g$, the Price identity
$\E[\xi\,\nabla V(X)^\top]=C^{1/2}A$, and $C\preceq\Lambda\Id$), giving
\[
    \calE_2(a_{n+1})-\calE_2(a)
    \leq-\Delta t\,g^\top Cg+\Tr(EM)+\frac{\LB\Delta t^2}2g^\top Cg+\frac{\LB}2\norm E_{\Frob}^2 .
\]
Because $\LB\Delta t\leq1$, the mean part obeys
$-\Delta t\,g^\top Cg+\tfrac{\LB\Delta t^2}2g^\top Cg\leq-\tfrac{\Delta t}2g^\top Cg$, which is exactly
$-\Delta t^{-1}$ times the mean part $\tfrac{\Delta t^2}2g^\top Cg$ of $\KL(\rho_{a_{n+1}}\|\rho_a)$. For
the covariance part, let $\mu_i$ be the eigenvalues of $M$ and $b_i=\sqrt{(1+\Delta t)/(1+\Delta t\,\mu_i)}$
the aligned eigenvalues of $B$; adding the entropy change, the covariance contribution is
$\sum_i\bigl[\mu_i(b_i-1)-\log b_i+\tfrac{\LB}2(b_i-1)^2\bigr]$, which by \Cref{lem:scalar-KL} is at most
$-\tfrac1{2\Delta t}\sum_i(b_i^2-1-2\log b_i)$. Since $C^{-1}C_{n+1}$ has eigenvalues $b_i^2$, this is
exactly $-\Delta t^{-1}$ times the covariance part
$\tfrac12\sum_i(b_i^2-1-\log b_i^2)$ of $\KL(\rho_{a_{n+1}}\|\rho_a)$. Adding the mean and covariance
parts proves \eqref{eq:KL-descent-app}.

\emph{Lower bound.} The mean part of the forward KL is
$\tfrac12(m_{n+1}-m)^\top C^{-1}(m_{n+1}-m)=\tfrac{\Delta t^2}2g^\top Cg\geq\tfrac{L_n\Delta t^2}2\norm g^2$.
For the covariance part, set $S:=C^{1/2}RC^{1/2}=\Id-M$ and
$Q:=C^{1/2}C_{n+1}^{-1}C^{1/2}=(\Id+\Delta t\,M)/(1+\Delta t)$, whose eigenvalues $q_i$ satisfy
$q_i-1=-\tfrac{\Delta t}{1+\Delta t}s_i$ (with $s_i=1-\mu_i$ the eigenvalues of $S$) and
$\max\{q_i,1\}\leq(1+\Delta t\,\LB)/(1+\Delta t)$. Using the scalar bound $q^{-1}-1+\log q\geq\tfrac{(q-1)^2}{2\max\{q,1\}^2}$
for $q>0$ (for $q\geq1$, $q^{-1}-1+\log q=\int_1^q(t-1)t^{-2}\dd t\geq\tfrac{(q-1)^2}{2q^2}$; for $q\leq1$,
$=\int_q^1(1-t)t^{-2}\dd t\geq\tfrac12(1-q)^2$) eigenvalue by eigenvalue,
\[
    \KL_C(\rho_{a_{n+1}}\|\rho_a)=\tfrac12\sum_i(q_i^{-1}-1+\log q_i)
    \geq\frac{\Delta t^2}{4(1+\Delta t\,\LB)^2}\sum_i s_i^2=\frac{\Delta t^2}{4(1+\Delta t\,\LB)^2}\norm S_{\Frob}^2 .
\]
Finally $\norm S_{\Frob}^2=\Tr(CRCR)=\Tr\bigl(C[RCR]\bigr)\geq L_n\Tr(RCR)$, since $RCR\succeq0$ and
$C\succeq L_n\Id$. Adding the mean and covariance parts gives \eqref{eq:KL-lower}.
\end{proof}

Together \Cref{lem:KL-argmin,lem:scalar-KL,lem:KL-onestep-app} prove \Cref{prop:KL-onestep}. Every
constant---the variational maps, the descent margin, and the lower-bound factor
$\tfrac{L_n\Delta t^2}{4(1+\Delta t\LB)^2}$---is derived here; nothing is imported.

\section{Non-Gaussian local modulus: detailed derivation}
\label{app:Kng}

This appendix gives the full auditable proofs of the score identities (\Cref{lem:second-score}) and
the non-Gaussian modulus bound (\Cref{lem:Kng}). We use the whitened notation of \Cref{sec:local}:
$v=C^{-1/2}u$, $S=C^{-1/2}UC^{-1/2}$, $t=\norm\zeta_a$, $X=m+C^{1/2}Z$, $Z\sim\N(0,\Id)$, and
$\bar\Sigma_\delta:=\min\{\Sigma_\delta,\beta_\star-\alpha_\star\}$. Recall the whitened score
generator $\mathcal L$ (Ornstein--Uhlenbeck), which acts as multiplication by $-k$ on the $k$-th
Hermite chaos, so that for $\phi=\sum_{k\geq1}\phi_k$ one has, by Gaussian integration by parts
($-\mathcal L=\nabla^*\nabla$),
$\E\norm{\nabla(-\mathcal L)^{-1}\phi}^2=\E[\phi\,(-\mathcal L)^{-1}\phi]=\sum_{k\geq1}\tfrac1k\norm{\phi_k}_{L^2}^2$.
In particular, when $\phi$ has only chaoses $\geq2$ (as does $\chi_\zeta$),
$\E\norm{\nabla(-\mathcal L)^{-1}\phi}^2\leq\tfrac12\norm\phi_{L^2}^2$; the bound fails for chaos-one
$\phi$ (e.g.\ $\phi=Z_1$ gives equality with constant $1$), but only chaoses $\geq2$ occur below.

\begin{proof}[Proof of \Cref{lem:second-score}]
\emph{Differentiation identities.} With $m_s=m+su$, $C_s=C+sU$ (positive definite for small $|s|$),
differentiate $\log\rho_{a_s}(x)=-\tfrac12\log\det C_s-\tfrac12(x-m_s)^\top C_s^{-1}(x-m_s)+\mathrm{const}$
at fixed $x=m+C^{1/2}Z$. The entropy term contributes $-\tfrac12\Tr S$ at first order and
$+\tfrac12\Tr S^2$ at second; the quadratic term contributes $v^\top Z+\tfrac12Z^\top SZ$ at first order
and $-\norm v^2-2v^\top SZ-Z^\top S^2Z$ at second. Hence $\partial_s\log\rho_{a_s}|_0=q_\zeta$ and
$\partial_s^2\log\rho_{a_s}|_0=\tfrac12\Tr S^2-\norm v^2-2v^\top SZ-Z^\top S^2Z=:p_\zeta$. From
$\partial_s\rho_{a_s}=\rho_{a_s}\partial_s\log\rho_{a_s}$ and
$\partial_s^2\rho_{a_s}=\rho_{a_s}\bigl((\partial_s\log\rho_{a_s})^2+\partial_s^2\log\rho_{a_s}\bigr)$,
differentiation under the integral (licensed by polynomial growth of $\phi$ and locally bounded
polynomial score factors) gives $D\Phi[\zeta]=\E[\phi\,q_\zeta]$ and
$D^2\Phi[\zeta,\zeta]=\E[\phi\,\chi_\zeta]$ with $\chi_\zeta=q_\zeta^2+p_\zeta$.

\emph{Chaos structure.} Write $\ell:=v^\top Z$ (chaos $1$) and $Q:=Z^\top SZ-\Tr S$ (chaos $2$), so
$q_\zeta=\ell+\tfrac12Q$ and, using $\operatorname{Var}(Z^\top SZ)=2\norm S_{\Frob}^2$,
$\E[q_\zeta^2]=\norm v^2+\tfrac14\cdot2\norm S_{\Frob}^2=t^2$. The chaos-two product formula gives
$Q^2=(Q^2)_0+(Q^2)_2+(Q^2)_4$ with $(Q^2)_0=2\norm S_{\Frob}^2$ and $(Q^2)_2=4(Z^\top S^2Z-\Tr S^2)=:4R_2$;
and $\ell Q=2v^\top SZ+G_3$ with $G_3$ pure chaos three, its chaos-one part being
$\sum_k\E[\ell QZ_k]Z_k=2v^\top SZ$ (from $\E[Z_i(Z_jZ_l-\delta_{jl})Z_k]=\delta_{ij}\delta_{lk}+\delta_{il}\delta_{jk}$).
Writing $Z^\top S^2Z=R_2+\Tr S^2$ in $p_\zeta$ gives $p_\zeta=-2v^\top SZ-R_2-\tfrac12\Tr S^2-\norm v^2$.
Collecting $\chi_\zeta=\ell^2+\ell Q+\tfrac14Q^2+p_\zeta$ by chaos order: the chaos-$0$ part is
$\norm v^2+\tfrac12\norm S_{\Frob}^2-\tfrac12\Tr S^2-\norm v^2=0$; the chaos-$1$ part is
$2v^\top SZ-2v^\top SZ=0$; and, with $\tfrac14(Q^2)_2=R_2$, one is left with
$\chi_2=\ell^2-\norm v^2$, $\chi_3=G_3=\ell Q-2v^\top SZ$, $\chi_4=\tfrac14(Q^2)_4$, mutually
orthogonal. In particular $\E[\chi_\zeta]=0$ and $\E[Z\chi_\zeta]=0$.

\emph{Sharp constant.} By orthogonality $\norm{\chi_\zeta}_{L^2}^2=\E[\chi_2^2]+\E[\chi_3^2]+\E[\chi_4^2]$.
With $a:=\norm v^2$, $b:=\norm S_{\Frob}^2$, $c:=v^\top S^2v$, $d_4:=\Tr S^4$: $\E[\chi_2^2]=\operatorname{Var}((v^\top Z)^2)=2a^2$;
$\E[\chi_3^2]=\E[(\ell Q)^2]-4\norm{v^\top SZ}_{L^2}^2=2ab+4c$; and
$\E[\chi_4^2]=\tfrac1{16}\norm{(Q^2)_4}_{L^2}^2=\tfrac12b^2+d_4$. Hence
$\norm{\chi_\zeta}_{L^2}^2=2a^2+2ab+4c+\tfrac12b^2+d_4$. Since $c=v^\top S^2v\leq\norm v^2\norm S_{\op}^2\leq ab$
and $d_4=\sum_i\lambda_i(S)^4\leq(\sum_i\lambda_i(S)^2)^2=b^2$,
$\norm{\chi_\zeta}_{L^2}^2\leq2a^2+6ab+\tfrac32b^2\leq6(a+\tfrac b2)^2=6t^4$. Equality holds at $v=0$
($a=c=0$) with $S$ rank one ($d_4=b^2$): then $\norm{\chi_\zeta}_{L^2}^2=\tfrac12b^2+b^2=\tfrac32b^2=6(b/2)^2=6t^4$.
\end{proof}

\begin{proof}[Proof of \Cref{lem:Kng}]
Write $x:=\norm v$, $y:=\norm S_{\Frob}/\sqrt2$, so $t^2=x^2+y^2$; also
$\norm U_{\op}\leq\sqrt2U_\delta y$, $\norm U_{\Frob}\leq\sqrt2U_\delta y$, $\norm C_{\op}\leq U_\delta$,
and convert at the end via $t\leq\ell_\delta^{-1}\norm\zeta_\star$.

\emph{Mean block.} Since $H_m(a)=C(g(a)+m)$ and $D^2m=0$,
$D^2H_m[\zeta,\zeta]=2U\,D(g+m)[\zeta]+C\,D^2g[\zeta,\zeta]$. Now
$g(a)+m=-\E[\nabla\widehat V(X)-X]$, and by the Ornstein--Uhlenbeck inverse of $q_\zeta$ (with
$\nabla\psi_1=v+\tfrac12SZ$), $D(g+m)[\zeta]=-\E[(\nabla^2\widehat V(X)-\Id)C^{1/2}(v+\tfrac12SZ)]$.
Splitting $\nabla^2\widehat V-\Id=(\nabla^2\widehat V-A(a))+(A(a)-\Id)$ and using
$\E[\nabla^2\widehat V-A(a)]=0$ and $\E[SZ]=0$,
\[
    \norm{D(g+m)[\zeta]}\leq\sqrt{U_\delta}\Bigl(\mathcal B_\delta\norm v+\tfrac1{\sqrt2}\Sigma_\delta y\Bigr),
\]
where $\norm{A(a)-\Id}_{\op}\leq\mathcal B_\delta$ controls the $(A-\Id)v$ term and
$(\E\norm{\nabla^2\widehat V-A}_{\Frob}^2)^{1/2}\leq\Sigma_\delta$ the $\tfrac12(\nabla^2\widehat V-A)SZ$
term. For the second derivative, $\chi_\zeta$ has chaoses $\geq2$, so with
$\psi_2=(-\mathcal L)^{-1}\chi_\zeta$ (whence $\E\norm{\nabla\psi_2}^2\leq\tfrac12\norm{\chi_\zeta}_{L^2}^2\leq3t^4$
and $\E\nabla\psi_2=0$), $D^2g[\zeta,\zeta]=-\E[(\nabla^2\widehat V(X)-A(a))C^{1/2}\nabla\psi_2]$, hence
$\norm{D^2g[\zeta,\zeta]}\leq\sqrt{3U_\delta}\,\bar\Sigma_\delta t^2$. Here $\bar\Sigma_\delta=\min\{\Sigma_\delta,\beta_\star-\alpha_\star\}$
arises because $\E\norm{(\nabla^2\widehat V-A)C^{1/2}\nabla\psi_2}$ may be bounded either by the pointwise
operator norm $\norm{\nabla^2\widehat V-A}_{\op}\leq\beta_\star-\alpha_\star$ (dimension-free) or by the
Frobenius dispersion $(\E\norm{\nabla^2\widehat V-A}_{\Frob}^2)^{1/2}\leq\Sigma_\delta$ via
Cauchy--Schwarz, whichever is smaller; no bound of the form $\norm{\nabla^2\widehat V-A}_{\Frob}\leq\beta_\star-\alpha_\star$
is used. Therefore
\[
    \norm{D^2H_m[\zeta,\zeta]}\leq U_\delta^{3/2}\bigl(2\sqrt2\,\mathcal B_\delta xy+2\Sigma_\delta y^2+\sqrt3\,\bar\Sigma_\delta t^2\bigr),
\]
and the quadratic form in $(x,y)$ formed by the first two terms has largest eigenvalue
$\Sigma_\delta+\sqrt{\Sigma_\delta^2+2\mathcal B_\delta^2}$; adding $\sqrt3\bar\Sigma_\delta t^2$ and
converting $t$ gives $K_{\mathrm{ng},m}$.

\emph{Covariance block.} $H_C(a)=-C(A(a)-\Id)C=-CWC$ with $W:=A(a)-\Id$; the product rule gives
\[
    D^2(CWC)[\zeta,\zeta]=2U(DW)C+2C(DW)U+2UWU+C(D^2W)C .
\]
With $\norm{DW[\zeta]}_{\Frob}=\norm{DA[\zeta]}_{\Frob}\leq\Sigma_\delta t$ (first score identity),
$\norm W_{\op}\leq\mathcal B_\delta$, and $\norm{D^2W[\zeta,\zeta]}_{\Frob}=\norm{D^2A[\zeta,\zeta]}_{\Frob}\leq\sqrt6\,\Sigma_\delta t^2$
(second score identity, \Cref{lem:second-score}),
\[
    \norm{D^2H_C[\zeta,\zeta]}_{\Frob}
    \leq4\sqrt2\,U_\delta^2\Sigma_\delta ty+4\mathcal B_\delta U_\delta^2y^2+\sqrt6\,U_\delta^2\Sigma_\delta t^2
    \leq U_\delta^2\bigl((4\sqrt2+\sqrt6)\Sigma_\delta+4\mathcal B_\delta\bigr)t^2,
\]
using $ty\leq t^2$ and $y^2\leq t^2$; converting $t$ gives $K_{\mathrm{ng},C}$. Finally the pair-norm
bound follows from $\norm{D^2H[\zeta,\zeta]}_\star^2=\norm{(D^2H)_m}^2+\tfrac12\norm{(D^2H)_C}_{\Frob}^2$.
Every constant is proportional to $\mathcal B_\delta$, $\Sigma_\delta$, or $\bar\Sigma_\delta$, all of
which vanish for Gaussian targets.
\end{proof}

\bibliographystyle{plainnat}
\nocite{roeder2017sticking,lambert2022variational,chen2023sampling,amari1998natural,opper2009variational,kim2024stl,domke2023provable}
\bibliography{note_refs}

\end{document}
